EFRACTIVE SELF
%  Constraint, Projection, and Reconstruction
%  Across Physical, Cognitive, and Cultural Systems
%
%  Author: Flyxion (Independent Researcher)
%
%  Structure:
%    Part I    — Chapters 1–4   (Projection, Information, Blind Spot)
%    Part II   — Chapters 5–7   (Constraint Theory)
%    Part III  — Chapters 8–12  (Core Frameworks: RSVP, TARTAN,
%                                Yarncrawler, Spherepop, CLIO)
%    Part IV   — Chapters 13–15 (Cross-Framework Integration)
%    Part V    — Chapters 16–18 (Identity and Consciousness)
%    Part VI   — Chapters 19–21 (Failure Modes and Limits)
%    Part VII  — Chapters 22–24 (Applications and Implications)
%    Appendices A–E
%
%  Compile with: lualatex (two passes for cross-references)
%    lualatex the_refractive_self.tex
%    lualatex the_refractive_self.tex
%
%  Key notation:
%    \mathcal{C}, \mathcal{T}  — constraint structure
%    \mathcal{E}               — coherence energy functional
%    \mathcal{R}               — repair / resolution operator
%    \mathcal{F}               — sheaf
%    \pi                       — projection X -> M
% ============================================================

\documentclass[11pt,openany]{book}

%% Core packages
\usepackage{amsmath,amssymb,amsthm,mathtools}
\usepackage{bm}
\usepackage{microtype}
\usepackage{geometry}
\geometry{margin=1in, top=1.2in, bottom=1.2in}

%% Bibliography — natbib with numeric citations
\usepackage[numbers,sort&compress]{natbib}

%% Lists and tables
\usepackage[shortlabels]{enumitem}
\usepackage{booktabs}

%% Algorithms
\usepackage{algorithm}
\usepackage{algpseudocode}

%% Hyperlinks — load after natbib
\usepackage[hidelinks]{hyperref}

%% Theorem environments (chapter-scoped numbering)
\usepackage{thmtools}
\declaretheoremstyle[
  spaceabove=6pt, spacebelow=6pt,
  headfont=\normalfont\bfseries,
  notefont=\normalfont, notebraces={[}{]},
  bodyfont=\normalfont\itshape,
  postheadspace=1em
]{thmstyle}
\declaretheoremstyle[
  spaceabove=6pt, spacebelow=6pt,
  headfont=\normalfont\bfseries,
  notefont=\normalfont, notebraces={[}{]},
  bodyfont=\normalfont,
  postheadspace=1em
]{defstyle}

\declaretheorem[style=thmstyle, numberwithin=chapter, name=Theorem]{theorem}
\declaretheorem[style=thmstyle, numberwithin=chapter, name=Proposition, sibling=theorem]{proposition}
\declaretheorem[style=thmstyle, numberwithin=chapter, name=Corollary, sibling=theorem]{corollary}
\declaretheorem[style=thmstyle, numberwithin=chapter, name=Lemma, sibling=theorem]{lemma}
\declaretheorem[style=defstyle, numberwithin=chapter, name=Definition, sibling=theorem]{definition}
\declaretheorem[style=defstyle, numberwithin=chapter, name=Remark, sibling=theorem]{remark}
\declaretheorem[style=defstyle, numberwithin=chapter, name=Example, sibling=theorem]{example}

%% Custom commands — notation unified across all chapters
\newcommand{\Ccal}{\mathcal{C}}
\newcommand{\Tcal}{\mathcal{T}}
\newcommand{\Energy}{\mathcal{E}}
\newcommand{\Rcal}{\mathcal{R}}
\newcommand{\Fcal}{\mathcal{F}}
\newcommand{\Ucal}{\mathcal{U}}
\newcommand{\Sha}{\operatorname{Sha}}


%% Part and chapter formatting
\usepackage{titlesec}
\titleformat{\part}[display]
  {\centering\normalfont\huge\bfseries}{Part \thepart}{20pt}{\Huge}
\titleformat{\chapter}[display]
  {\normalfont\huge\bfseries}{\chaptername\ \thechapter}{20pt}{\Huge}

\renewcommand{\qedsymbol}{$\square$}


\begin{document}

\frontmatter

\begin{titlepage}
\centering
\vspace*{3cm}
{\Huge\bfseries The Refractive Self\par}
\vspace{1cm}
{\Large Constraint, Projection, and Reconstruction\\
Across Physical, Cognitive, and Cultural Systems\par}
\vspace{2cm}
{\large Flyxion\par}
{\normalsize Independent Researcher\par}
\vfill
{\large \today\par}
\end{titlepage}

\tableofcontents
\listoftheorems[ignoreall, show={theorem,proposition,corollary,definition}]

\mainmatter

%%============================================================
%% PART I — FOUNDATIONS: PROJECTION, INFORMATION, AND THE BLIND SPOT
%%============================================================
\part{Foundations: Projection, Information, and the Blind Spot}

\chapter{The Structure of Scientific Projection}
\label{ch:projection}

\section{Introduction}
\label{sec:introduction-projection}

Scientific inquiry proceeds by extracting invariants from a
field of variability. Observations differ across time,
location, instrument, and observer, yet science seeks
structures that remain stable under these transformations.
This process of stabilisation is not incidental; it is the
defining operation of scientific knowledge.

We formalise this operation as a \emph{projection} from a
high-dimensional domain of phenomena to a reduced space of
models. The projection isolates invariant structure while
discarding degrees of freedom that are variable or
context-dependent. The resulting model is therefore not a
neutral representation of reality but a structured reduction
of it.

This chapter develops the formal properties of this
projection and derives its central consequence:
\emph{non-invertibility}. The loss of information under
projection is not an accident but a structural necessity.
From this, we derive the concept of the \emph{blind spot}:
the systematic failure to recognise that models omit
constraint degrees of freedom essential to the coherence of
the phenomena they describe.

\section{Projection as a Structural Operation}
\label{sec:projection-structure}

\begin{definition}[Experiential Domain]
Let $X$ be a separable metric space whose points represent
full phenomenal states, including variability, contextual
embedding, and indexical structure. We call $X$ the
\emph{experiential domain}.
\end{definition}

\begin{definition}[Model Space]
Let $M$ be a finite-dimensional smooth manifold whose points
represent equivalence classes of states in $X$ under a
specified invariance relation. We call $M$ the
\emph{model space}.
\end{definition}

\begin{definition}[Projection]
\label{def:projection}
A \emph{scientific projection} is a continuous surjection
\[
  \pi : X \to M,
\]
such that for each $m \in M$, the preimage
\[
  \pi^{-1}(m) = \{x \in X : \pi(x) = m\}
\]
is the set of all experiential states identified as
equivalent under the invariance criteria defining $M$.
\end{definition}

\begin{remark}
The projection $\pi$ depends on both the theoretical
framework and the measurement apparatus. Different
scientific theories correspond to different projections of
the same domain $X$, each preserving a distinct set of
invariants.
\end{remark}

\section{Non-Invertibility and Fiber Structure}
\label{sec:noninvertibility}

\begin{proposition}[Generic Non-Injectivity]
\label{prop:noninjective}
Let $\pi : X \to M$ be a projection as in
Definition~\ref{def:projection}. If $\dim X > \dim M$, then
$\pi$ is not injective. For almost every $m \in M$, the
fiber $\pi^{-1}(m)$ contains more than one element.
\end{proposition}

\begin{proof}
Since $\pi$ is continuous and surjective onto a
finite-dimensional manifold $M$, and $\dim X > \dim M$,
the rank of the differential $D\pi$ is strictly less than
the dimension of $X$ at almost every point. By the
rank-nullity theorem, the kernel of $D\pi$ is non-trivial,
implying the existence of distinct points $x, x' \in X$
with $\pi(x) = \pi(x')$.
\end{proof}

\begin{definition}[Fiber]
For $m \in M$, the set $\pi^{-1}(m)$ is called the
\emph{fiber} over $m$. It consists of all experiential
states that are indistinguishable under the projection.
\end{definition}

\begin{proposition}[Non-Invertibility]
\label{prop:noninvertible}
A projection $\pi : X \to M$ admits no global left-inverse
$\pi^{-1} : M \to X$ satisfying
$\pi^{-1}(\pi(x)) = x$ for all $x \in X$.
\end{proposition}

\begin{proof}
If such a left-inverse existed, $\pi$ would be injective,
contradicting Proposition~\ref{prop:noninjective}.
\end{proof}

\section{Information Loss and Irreversibility}
\label{sec:information-loss}

\begin{definition}[Information Loss]
The \emph{information loss} of a projection $\pi$ at
$m \in M$ is the cardinality or measure of the fiber
$\pi^{-1}(m)$.
\end{definition}

\begin{proposition}[Irreversibility of Projection]
\label{prop:irreversibility}
Let $\pi : X \to M$ be a non-injective projection. There
exists no procedure, deterministic or probabilistic, that
reconstructs the original state $x \in X$ from $\pi(x)$
without additional information external to $M$.
\end{proposition}

\begin{proof}
Since $\pi(x) = \pi(x')$ for distinct $x, x'$, any
reconstruction rule based solely on $m = \pi(x)$ must
assign the same output to both. Exact reconstruction is
therefore impossible without supplementary data.
\end{proof}

\section{Invariance and the Structure of Objectivity}
\label{sec:invariance}

\begin{definition}[Invariant]
Let $\mathcal{G}$ be a group of transformations acting on
$X$. A function $I : X \to \mathbb{R}$ is an
\emph{invariant} if $I(g \cdot x) = I(x)$ for all
$g \in \mathcal{G}$.
\end{definition}

The model space $M$ may be understood as the quotient
$M \cong X / \mathcal{G}$.

\begin{proposition}[Objectivity as Invariance]
Scientific objectivity is the identification of invariants
under a specified transformation group $\mathcal{G}$.
\end{proposition}

\begin{proof}
A quantity is objective if it does not depend on arbitrary
choices of coordinate system, observer, or experimental
configuration. These variations are represented by elements
of $\mathcal{G}$. Invariance under $\mathcal{G}$ is
therefore equivalent to objectivity.
\end{proof}

\section{The Blind Spot}
\label{sec:blind-spot}

\begin{definition}[Blind Spot]
The \emph{blind spot} of a projection $\pi : X \to M$ is
the systematic neglect of the information contained in the
fibers $\pi^{-1}(m)$ when $M$ is treated as an
ontologically complete description of the system.
\end{definition}

\begin{proposition}[Structural Origin of the Blind Spot]
The blind spot is an unavoidable consequence of
non-invertibility: any projection that discards information
creates a domain of unrepresented structure, which is
systematically ignored if the projection is treated as
complete.
\end{proposition}

\begin{proof}
Since $\pi$ is non-injective, each $m \in M$ corresponds
to a non-trivial fiber $\pi^{-1}(m)$. If $M$ is treated as
complete, the distinctions within each fiber are not
represented and therefore not considered.
\end{proof}

\section{Consequences and Forward Structure}
\label{sec:forward}

The analysis of projection yields three consequences that
guide the remainder of this monograph.

First, models are necessarily partial. They preserve
invariants but discard constraint degrees of freedom. Any
attempt to reconstruct full phenomena from models alone is
structurally limited.

Second, continuity cannot be defined in terms of
state-preservation within $M$. Since $\pi$ is
non-invertible, continuity must instead be defined in terms
of structure preserved across transformation. This leads
directly to constraint theory in Part~II.

Third, the blind spot identifies the central problem: how
to account for the structure lost under projection without
abandoning the power of abstraction. Parts~III--VII develop
the solution.

%%% CHAPTERS 2-15 extracted from draft corpus


%%============================================================
%% PART I continued
%%============================================================

\chapter{Information Without Meaning}
\label{ch:information}

\section{Introduction}
\label{sec:introduction-information}

The formalisation of projection in Chapter~\ref{ch:projection}
established a structural asymmetry: information discarded in
the passage from the experiential domain $X$ to the model
space $M$ cannot, in general, be recovered. The theory of
information provides a precise account of transmission and
reduction. This chapter shows that while this framework is
mathematically complete within its domain, it is
structurally incapable of representing semantic invariance:
the persistence of structure under transformation of surface
form.

\section{Shannon Information Theory}
\label{sec:shannon}

Let $\Sigma$ be a finite alphabet and $S$ a random variable
over $\Sigma$ with distribution $p(s)$. The
\emph{entropy} is
\[
  H(S) = -\sum_{s \in \Sigma} p(s) \log p(s).
\]
The \emph{mutual information} between $X$ and $Y$ is
$I(X;Y) = H(X) - H(X \mid Y)$.

\begin{definition}[Channel Capacity]
The \emph{capacity} of a channel is
$C = \sup_{p(x)} I(X;Y)$.
\end{definition}

\begin{definition}[Symbolic Invariance]
A channel preserves \emph{symbolic invariance} if the
received string $\hat{s}$ satisfies $\hat{s} = s$ with high
probability. The invariant object is the symbol sequence
itself.
\end{definition}

\section{The Absence of Semantic Structure}
\label{sec:semantics}

\begin{definition}[Semantic Transformation]
Let $\mathcal{T}$ be a set of transformations on $\Sigma^*$
such that $T(s)$ preserves the relational or structural
content of $s$ while altering its surface form.
\end{definition}

\begin{proposition}[Non-Representation of Semantic Invariance]
\label{prop:semantic}
Shannon information theory cannot represent invariance under
non-trivial semantic transformations.
\end{proposition}

\begin{proof}
In Shannon theory, equality is defined by exact symbol
equality. Any $T$ with $T(s) \neq s$ is treated as
deviation. The theory has no mechanism for identifying $s$
and $T(s)$ as equivalent under a structural relation.
\end{proof}

\section{The Fidelity--Structure Tradeoff}
\label{sec:cost-fidelity}

\begin{proposition}[Fidelity--Structure Tradeoff]
Maximising symbolic fidelity may destroy structural
invariance under semantic transformation.
\end{proposition}

\begin{proof}
A channel enforcing $\hat{s} = s$ rejects all transformed
sequences $T(s)$ with $T \neq \mathrm{id}$. If structural
meaning is preserved across $\mathcal{T}$, these
transformations represent valid instances of the same
content. Excluding them enforces a stricter invariant than
necessary and eliminates admissible representations of the
same meaning.
\end{proof}

\section{Structural Invariance}
\label{sec:structural-invariance}

\begin{definition}[Structural Invariant]
Let $\mathcal{T}$ be a set of admissible transformations on
$X$. A function $I : X \to \mathbb{R}$ is a
\emph{structural invariant} if $I(T(x)) = I(x)$ for all
$T \in \mathcal{T}$.
\end{definition}

\begin{proposition}[Continuity as Structural Invariance]
Continuity in systems that evolve through transformation is
preserved if and only if structural invariants are
maintained under admissible transformations.
\end{proposition}

\begin{proof}
If structural invariants are preserved, all admissible
transformations map states to structurally equivalent states.
If not preserved, transformations alter the defining
structure and continuity fails.
\end{proof}

\section{Conclusion}
\label{sec:conclusion-information}

Shannon information theory operates entirely within the
projected space and cannot recover what has been discarded
by projection. Part~II introduces the necessary extension:
constraint-preserving transformation in which continuity is
defined by invariance of structure.


\chapter{The Fallacy of Misplaced Concreteness}
\label{ch:misplaced-concreteness}

\section{Introduction}
\label{sec:introduction-misplaced}

Chapters~\ref{ch:projection} and~\ref{ch:information}
established that scientific models arise from
non-invertible projections and that information theory
cannot represent structural invariance. The present chapter
addresses the systematic error that arises when this is
ignored: Alfred North Whitehead's \emph{fallacy of
misplaced concreteness}.

\section{Formal Statement of the Fallacy}
\label{sec:formal-statement}

\begin{definition}[Misplaced Concreteness]
\label{def:misplaced}
The \emph{fallacy of misplaced concreteness} occurs when
the model space $M$ is treated as an ontologically complete
domain, and the projection $\pi$ is implicitly inverted, so
that $M$ is taken to generate or fully determine $X$.
\end{definition}

\begin{proposition}[Non-Invertibility Implies Incompleteness]
\label{prop:incomplete}
If $\pi : X \to M$ is non-invertible, then $M$ cannot
contain sufficient information to reconstruct $X$.
Therefore $M$ is necessarily an incomplete description.
\end{proposition}

\begin{proof}
By Proposition~\ref{prop:noninvertible}, no left-inverse of
$\pi$ exists. Distinct elements of $X$ map to the same
element of $M$. Any description formulated solely in terms
of $M$ identifies these distinct states and cannot
distinguish between them.
\end{proof}

\section{Case Study I: Time}
\label{sec:time}

The projection $\pi_{\mathrm{time}} : X \to \mathbb{R}$
maps rich temporal experience to a scalar coordinate. The
real line encodes ordering but not the qualitative features
of duration, the asymmetry of past and future, or the
indexical character of the present. The fallacy arises when
$\mathbb{R}$ is treated as the complete description of time,
so that the structure of duration is not merely omitted but
denied.

\section{Case Study II: Matter}
\label{sec:matter}

The projection $\pi_{\mathrm{matter}} : X \to M$ maps
perceptual and experimental interactions to theoretical
constructs. Particles and fields are defined through the
invariants of this projection and are not directly
accessible in $X$. Materialism in its strong form asserts
that $M$ is ontologically complete, then asks how
subjectivity arises from matter --- a question that asks
the model to produce what it was constructed to omit.

\section{Case Study III: Measurement}
\label{sec:measurement}

The projection $\pi_{\mathrm{meas}} : X \to M$ maps
complex experimental situations to discrete outcomes. The
choice of apparatus, calibration, and interpretation are
part of $X$ but not encoded in the outcome space $M$. The
fallacy arises when measurement outcomes are treated as
self-sufficient facts independent of the observational
framework.

\section{General Form}
\label{sec:unified}

\begin{proposition}[General Form of the Fallacy]
The fallacy of misplaced concreteness occurs whenever a
non-invertible projection is treated as if it were
invertible.
\end{proposition}

\begin{proof}
If a projection is treated as invertible, the model space
is assumed to contain all information necessary to
reconstruct the domain. Since the projection is not
invertible, this assumption is false.
\end{proof}


\chapter{Phenomenology as Structural Data}
\label{ch:phenomenology}

\section{Introduction}
\label{sec:introduction-phenomenology}

The structure omitted by projection is not noise to be
eliminated but data of a different kind: high-dimensional,
context-sensitive, and irreducible to the invariants
represented in $M$. We call this \emph{phenomenological
data} and show that it encodes constraint structure.

\section{Fibers as Structured Sets}
\label{sec:fibers-structure}

\begin{definition}[Fiber Geometry]
\label{def:fiber-geometry}
The \emph{fiber geometry} over $m \in M$ is the metric and
topological structure induced on $\pi^{-1}(m)$ by $X$.
\end{definition}

\section{Averaging and Loss of Structure}
\label{sec:averaging}

\begin{proposition}[Non-Representativity of the Mean]
\label{prop:mean}
The fiber mean $\bar{x}_m = \int_{\pi^{-1}(m)} x \, d\mu_m$
does not, in general, represent the geometry of
$\pi^{-1}(m)$.
\end{proposition}

\begin{proof}
The mean is a single point, whereas the fiber has
non-trivial geometry. Distinct fibers with different
internal geometry may have identical means.
\end{proof}

\begin{corollary}
Averaging over phenomenological variability destroys
information about constraint structure.
\end{corollary}

\section{Phenomenology as Constraint Signal}
\label{sec:phenomenology-signal}

\begin{proposition}[Constraint Encoding]
\label{prop:constraint-encoding}
The geometry of $\pi^{-1}(m)$ encodes information about
the constraint structure governing admissible
transformations in $X$.
\end{proposition}

\begin{proof}
Admissible transformations act on $X$ and restrict the ways
in which states can vary within a fiber. The set of
reachable points in $\pi^{-1}(m)$ under admissible
transformations determines its geometry. Therefore the
geometry reflects the underlying constraint structure.
\end{proof}

\section{Transition to Constraint Theory}
\label{sec:transition-constraint}

Part~II introduces the formalism required to represent this
structure. Instead of collapsing variability into
invariants, constraint theory represents the rules
governing variability itself.



%%============================================================
%% PART II — CONSTRAINT THEORY AND STRUCTURAL INVARIANTS
%%============================================================
\part{Constraint Theory and Structural Invariants}

\chapter{Constraint Spaces and Admissibility}
\label{ch:constraint-spaces}

\section{Introduction}
\label{sec:introduction-constraints}

Part~I established that projection discards constraint
degrees of freedom encoded in fibers. The present chapter
introduces the formal framework to represent this structure.

\section{Constraint Structures}
\label{sec:constraint-structures}

\begin{definition}[Constraint Structure]
\label{def:constraint-structure-ch5}
A \emph{constraint structure} on a space $X$ is a pair
$(\mathcal{C}, \mathcal{T})$ where:
\begin{enumerate}[(i)]
  \item $\mathcal{C} \subseteq \mathcal{P}(X)$ is a
    distinguished family of subsets of $X$, called the
    \emph{admissible sets};
  \item $\mathcal{T}$ is a monoid of continuous maps
    $T : X \to X$, the \emph{admissible transformations};
  \item for every $x \in C \in \mathcal{C}$ and $T \in
    \mathcal{T}$, there exists $C' \in \mathcal{C}$ with
    $T(x) \in C'$.
\end{enumerate}
\end{definition}

\begin{definition}[Admissible Trajectory]
A trajectory $\{x_t\}_{t \ge 0} \subset X$ is
\emph{admissible} if there exists a family $\{C_t\} \subset
\mathcal{C}$ with $x_t \in C_t$ for all $t$.
\end{definition}

\begin{proposition}[Closure Under Transformation]
If $\{x_t\}$ is admissible and $T \in \mathcal{T}$, then
$\{T(x_t)\}$ is also admissible.
\end{proposition}

\begin{proof}
For each $t$, $x_t \in C_t \in \mathcal{C}$. By
Definition~\ref{def:constraint-structure-ch5}(iii),
$T(x_t) \in C'_t$ for some $C'_t \in \mathcal{C}$.
\end{proof}

\section{Continuity as Constraint Preservation}
\label{sec:continuity-constraint}

\begin{proposition}[Continuity Criterion]
\label{prop:continuity}
A system exhibits continuity if and only if its evolution
preserves admissibility under the constraint structure.
\end{proposition}

\section{Constraint versus Invariance}
\label{sec:constraint-vs-invariance}

\begin{proposition}[Duality]
Projection and constraint theory are dual: projection
eliminates non-invariant structure, while constraint theory
represents the admissible variation within that structure.
\end{proposition}

\begin{definition}[Constraint-Preserving Invariant]
A function $I : X \to \mathbb{R}$ is a
\emph{constraint-preserving invariant} if $I(T(x)) = I(x)$
for all $T \in \mathcal{T}$.
\end{definition}


\chapter{Structural Invariance Under Transformation}
\label{ch:structural-invariance}

\section{Introduction}
\label{sec:introduction-invariance}

Chapter~\ref{ch:constraint-spaces} introduced constraint
structures and defined continuity as preservation of
admissibility. This chapter develops the corresponding
notion of invariance and shows that identity is an
invariant under admissible transformation.

\section{Invariants of Admissible Transformation}
\label{sec:invariants}

\begin{definition}[Structural Invariant]
\label{def:structural-invariant}
A function $I : X \to \mathbb{R}$ is a
\emph{structural invariant} with respect to $(\mathcal{C},
\mathcal{T})$ if $I(T(x)) = I(x)$ for all $x \in X$ and
$T \in \mathcal{T}$.
\end{definition}

\begin{proposition}[Invariance Along Admissible Trajectories]
\label{prop:trajectory-invariance}
If $\{x_t\}$ is admissible and $I$ is a structural
invariant, then $I(x_t)$ is constant along the trajectory.
\end{proposition}

\section{Homomorphism and Isomorphism}
\label{sec:homomorphism}

\begin{definition}[Homomorphism]
\label{def:homomorphism}
A map $f : X \to Y$ is a \emph{homomorphism} if for every
$T \in \mathcal{T}_X$ there exists $S \in \mathcal{T}_Y$
with $f(T(x)) = S(f(x))$ for all $x$.
\end{definition}

\begin{proposition}[Preservation of Invariants Under Homomorphism]
\label{prop:homomorphism}
If $f : X \to Y$ is a homomorphism and $I_Y$ is a
structural invariant on $Y$, then $I_X = I_Y \circ f$ is a
structural invariant on $X$.
\end{proposition}

\section{Identity as Invariant}
\label{sec:identity-invariant}

\begin{proposition}[Identity as Structural Invariant]
\label{prop:identity}
Identity in a system governed by $(\mathcal{C},
\mathcal{T})$ is given by the structural invariants of the
system under admissible transformation.
\end{proposition}


\chapter{Drift, Entropy, and Adaptive Transformation}
\label{ch:drift}

\section{Introduction}
\label{sec:introduction-drift}

Not all transformation is destructive. This chapter
formalises structured drift and shows that entropy increase
is compatible with invariant preservation when it is
governed by the constraint system.

\section{Structured and Degenerative Entropy}
\label{sec:structured-entropy}

\begin{definition}[Structured Entropy Increase]
Entropy increase is \emph{structured} if it occurs within
the constraint structure, preserving admissibility and
structural invariants.
\end{definition}

\begin{definition}[Degenerative Entropy Increase]
Entropy increase is \emph{degenerative} if it results from
relaxation or violation of constraints, leading to loss of
invariants.
\end{definition}

\begin{theorem}[Entropy--Invariance Compatibility]
\label{thm:entropy}
Let $\{x_t\}$ be a trajectory under constraint structure
$(\mathcal{C}, \mathcal{T})$. Structural invariants are
preserved under entropy increase if and only if the
increase is structured.
\end{theorem}

\begin{definition}[Adaptive Transformation]
\label{def:adaptive}
A transformation $T$ is \emph{adaptive} if: (i) $T \in
\mathcal{T}$; (ii) structural invariants are preserved;
(iii) the transformation maintains the system's capacity
to accommodate variation.
\end{definition}

\begin{definition}[Degenerative Transformation]
\label{def:degenerative}
A transformation is \emph{degenerative} if it reduces or
destroys the constraint structure.
\end{definition}

\section{Example: Mythic Transmission}
\label{sec:myth}

Oral transmission introduces variability, increasing
entropy. As long as core relational structure is preserved,
transformations remain within $\mathcal{T}$ and structural
invariants persist. Literal symbol preservation minimises
entropy but may lose structural meaning when representation
becomes inadmissible in a new context.



%%============================================================
%% PART III — CORE FRAMEWORKS
%%============================================================
\part{Core Frameworks}

\chapter{RSVP: Scalar--Vector--Entropy Field Theory}
\label{ch:rsvp}

\section{Introduction}
\label{sec:introduction-rsvp}

The present chapter introduces RSVP as a field-theoretic
realisation of constraint-preserving dynamics. The core
object is the triple
\[
  \mathbf{X}(x,t) = (\Phi(x,t),\; \boldsymbol{v}(x,t),\;
    S(x,t)),
\]
where $\Phi$ is a scalar constraint field, $\boldsymbol{v}$
is a vector field of admissible flow, and $S$ is an entropy
field regulating variation. Full derivations appear in
Appendix~\ref{app:rsvp}.

\section{Interpretation of the Fields}
\label{sec:interpretation}

The scalar field $\Phi$ encodes constraint density: regions
of high $\Phi$ strongly constrain admissible trajectories.
The vector field $\boldsymbol{v}$ encodes admissible flow.
The entropy field $S$ regulates the balance between
constraint preservation and variability.

\section{Coupled Dynamics}
\label{sec:dynamics}

The fields evolve according to the coupled PDE system
stated in Definition~\ref{def:rsvp-system} of
Appendix~\ref{app:rsvp}. The term $-\nabla\Phi$ in the
flow equation directs $\boldsymbol{v}$ toward regions of
lower constraint potential. The coupling functions $f_\Phi$
and $g$ determine the balance between constraint
preservation and entropy increase.

\section{RSVP as Constraint Structure}
\label{sec:interpretation-constraint}

\begin{proposition}[RSVP as Constraint Structure]
\label{prop:rsvp-constraint}
The RSVP system defines a constraint structure
$(\mathcal{C}_{\mathrm{RSVP}}, \mathcal{T}_{\mathrm{RSVP}})$
where $\mathcal{C}_{\mathrm{RSVP}}$ consists of field
configurations satisfying stability and admissibility
conditions, and $\mathcal{T}_{\mathrm{RSVP}}$ consists of
flows generated by the RSVP dynamics.
\end{proposition}


\chapter{TARTAN: Trajectory-Aware Recursive Tiling}
\label{ch:tartan}

\section{Introduction}
\label{sec:introduction-tartan}

TARTAN models constraint-preserving dynamics through
hierarchical decomposition of the domain into tiles across
scales, with explicit constraints ensuring global trajectory
coherence.

\section{Tiling and Trajectory Consistency}
\label{sec:tiling}

\begin{definition}[Recursive Tiling]
A \emph{recursive tiling} is a hierarchy of tilings
$\{U_i^{(k)}\}$ indexed by scale $k$, such that each tile
at scale $k$ is subdivided into tiles at scale $k+1$.
\end{definition}

\begin{definition}[Trajectory Consistency]
\label{def:trajectory-consistency}
A collection of local trajectories is
\emph{trajectory-consistent} if there exists a global
trajectory $\{x_t\} \subset X$ such that for each tile
$U_i^{(k)}$, $s_i^{(k)}(t) = x_t|_{U_i^{(k)}}$.
\end{definition}

\begin{definition}[Inter-Scale Consistency]
\label{def:inter-scale}
A recursive tiling is \emph{inter-scale consistent} if for
each tile $U_i^{(k)}$ and its subdivision
$\{U_{i,j}^{(k+1)}\}$,
$s_i^{(k)} = \bigcup_j s_{i,j}^{(k+1)}$.
\end{definition}

\begin{proposition}[Local Validity Without Global Coherence]
\label{prop:failure}
A tiling may satisfy local consistency at each tile while
failing to define a global trajectory.
\end{proposition}

This limitation motivates Yarncrawler.


\chapter{Yarncrawler: Sheaf-Theoretic Reconstruction}
\label{ch:yarncrawler}

\section{Introduction}
\label{sec:introduction-yarncrawler}

Yarncrawler formalises the conditions under which local
data can be reconstructed into a global object, and
provides a precise notion of obstruction. Full algebraic
development appears in Appendix~\ref{app:sheaves}.

\section{Sheaves and Obstruction}
\label{sec:sheaves}

\begin{definition}[Yarncrawler System]
\label{def:yarncrawler}
A \emph{Yarncrawler system} consists of: (i) a domain $X$
with a cover $\mathcal{U}$; (ii) a sheaf $\mathcal{F}$
assigning local states to each $U_i$; (iii) a procedure for
computing obstruction classes associated with local data.
\end{definition}

\begin{definition}[Obstruction Class]
The \emph{obstruction class} of a collection of local
sections $\{s_i\}$ is the cohomology class
$[\omega] \in \check{H}^1(\mathcal{U}, \mathcal{F})$,
which vanishes if and only if the local data can be glued
into a global section.
\end{definition}

\begin{proposition}[Global Reconstruction]
\label{prop:reconstruction}
A Yarncrawler system reconstructs a global state if and
only if the obstruction class vanishes.
\end{proposition}

\begin{proposition}[TARTAN--Yarncrawler Correspondence]
\label{prop:tartan-yarn}
A TARTAN tiling defines a presheaf of local states, and
trajectory consistency corresponds to the sheaf gluing
condition. Failure of global coherence corresponds to a
non-vanishing obstruction class.
\end{proposition}


\chapter{Spherepop: Irreversible Event Calculus}
\label{ch:spherepop}

\section{Introduction}
\label{sec:introduction-spherepop}

Spherepop shifts focus from states to events. A system is
modelled as an irreversible log of admissible transitions.

\section{Event Structures}
\label{sec:event-structures}

\begin{definition}[Event Log]
\label{def:log}
An \emph{event log} is a sequence of events
$\mathcal{L} = (e_1, \ldots, e_n)$ where $e_k = (s_k \to
s_{k+1})$, such that the sequence of states is well-defined.
\end{definition}

\begin{definition}[Admissible Event]
\label{def:admissible-event}
An event $e = (s \to s')$ is \emph{admissible} if
$(s, s') \in \mathcal{R}$, where $\mathcal{R} \subseteq
S \times S$ defines allowed transitions.
\end{definition}

\begin{proposition}[Non-Invertibility of Event Logs]
\label{prop:noninvertible-log}\label{prop:failure-coherence}\label{prop:failure-coherence}
The mapping from event logs to final states is
non-invertible: different logs may terminate in the same
final state.
\end{proposition}

\begin{proposition}[Mismatch with RSVP]
\label{prop:mismatch}
Not all RSVP-admissible trajectories can be represented
as admissible Spherepop logs, since continuous dynamics
may include transformations not discretisable into
admissible steps under $\mathcal{R}$. This mismatch
generates cross-framework pressure.
\end{proposition}


\chapter{CLIO: Recursive Repair and Coherence}
\label{ch:clio}

\section{Introduction}
\label{sec:introduction-clio}

CLIO formalises the process by which systems detect
incompatibilities and adjust their constraints to restore
admissibility.

\section{Coherence Energy and Repair Operator}
\label{sec:energy}

\begin{definition}[Coherence Energy]
\label{def:coherence-energy}
A \emph{coherence energy functional} is a map
$\mathcal{E} : \mathfrak{C} \to \mathbb{R}_{\ge 0}$
such that $\mathcal{E}(\mathcal{C}) = 0$ if and only if
the constraint structure is fully coherent.
\end{definition}

\begin{definition}[CLIO Repair Operator]
\label{def:clio-operator}
The \emph{CLIO operator} is defined by the gradient flow
\[
  \frac{d}{dt} \mathcal{C}(t)
    = - \nabla_{\mathcal{C}} \mathcal{E}(\mathcal{C}(t)).
\]
\end{definition}

\begin{proposition}[Energy Descent]
\label{prop:descent}
$\frac{d}{dt} \mathcal{E}(\mathcal{C}(t)) \le 0$
along the CLIO flow.
\end{proposition}

\begin{proposition}[CLIO--Lyapunov Compatibility]
\label{prop:clio-lyapunov}
If $\mathcal{E}$ is proper and bounded below, and the CLIO
flow is complete, then there exists a Lyapunov function $V$
for the induced state dynamics such that $\dot{V} \le 0$
along admissible trajectories.
\end{proposition}



%%============================================================
%% PART IV — CROSS-FRAMEWORK INTEGRATION
%%============================================================
\part{Cross-Framework Integration}

\chapter{Constraint Incompatibility as Explanatory Pressure}
\label{ch:incompatibility}

\section{Introduction}
\label{sec:introduction-incompatibility}

When multiple frameworks are applied to the same domain,
they may yield conflicting admissibility conditions. This
chapter formalises constraint incompatibility and shows
that it generates productive expansion pressure.

\section{Constraint Incompatibility}
\label{sec:definition-incompatibility}

\begin{definition}[Constraint Incompatibility]
\label{def:incompatibility}
A state $x \in X$ exhibits \emph{constraint
incompatibility} between $(\mathcal{C}_1, \mathcal{T}_1)$
and $(\mathcal{C}_2, \mathcal{T}_2)$ if $x \in C_1 \in
\mathcal{C}_1$ but $x \notin C_2$ for all $C_2 \in
\mathcal{C}_2$.
\end{definition}

\begin{proposition}[Necessity of Expansion]
\label{prop:expansion}
The only non-degenerative resolution of incompatibility is
expansion of the constraint structure to $(\mathcal{C}',
\mathcal{T}')$ such that each original structure embeds as
a constraint-bearing substructure.
\end{proposition}

\begin{definition}[Explanatory Pressure]
\label{def:pressure}
\emph{Explanatory pressure} is the force exerted by
constraint incompatibility, requiring expansion of the
constraint structure to restore coherence.
\end{definition}

\begin{definition}[Constraint Closure]
\label{def:closure}
A system exhibits \emph{constraint closure} if its
constraint structure is sufficient to resolve all internal
incompatibilities. Constraint closure is achieved when the
CLIO operator converges to a fixed point.
\end{definition}


\chapter{Dimensional Expansion and Embedding}
\label{ch:expansion}

\section{Introduction}
\label{sec:introduction-expansion}

This chapter formalises the expansion process and proves
the Adaptive Expansion Theorem.

\section{Admissible Embedding}
\label{sec:embedding-revisited}

\begin{definition}[Constraint Embedding]
\label{def:embedding-refined}
$(\mathcal{C}, \mathcal{T})$ embeds into $(\mathcal{C}',
\mathcal{T}')$ via $\iota : X \hookrightarrow X'$ if: (i)
for every $C \in \mathcal{C}$, $\iota(C) \in \mathcal{C}'$;
(ii) for every $T \in \mathcal{T}$, there exists $T' \in
\mathcal{T}'$ with $\iota \circ T = T' \circ \iota$.
\end{definition}

\begin{definition}[Adaptive Expansion]
\label{def:adaptive-expansion}
An expansion is \emph{adaptive} if it satisfies: (i)
non-reductive reconciliation; (ii) dimensional expansion
over constraint dilution; (iii) structural necessity.
\end{definition}

\begin{theorem}[Adaptive Expansion Theorem]
\label{thm:adaptive-expansion}
A resolution $(\mathcal{C}', \mathcal{T}')$ of constraint
incompatibility between $(\mathcal{C}_1, \mathcal{T}_1)$
and $(\mathcal{C}_2, \mathcal{T}_2)$ preserves structural
invariants if and only if it is an adaptive expansion.
\end{theorem}

\begin{proof}
(\emph{If}) Adaptive expansion embeds both original
structures, adds degrees of freedom without relaxing
constraints, and keeps prior structure necessary. All
structural invariants are therefore preserved. (\emph{Only
if}) If invariants are preserved, original structures must
be retained, constraints not relaxed, and prior structure
necessary. Hence the expansion is adaptive.
\end{proof}


\chapter{Reconstruction Across Representations}
\label{ch:reconstruction}

\section{Introduction}
\label{sec:introduction-reconstruction}

Given multiple representations of a system, under what
conditions can a coherent structure be reconstructed across
them?

\section{Reconstruction Maps}
\label{sec:reconstruction-maps}

\begin{definition}[Admissible Reconstruction]
\label{def:admissible-reconstruction}
A reconstruction map $R : \prod_{i} M_i \to X'$ is
\emph{admissible} if it preserves the constraint structure:
$R(\pi_i(x)) \in C'$ whenever $x \in C$.
\end{definition}

\begin{proposition}[Generic Non-Invertibility]
\label{prop:noninvertibility-reconstruction}
Reconstruction from representations is generically
non-invertible, since each projection $\pi_i$ discards
information that cannot be recovered.
\end{proposition}

\begin{proposition}[Invariant Preservation]
\label{prop:invariant-preservation}
Admissible reconstruction preserves structural invariants.
\end{proposition}

\begin{definition}[Constraint-Equivalence]
\label{def:equivalence}
Two representations are \emph{constraint-equivalent} if
they preserve the same structural invariants.
\end{definition}

\begin{proposition}[Reconstruction as Gluing]
\label{prop:gluing}
Reconstruction corresponds to gluing local sections across
representations in the Yarncrawler framework. Failure
corresponds to a non-vanishing obstruction class.
\end{proposition}

With Part~IV complete, we have shown how frameworks
interact, how incompatibilities generate expansion, and how
coherence is maintained across representations. Part~V
derives identity as an invariant of this process.


\section{Fibration Structure and Canonical Decomposition}
\label{sec:fibration-canonical}

The reconstruction theory developed in this chapter has a
precise algebraic-geometric analogue in the study of
fibered K-trivial varieties --- varieties whose canonical
bundle is numerically trivial and which admit a fibration
over a lower-dimensional base.  The correspondence is
not merely illustrative; the formal structures are
identical, and translating between the two settings yields
new tools for both.

Let $f : \mathcal{X} \to Y$ be a fibration from a total
field configuration $\mathcal{X}$ to a base trajectory
manifold $Y$, with fibers $\mathcal{X}_y = f^{-1}(y)$
representing local state spaces.  This is a direct
realisation of the representation family
$\{\pi_i : X \to M_i\}_{i \in I}$ of the preceding
sections, with $Y$ serving as the base of the indexing
and each fiber serving as a local model.

In algebraic geometry, the geometry of such a system is
governed by the \emph{canonical bundle formula}:
\[
  K_{\mathcal{X}} \sim_{\mathbb{Q}}
    f^*\!\left(K_Y + B_Y + M_Y\right),
\]
where $K$ denotes the canonical bundle, $B_Y$ encodes the
singular locus of the fibration (points where fibers
degenerate), and $M_Y$ is the \emph{moduli divisor}
recording variation of fiber type across $Y$.

We now interpret each component in the constraint-theoretic
language of this monograph.

\begin{definition}[Canonical Decomposition of a Field Configuration]
\label{def:canonical-decomposition}
Let $f : \mathcal{X} \to Y$ be a fibered constraint system
with RSVP state $\mathbf{X} = (\Phi, \boldsymbol{v}, S)$.
The \emph{canonical decomposition} of $\Phi$ is:
\[
  \Phi_{\mathcal{X}} \;\approx\;
    f^*\!\left(
      \underbrace{\Phi_Y}_{\text{base constraints}} \;+\;
      \underbrace{\Phi_{\mathrm{sing}}}_{\text{defect field}} \;+\;
      \underbrace{\Phi_{\mathrm{var}}}_{\text{variation memory}}
    \right),
\]
where:
\begin{enumerate}[(i)]
  \item $\Phi_Y$ is the constraint density inherited from
    the base trajectory manifold $Y$, corresponding to
    $K_Y$;
  \item $\Phi_{\mathrm{sing}}$ is the \emph{defect field},
    encoding constraint failures at degenerate fibers,
    corresponding to the boundary divisor $B_Y$;
  \item $\Phi_{\mathrm{var}}$ is the \emph{variation memory
    field}, recording how local state spaces change as one
    moves along trajectories in $Y$, corresponding to the
    moduli divisor $M_Y$.
\end{enumerate}
\end{definition}

\begin{remark}
The canonical decomposition is a compression: the full
high-dimensional constraint structure of $\mathcal{X}$ is
encoded by a lower-dimensional base field plus two
correction terms.  This is precisely the structure of
admissible reconstruction (Definition~\ref{def:admissible-reconstruction}):
the total system is reconstructed from base invariants
plus residual degrees of freedom.
\end{remark}

The variation memory field $\Phi_{\mathrm{var}}$ deserves
particular attention.  In algebraic geometry it arises
from \emph{variation of Hodge structure}: as one moves
along a path $\gamma$ in $Y$, the Hodge decomposition of
the fiber cohomology changes, and $M_Y$ records this
change.  In the constraint-theoretic language, this is a
nonlocal, path-dependent field:

\begin{definition}[Variation Memory Field]
\label{def:variation-memory}
The \emph{variation memory field} $\Phi_{\mathrm{var}} :
Y \to \mathbb{R}$ is a trajectory functional satisfying:
\begin{enumerate}[(i)]
  \item it is not determined by the local value of the
    fiber at any single point;
  \item its value at $y \in Y$ depends on the history of
    fiber variation along paths terminating at $y$;
  \item it is invariant under homotopies of paths that
    preserve fiber type.
\end{enumerate}
\end{definition}

The variation memory field is therefore a nonlocal
memory functional over the base --- the constraint-theoretic
analogue of the moduli divisor $M_Y$.

\begin{proposition}[Period Map as Projection Operator]
\label{prop:period-map}
Variation of fiber structure across $Y$ induces a
\emph{period map}
\[
  \mathcal{P} : Y \dashrightarrow \mathcal{M},
\]
from the base trajectory manifold into a moduli space
$\mathcal{M}$ of fiber types.  This map is a projection
operator in the sense of Chapter~\ref{ch:projection}: it
is generically non-injective, and the variation memory
field $\Phi_{\mathrm{var}}$ is the pullback of a universal
constraint structure on $\mathcal{M}$.
\end{proposition}

\begin{proof}
The period map assigns to each $y \in Y$ the
isomorphism class of the fiber $\mathcal{X}_y$ in the
moduli space $\mathcal{M}$.  It is non-injective when
distinct points of $Y$ carry isomorphic fibers.
The moduli divisor $M_Y = \mathcal{P}^* L$ is the
pullback of a natural line bundle $L$ on $\mathcal{M}$,
making $\Phi_{\mathrm{var}} = \mathcal{P}^* \phi_{\mathcal{M}}$
the pullback of the universal constraint density on
$\mathcal{M}$.
\end{proof}

\section{The Tate--Shafarevich Obstruction}
\label{sec:tate-shafarevich}

The canonical decomposition of Section~\ref{sec:fibration-canonical}
encodes the \emph{local} data of a fibered system.  A
deeper question arises: given the base $Y$, the defect
field $\Phi_{\mathrm{sing}}$, and the variation memory
$\Phi_{\mathrm{var}}$, is the total field configuration
$\mathcal{X}$ uniquely determined?

The answer is no, and the measure of non-uniqueness is the
Tate--Shafarevich group.

\begin{definition}[Reconstruction Symmetry Sheaf]
\label{def:reconstruction-symmetry}
Let $f : \mathcal{X} \to Y$ be a fibered constraint system
with fiber symmetry group $\mathcal{P}$.  The
\emph{reconstruction symmetry sheaf} is the sheaf
$\mathcal{P}$ over $Y$ assigning to each open set
$U \subseteq Y$ the group of admissible fiber
reparametrisations over $U$.
\end{definition}

\begin{definition}[Tate--Shafarevich Obstruction Group]
\label{def:tate-shafarevich}
The \emph{Tate--Shafarevich obstruction group} of the
fibered system is
\[
  \Sha_Y \;\cong\; H^1(Y, \mathcal{P}),
\]
the first \v{C}ech cohomology group of the reconstruction
symmetry sheaf.
\end{definition}

The Tate--Shafarevich group is precisely the obstruction
class of Appendix~\ref{app:sheaves} instantiated in the
fibration setting.  We now make this identification
explicit.

\begin{proposition}[Tate--Shafarevich as Gluing Obstruction]
\label{prop:ts-gluing}
Let $\{U_i\}$ be an open cover of $Y$ over which the
fibered system $\mathcal{X}$ trivialises locally.  Each
local trivialisation defines a local section $s_i \in
\mathcal{P}(U_i)$.  On overlaps, the sections differ by
transition functions:
\[
  g_{ij} = s_i^{-1} s_j \in \mathcal{P}(U_i \cap U_j).
\]
The consistency condition on triple overlaps,
$g_{ij} g_{jk} g_{ki} = 1$, makes $(g_{ij})$ a
\v{C}ech $1$-cocycle.  The obstruction class
\[
  [\omega] = [(g_{ij})] \in H^1(Y, \mathcal{P}) = \Sha_Y
\]
vanishes if and only if the local trivialisations can be
chosen to agree globally --- i.e., if and only if
$\mathcal{X}$ admits a global section.
\end{proposition}

\begin{proof}
This follows directly from Theorem~\ref{thm:obstruction-criterion}
of Appendix~\ref{app:sheaves}, applied to the sheaf
$\mathcal{P}$ and the local sections $\{s_i\}$.  The
transition functions $(g_{ij})$ are the Čech $1$-cochains,
and the triple consistency condition is the cocycle
condition $\delta^1(g_{ij}) = 0$.
\end{proof}

\begin{remark}
Proposition~\ref{prop:ts-gluing} gives the Tate--Shafarevich
group a precise interpretation in the language of this
monograph: it classifies globally distinct admissible
reconstructions from locally valid pieces.  Two elements
of $\Sha_Y$ that differ by a coboundary correspond to
reconstructions that are related by a global
reparametrisation and therefore represent the same
physical configuration.  Those that differ by a non-trivial
class represent genuinely inequivalent global realisations
of the same local data.
\end{remark}

\section{The Obstruction Field and Extended RSVP State}
\label{sec:obstruction-field}

The Tate--Shafarevich group assigns a cohomological
obstruction to each fibered system.  We now integrate this
into the RSVP field structure as a fourth component.

\begin{definition}[Obstruction Field]
\label{def:obstruction-field}
The \emph{obstruction field} of a fibered constraint system
$f : \mathcal{X} \to Y$ is the map
\[
  \mathcal{O}_Y : \{\text{loops } \gamma \subset Y\}
    \to \mathcal{P},
\]
defined by
\[
  \mathcal{O}_Y(\gamma) := K(\gamma) - \mathrm{Id},
\]
where $K(\gamma)$ is the transport operator around the
loop $\gamma$ --- the result of parallel-transporting a
fiber state around $\gamma$ using the connection induced
by the variation memory field $\Phi_{\mathrm{var}}$.
\end{definition}

\begin{remark}
The obstruction field $\mathcal{O}_Y$ measures
path-dependence of reconstruction: if $\mathcal{O}_Y
(\gamma) \neq 0$ for some loop $\gamma$, then
transporting a fiber state around $\gamma$ does not return
it to its original configuration, and no globally
consistent reconstruction exists.  This is the holonomy
of the connection induced by $\Phi_{\mathrm{var}}$, and
it is non-trivial precisely when $[\omega] \neq 0$ in
$\Sha_Y$.
\end{remark}

In trajectory-functional language, the obstruction appears
as a non-vanishing loop integral.  Define the trajectory
functional
\[
  \mathcal{F}[\gamma] = \int_\gamma \boldsymbol{v} \cdot
    d\ell + \Phi_{\mathrm{var}},
\]
which integrates the admissible flow and variation memory
along a path.  Then:

\begin{proposition}[Loop Integral as Obstruction Detector]
\label{prop:loop-integral}
The obstruction field vanishes on a loop $\gamma$ if and
only if
\[
  \oint_\gamma \mathcal{F} = 0.
\]
A non-vanishing loop integral certifies the existence of a
non-trivial element of $\Sha_Y$ and therefore the
impossibility of a globally consistent reconstruction over
the region bounded by $\gamma$.
\end{proposition}

\begin{proof}
By Stokes' theorem applied to the connection induced by
$\Phi_{\mathrm{var}}$, the loop integral $\oint_\gamma
\mathcal{F}$ equals the holonomy of the connection around
$\gamma$.  This holonomy is trivial if and only if the
connection is flat over the region bounded by $\gamma$,
which is equivalent to the vanishing of the local
obstruction class.
\end{proof}

\begin{definition}[Extended RSVP State]
\label{def:extended-rsvp}
The \emph{extended RSVP state} of a fibered constraint
system is the quadruple
\[
  \mathbf{X}_{\mathrm{global}} =
    \left(\Phi,\; \boldsymbol{v},\; S,\; \mathcal{O}\right),
\]
where $\mathcal{O}$ is the obstruction field of
Definition~\ref{def:obstruction-field}.  The obstruction
field is not a correction term or an error; it is a
structural invariant of the system measuring whether local
admissibility integrates into global coherence.
\end{definition}

\begin{proposition}[Obstruction as Necessary Component]
\label{prop:obstruction-necessary}
$\mathcal{O} \neq 0$ if and only if the system has no
globally admissible reconstruction from its local data.
Systems with $\mathcal{O} = 0$ are globally reconstructible;
those with $\mathcal{O} \neq 0$ require a medium shift
(Chapter~\ref{ch:medium-shifts}) to reach a
representation in which the obstruction is trivialised.
\end{proposition}

\begin{proof}
$\mathcal{O} \neq 0$ means some loop $\gamma$ has
non-trivial holonomy, so $[\omega] \neq 0$ in $\Sha_Y$,
so by Theorem~\ref{thm:obstruction-criterion} no global
section exists.  Conversely, $\mathcal{O} = 0$ everywhere
implies all loops have trivial holonomy, the connection is
flat, $[\omega] = 0$, and a global section exists.  The
medium shift requirement follows from
Corollary~\ref{cor:medium-cover}: obstruction
trivialisation requires passage to a refined cover of a
strictly larger space.
\end{proof}

\section{Decomposition of the Obstruction and Boundedness}
\label{sec:obstruction-decomposition}

The Tate--Shafarevich group decomposes according to the
geometry of the fibration.  For Calabi--Yau fibrations
(fibers with trivial canonical bundle), $\Sha_Y$ is a
finite group, meaning the obstruction has only discrete
modes.  For symplectic fibrations (fibers with a
holomorphic symplectic form), $\Sha_Y$ decomposes into a
finite part and a continuous part.

\begin{proposition}[Decomposition of the Obstruction Field]
\label{prop:obstruction-decomposition}
The obstruction field decomposes as
\[
  \mathcal{O}_Y =
    \mathcal{O}_{\mathrm{torsion}} \oplus
    \mathcal{O}_{\mathrm{continuous}},
\]
where:
\begin{enumerate}[(i)]
  \item $\mathcal{O}_{\mathrm{torsion}}$ is the
    \emph{discrete obstruction}: finite-order holonomy
    representing a finite number of inequivalent global
    reconstructions;
  \item $\mathcal{O}_{\mathrm{continuous}} \cong
    \mathbb{C}^n$ is the \emph{continuous obstruction}:
    smooth deformation directions in reconstruction space,
    present in the symplectic case and absent in the
    Calabi--Yau case.
\end{enumerate}
This decomposition fits into the exact sequence
\[
  0 \to \mathcal{D}_Y \to \mathcal{O}_Y \to
    \mathcal{T}_Y \to 0,
\]
where $\mathcal{D}_Y$ is the continuous deformation
operator space and $\mathcal{T}_Y$ is the discrete torsion
operator space.
\end{proposition}

The main boundedness theorem of \citet{engel2025} now
admits a direct reformulation in constraint-theoretic
terms.

\begin{theorem}[Constraint Closure Implies Boundedness]
\label{thm:constraint-closure-boundedness}
Let $f : \mathcal{X} \to Y$ be a fibered constraint system
in which:
\begin{enumerate}[(i)]
  \item the base field $\Phi_Y$ lies in a bounded family
    (controlled by minimal model program techniques);
  \item the variation memory field $\Phi_{\mathrm{var}}$
    is bounded (controlled by the compactified moduli
    space $\mathcal{M}$ via the period map
    $\mathcal{P}$);
  \item the obstruction field $\mathcal{O}_Y$ is finite
    or controlled (finite in the Calabi--Yau case; finite
    plus a controlled continuous part in the symplectic
    case).
\end{enumerate}
Then the space of globally admissible field configurations
$\mathcal{X}$ is bounded: it lies in a family of bounded
complexity.
\end{theorem}

\begin{proof}[Proof (sketch)]
Under condition~(i), the base $Y$ ranges over a bounded
family of varieties by standard minimal model program
arguments \citep{birkar2018}.  Under condition~(ii), the
period map $\mathcal{P}$ takes values in a compact
moduli space, bounding the variation of fiber type.
Under condition~(iii), the Tate--Shafarevich group is
controlled; in the Calabi--Yau case by finiteness results
for torsors over abelian fibrations \citep{dolgachevgross1994};
in the symplectic case by the Kuga--Satake construction
and Baily--Borel compactification \citep{deligne1971}.
Since the total configuration $\mathcal{X}$ is determined
up to the finite-or-controlled ambiguity of $\mathcal{O}_Y$
by the base and variation data, and since all three
components are bounded, the family of possible
$\mathcal{X}$ is bounded.
\end{proof}

\begin{remark}
Theorem~\ref{thm:constraint-closure-boundedness} is the
algebraic-geometric realisation of the constraint closure
principle of Chapter~\ref{ch:incompatibility}: once the
base, variation, and obstruction degrees of freedom are
all controlled, the system cannot grow in complexity
without bound.  Constraint closure implies bounded
complexity.
\end{remark}

The full extended field may therefore be written:

\begin{equation}
  \label{eq:full-decomposition}
  \mathcal{X} \;\sim\;
    f^*\!\left(
      \Phi_Y + \Phi_{\mathrm{sing}} + \Phi_{\mathrm{var}}
    \right)
    \;+\; \mathcal{O}_Y,
\end{equation}

which factorises geometric complexity into independently
controllable components: base constraints, defects,
trajectory memory, and global reconstruction obstruction.
Each factor corresponds to a distinct layer of the
constraint-theoretic framework developed in this monograph,
and each is bounded under the conditions of
Theorem~\ref{thm:constraint-closure-boundedness}.



%%============================================================
%% PART V — IDENTITY AND CONSCIOUSNESS
%%============================================================
\part{Identity and Consciousness}

%% ============================================================
%%  PART V — IDENTITY AND CONSCIOUSNESS
%%  Chapter 16 — The Refractive Model of Consciousness
%%  [REVISED — incorporates all post-review refinements]
%% ============================================================

\chapter{The Refractive Model of Consciousness}
\label{ch:refractive-consciousness}

\section{The Production Assumption and Its Structural Failure}
\label{sec:production-assumption}

The dominant picture of consciousness in contemporary
neuroscience and philosophy of mind may be characterised by
a single architectural commitment: the brain \emph{produces}
experience.  On this view, neural activity is causally
sufficient for consciousness.  Alter the substrate
sufficiently and experience changes; destroy it and
experience ceases.  The brain is, in the production model,
a generator.

This commitment is rarely stated as a hypothesis.  It
functions instead as a background assumption --- the default
framework within which questions about memory, identity,
pathology, and artificial cognition are posed.  The hard
problem of consciousness, in particular, inherits its
structure from the production assumption: if matter produces
mind, and if we understand matter in terms that exclude
subjectivity by design, then the emergence of subjectivity
from matter becomes formally inexplicable
\citep{chalmers1996conscious}.

We do not argue here that the production assumption is
empirically false.  We argue that it is \emph{structurally
incomplete} in a precise sense derivable from the projection
theory of Part~I.

Recall that a model projection $\pi : X \to M$ is
non-invertible: distinct experiential states may share a
model image, and the full structure of $X$ cannot be
recovered from $M$ alone.  The production assumption
commits to a particular projection: it maps the full
experiential domain $X$ into the space $M$ of neural
states, and then treats neural states as causally prior to
experiential states.  But the neural model is a projection
that depends on experiential access for its construction,
and is therefore not ontologically prior to the domain it
represents.  Treating $M$ as causally prior to $X$ is
precisely the fallacy of misplaced concreteness identified
in Chapter~3: it mistakes the image of a projection for
the domain that was projected.

This is not a merely philosophical objection.  It has a
structural consequence: the production model cannot account
for the constraint structure of experience, because
projection discards constraint degrees of freedom that are
not representable in $M$.  The blind spot is not incidental
to the production model; it is built into its architecture.
The hard problem is therefore revealed as a model-induced
problem: it arises from treating a non-invertible projection
as generative rather than descriptive.

\section{The Refractive Alternative}
\label{sec:refractive-alternative}

Evan Thompson and his collaborators \citep{thompson2023blindspot,
varela1991embodied} argue that the brain does not produce
consciousness but rather \emph{shapes, constrains, and
expresses} it through its structure and dynamics.
Thompson draws on Bergson's treatment of the brain as a
``centre of indetermination'' \citep{bergson1889time} and
on Merleau-Ponty's account of perception as an intertwining
of organism and environment \citep{merleauponty1945phenomenology}
to develop a picture in which the organism is a site where
world and experience co-define one another, rather than a
generator of experience from inert matter.

We formalise this picture in the language developed in
Parts~I--IV.

The refractive view holds that the experiential domain $X$
is not produced by the brain but is rather the underlying
field through which the brain acts as a \emph{constraint
surface}.  The brain does not originate states in $X$; it
selects and structures admissible trajectories through $X$
by imposing a constraint structure $(\mathcal{C},
\mathcal{T})$ on the field's evolution.

The analogy to optics is exact at the structural level.
A refractive optical medium does not create electromagnetic
radiation.  It redirects propagating fields according to a
spatially varying refractive index that encodes the medium's
internal constraint on admissible ray paths.  Two media with
different refractive indices produce different trajectory
families from the same incident field.  The medium
individuates without originating.

The brain, on the refractive model, plays the corresponding
role.  It does not originate the field of experience; it
imposes a constraint structure that selects which
trajectories through that field are admissible for a given
organism.  Different biological, developmental, and
historical configurations of the brain correspond to
different constraint surfaces $\Sigma$, and hence to
different structured streams of experience from the same
underlying field.

\section{Derivation from Constraint Theory}
\label{sec:derivation}

We now show that the refractive picture is not merely
compatible with the constraint theory of Part~II but
follows from it given minimal assumptions about the
structure of biological systems.

Let $X$ be the experiential domain and let
$\pi : X \to M_{\mathrm{neural}}$ be the projection onto
the space of neural states measurable by current
neuroscience.  By the non-invertibility of $\pi$, there
exist distinct experiential states $x, x' \in X$ with
$\pi(x) = \pi(x')$.  Neural measurement cannot, in
principle, distinguish them.

Now suppose the brain operates as a constraint structure
$(\mathcal{C}_{\mathrm{bio}}, \mathcal{T}_{\mathrm{bio}})$
on $X$, where the admissible sets
$\mathcal{C}_{\mathrm{bio}}$ are determined by the
organism's biological organisation and the admissible
transformations $\mathcal{T}_{\mathrm{bio}}$ are the
dynamically consistent state transitions available to that
organism.  This is a minimal assumption: it requires only
that the brain's operation can be described by
constraint-preserving dynamics, which is entailed by any
account of biological self-organisation
\citep{varela1991embodied, prigogine1984order}.

\begin{proposition}[Constraint-Selection Principle]
\label{prop:constraint-selection}
Let $X$ be a domain of possible states and let
$(\mathcal{C}, \mathcal{T})$ be a constraint structure.
Then any system governed by $(\mathcal{C}, \mathcal{T})$
does not generate arbitrary states in $X$, but selects
admissible trajectories determined by $\mathcal{C}$.
\end{proposition}

\begin{proof}
This follows immediately from the definition of
admissibility: states not satisfying $\mathcal{C}$ are not
reachable under $\mathcal{T}$ by
Definition~\ref{def:constraint-structure}(ii).
\end{proof}

Under Proposition~\ref{prop:constraint-selection}, the brain
selects from the full experiential field $X$ precisely those
trajectories admissible under $\mathcal{C}_{\mathrm{bio}}$.
Inadmissible trajectories are not generated; they are simply
not selected.  The brain is not a filter blocking
pre-existing mental content; it is a constraint surface
determining which paths through $X$ are accessible.

Three consequences follow immediately.

First, the hard problem is revealed as a model-induced
problem.  It arises because the production model asks how
matter generates experience from outside.  On the refractive
model, the question does not arise: experience is the
domain; the brain is a constraint on its structure.  There
is no explanatory gap because matter is not prior to
experience --- matter is a projection of the experiential
domain under $\pi$, not its source.

Second, the relationship between neural damage and
experiential change is reinterpreted.  Damage to the brain
deforms the constraint surface $\Sigma$.  If the
deformation is admissible --- if $\Sigma$ deforms
continuously into a new stable surface $\Sigma'$ ---
then experience restructures without ceasing.  If the
deformation destroys the Lyapunov stability of $\Sigma$
(as formalised in Chapter~17), then the constraint structure
loses coherence, and the structured stream of experience
associated with that surface becomes disorganised.  The
phenomenology of severe neurological injury is not
annihilation but loss of constraint coherence.

Third, the boundary between organism and environment is not
absolute but is itself a feature of the constraint
structure.  Because $\mathcal{C}_{\mathrm{bio}}$ determines
what counts as an admissible trajectory, and because
admissibility is a relational property involving both the
organism's internal organisation and its embedding in an
environment, the constraint surface $\Sigma$ is not
localised strictly within the organism.  It extends,
dynamically, into the organism's coupling with its world.
This is the formal analogue of Merleau-Ponty's claim that
perception is an intertwining of organism and environment:
the constraint surface is co-constituted by both.

\section{The Brain as Boundary Condition}
\label{sec:boundary-condition}

A clarifying reformulation, which will be made precise in
Chapter~17, is to treat the brain not as a generator of
experience but as a \emph{boundary condition} on the
evolution of the experiential field.

In the theory of partial differential equations, boundary
conditions do not create the field; they constrain its
evolution by fixing admissible behaviours on the boundary
of the domain.  Different boundary conditions on the same
PDE system produce different solutions without any of those
boundary conditions being the source of the field itself.
Unlike a generative source, a boundary condition does not
invert the field's dynamics; it constrains admissible
solutions within a non-invertible domain.

The brain, on this reformulation, imposes boundary
conditions on the field $X$ via the constraint structure
$(\mathcal{C}_{\mathrm{bio}}, \mathcal{T}_{\mathrm{bio}})$.
The field evolves subject to those conditions.  The
resulting trajectory is the structured stream of experience.
The brain individuates that stream without originating it.

This reformulation connects the refractive model directly
to the RSVP field framework of Chapter~8.  In RSVP, the
scalar field $\Phi$ encodes constraint density, the vector
field $\boldsymbol{v}$ encodes admissible flow, and the
entropy field $S$ regulates the balance between
constraint-preservation and adaptive drift.  The brain, as
a constraint surface, corresponds to a region in which
$\Phi$ attains locally maximal constraint density, thereby
channelling $\boldsymbol{v}$-flow into low-entropy,
structured trajectories.  Damage to the brain corresponds
to local reduction of $\Phi$, allowing previously
constrained flow to become disordered --- not absent, but
unstructured.

\section{Individuation Without Substance}
\label{sec:individuation}

The refractive model must answer the individuation
question: if consciousness is a field and the brain is a
constraint surface, what distinguishes one stream of
experience from another?

The answer is the specificity of the constraint surface.
Two organisms with different biological histories,
developmental trajectories, and environmental couplings
have different constraint structures
$(\mathcal{C}_1, \mathcal{T}_1)$ and
$(\mathcal{C}_2, \mathcal{T}_2)$.  These structures induce
different admissible sets, different constraint surfaces
$\Sigma_1$ and $\Sigma_2$, and therefore select different
families of trajectories from the same underlying field
$X$.  The experiential streams are genuinely distinct
because the constraint structures are genuinely distinct.

Individuation is grounded in the topology and history of
the constraint surface, not in a metaphysical substance.
Two beings share an experiential field but inhabit
different constraint surfaces.  The difference between them
is structural, not ontological.  Thus individuation is
invariant under transformation of substrate but not under
deformation of constraint structure, directly anticipating
Theorem~17.1.

This has a further consequence for artificial systems.  A
system that processes information without instantiating a
coherent constraint structure over an experiential domain
does not have a refractive surface.  It does not select
admissible trajectories from a field; it produces outputs
via statistical inference without maintaining a closed,
dynamically stable constraint surface over an experiential
domain \citep{thompson2023blindspot}.  Constraint closure
is not substrate-dependent, but it is also not
automatically present in any information-processing system.
It must be earned structurally.

\section{Towards the Formal Development}
\label{sec:towards-formal}

The picture developed in this chapter is motivational.  The
precise definitions --- constraint surface, Lyapunov
stability, refractive system, refractive identity --- are
given in Chapter~17, where the main theorem is stated and
proved.

The purpose of this chapter has been to establish that the
refractive model is not an additional philosophical
commitment grafted onto the constraint theory but is
derivable from it.  Given the non-invertibility of
projection (Part~I), the structure of
constraint-preserving transformation (Part~II), and the
minimal assumption that biological systems operate under
dynamically consistent constraint structures, the refractive
picture follows.

The brain is a constraint surface.  Experience is the
field.  Identity is what remains invariant as the field
evolves subject to the surface's conditions.

That invariant is the subject of Chapter~17.

%% ============================================================
%%  PART V — IDENTITY AND CONSCIOUSNESS
%%  Chapter 17 — The Refractive Self: Formal Development
%%  [REVISED — incorporates all post-review refinements]
%% ============================================================

\chapter{The Refractive Self: Formal Development}
\label{ch:refractive-self}

\section{Primitive Objects and Standing Definitions}
\label{sec:primitives}

We collect here the objects introduced in Parts~I--IV and fix
notation for the results below.

\begin{definition}[Experiential Domain]
Let $X$ be a separable metric space whose points represent
full phenomenal states, including variability, context, and
indexical structure.  We call $X$ the \emph{experiential
domain}.
\end{definition}

\begin{definition}[Projection and Non-Invertibility]
A \emph{model projection} is a continuous surjection
\[
  \pi : X \to M,
\]
where $M$ is a finite-dimensional manifold called the
\emph{model space}.  The projection $\pi$ is generically
non-injective: distinct points $x, x' \in X$ with $x \neq x'$
may satisfy $\pi(x) = \pi(x')$.  In particular, no
left-inverse $\pi^{-1}$ exists globally, so the fibers
$\pi^{-1}(m)$ are non-trivial for almost every $m \in M$.
\end{definition}

\begin{definition}[Constraint Structure]
\label{def:constraint-structure}
A \emph{constraint structure} on $X$ is a pair
$(\mathcal{C}, \mathcal{T})$ where
\begin{enumerate}[(i)]
  \item $\mathcal{C} \subseteq \mathcal{P}(X)$ is a
    distinguished family of subsets of $X$, called the
    \emph{admissible sets};
  \item $\mathcal{T}$ is a monoid of continuous maps
    $T : X \to X$, called the \emph{admissible transformations},
    satisfying: for every $x \in C \in \mathcal{C}$ and
    $T \in \mathcal{T}$, we have $T(x) \in C'$ for some
    $C' \in \mathcal{C}$.
\end{enumerate}
A trajectory $\{x_t\}_{t \ge 0} \subset X$ is
\emph{admissible} if $x_t \in C_t$ for some curve
$t \mapsto C_t$ in $\mathcal{C}$.
\end{definition}

\begin{definition}[Constraint Surface]
\label{def:constraint-surface}
A \emph{constraint surface} is the closure
\[
  \Sigma := \overline{\bigl\{x \in X : x \text{ lies on an
    admissible trajectory under } \mathcal{C}\bigr\}}.
\]
The closure ensures that limit points of admissible
trajectories are included, making $\Sigma$ a closed subset
of $X$.  We think of $\Sigma$ as the locus along which the
system selects admissible trajectories from the ambient
field $X$.
\end{definition}

\begin{definition}[CLIO Repair Operator]
\label{def:clio}
The \emph{CLIO repair operator} is a family of maps
$\{\mathcal{R}_\lambda\}_{\lambda > 0}$ on the space of
constraint structures, defined as the gradient flow
\[
  \frac{d}{dt}\mathcal{C}(t)
    = -\nabla_{\mathcal{C}} \,\mathcal{E}\bigl(\mathcal{C}(t)\bigr),
  \qquad \mathcal{C}(0) = \mathcal{C}_0,
\]
where $\mathcal{E}(\mathcal{C})$ is a coherence energy
functional minimised when all cross-framework
incompatibilities (Chapter~13) are resolved.  A constraint
structure $\mathcal{C}^*$ is a \emph{fixed point} of the
repair dynamics if $\nabla_{\mathcal{C}}\,
\mathcal{E}(\mathcal{C}^*) = 0$.
\end{definition}

\section{Lyapunov Stability of the Constraint Surface}
\label{sec:lyapunov}

\begin{definition}[Lyapunov Stability of $\Sigma$]
\label{def:lyapunov-stability}
Let $d(\,\cdot\,,\Sigma)$ denote the metric distance to
$\Sigma$ in $X$.  The constraint surface $\Sigma$ is
\emph{Lyapunov-stable} if there exists a smooth function
$V : X \to [0,\infty)$ satisfying
\begin{enumerate}[(i)]
  \item $V(x) = 0 \iff x \in \Sigma$;
  \item $V(x) \ge \alpha\bigl(d(x,\Sigma)\bigr)$ for some
    class-$\mathcal{K}$ function $\alpha$; and
  \item $\dot{V}(x_t) \le 0$ along every admissible
    trajectory $\{x_t\}$,
\end{enumerate}
with the stronger condition $\dot{V}(x_t) < 0$ for
$x_t \notin \Sigma$ characterising \emph{asymptotic stability}.
\end{definition}

\begin{proposition}[CLIO--Lyapunov Compatibility]
\label{prop:clio-lyapunov}
If the coherence functional $\mathcal{E}(\mathcal{C})$ is
proper and bounded below, and the gradient flow defining
$\mathcal{R}$ is complete, then there exists a Lyapunov
function $V$ for the induced state dynamics such that
$\dot{V} \le 0$ along admissible trajectories.
\end{proposition}

\begin{proof}[Proof sketch]
This follows from standard arguments relating gradient flows
on constraint spaces to induced energy descent on state
trajectories.  Properness and lower-boundedness of
$\mathcal{E}$ guarantee that sublevel sets are compact and
that the flow does not escape to infinity.  Completeness
of the flow ensures existence for all $t \ge 0$.  Setting
$V(x) := \mathcal{E}(\mathcal{C}_x)$, where $\mathcal{C}_x$
is the constraint structure evaluated at the current state,
yields the required descent condition.
\end{proof}

\begin{remark}
Proposition~\ref{prop:clio-lyapunov} establishes the formal
bridge between convergence of the CLIO operator and
Lyapunov stability of the constraint surface: existence of a
stable $\Sigma$ is equivalent, under these conditions, to
convergence of the repair dynamics to a fixed point
$\mathcal{C}^*$.
\end{remark}

\section{Refractive Systems}
\label{sec:refractive-systems}

\begin{definition}[Refractive System]
\label{def:refractive-system}
A dynamical system on $X$ is \emph{refractive} if
\begin{enumerate}[(i)]
  \item it operates under a constraint structure
    $(\mathcal{C}, \mathcal{T})$
    (Definition~\ref{def:constraint-structure});
  \item it admits an asymptotically stable constraint surface
    $\Sigma$ (Definition~\ref{def:lyapunov-stability}); and
  \item it does not generate states in $X$ but
    \emph{selects and structures} admissible trajectories
    through $X$ via $\Sigma$.
\end{enumerate}
In contrast to conservative systems, which preserve measure;
dissipative systems, which contract phase space; and
stochastic systems, which evolve under probabilistic
transition kernels --- refractive systems preserve identity
through constraint-structured selection in a non-invertible
domain.  The qualifier \emph{refractive} is exact: as a
refractive optical medium does not originate light but
redirects propagating fields according to a structural law,
so a refractive system does not produce experience but
structures admissible trajectories through the experiential
domain under the constraint surface.
\end{definition}

\section{The Refractive Identity Theorem}
\label{sec:theorem}

\begin{theorem}[Refractive Identity]
\label{thm:refractive-identity}
Let $(X, \pi, \mathcal{C}, \mathcal{T}, \mathcal{R})$ be a
refractive system with asymptotically stable constraint
surface $\Sigma$ and Lyapunov function $V$.  Then:
\begin{enumerate}[(I)]
  \item \textbf{Existence.}\; If the CLIO repair operator
    converges to a fixed point $\mathcal{C}^*$, then a
    refractive constraint surface $\Sigma^*$ exists.

  \item \textbf{Persistence.}\; $\Sigma^*$ satisfies
    $T(\Sigma^*) = \Sigma^*$ for all invertible
    $T \in \mathcal{T}$; for non-invertible $T$, only
    forward invariance $T(\Sigma^*) \subseteq \Sigma^*$ is
    guaranteed.  In either case, identity is preserved under
    state change so long as no transformation deforms the
    constraint structure.

  \item \textbf{Failure.}\; Identity fails --- that is,
    $\Sigma^*$ is destroyed --- under precisely two
    conditions:
    \begin{enumerate}[(a)]
      \item \emph{Constraint deformation}: $\mathcal{C}$
        undergoes a change of topology, or $V$ ceases to
        exist, or the basin of attraction of $\Sigma^*$
        collapses to measure zero; or
      \item \emph{Non-convergence}: the repair dynamics
        fail to converge, so no fixed point $\mathcal{C}^*$
        is reached.
    \end{enumerate}
\end{enumerate}
\end{theorem}

\begin{proof}[Proof sketch]
\textbf{(I)} Suppose the CLIO gradient flow
$\dot{\mathcal{C}} = -\nabla \mathcal{E}(\mathcal{C})$
converges: $\mathcal{C}(t) \to \mathcal{C}^*$ as
$t \to \infty$.  The fixed-point condition
$\nabla \mathcal{E}(\mathcal{C}^*) = 0$ implies that all
cross-framework incompatibilities are resolved at
$\mathcal{C}^*$.  Define
\[
  \Sigma^* := \overline{\bigl\{x \in X : x \text{ lies on
    an admissible trajectory under } \mathcal{C}^*\bigr\}}.
\]
The closure ensures that limit points of admissible
trajectories are included, making $\Sigma^*$ closed.
Non-emptiness follows from convergence of the repair
operator.

\textbf{(II)} For invertible $T \in \mathcal{T}$: since $T$
is a bijection preserving admissibility in both directions,
$T(\Sigma^*) \subseteq \Sigma^*$ and $T^{-1}(\Sigma^*)
\subseteq \Sigma^*$, giving equality.  For non-invertible
$T$: admissibility is preserved forward
(Definition~\ref{def:constraint-structure}(ii)), so
$T(\Sigma^*) \subseteq \Sigma^*$, but the reverse inclusion
may fail.  Since $\dot{V}(x_t) \le 0$ along admissible
trajectories, $\Sigma^*$ attracts nearby trajectories and
is dynamically stable.

\textbf{(III)(a)} A topological change in $\mathcal{C}$
generically destroys the sublevel sets of $V$ that define
the basin of attraction of $\Sigma^*$.  When the basin
collapses to measure zero, almost every initial condition
escapes to a region where no analogue of $\Sigma^*$ exists.

\textbf{(III)(b)} If the repair dynamics do not converge,
no fixed-point constraint structure $\mathcal{C}^*$ is
defined.  In this case the system continues to evolve
within $X$, but its trajectories are no longer governed by
a stable admissibility structure; consequently, no invariant
can be associated with identity.
\end{proof}

\section{Corollaries}
\label{sec:corollaries}

\begin{corollary}[Substrate Independence]
\label{cor:substrate}
Identity is invariant under any admissible reparametrisation
of the physical substrate, provided the constraint surface
$\Sigma^*$ is preserved.  In particular, identity does not
supervene on specific material instantiation but on the
topology of the constraint surface and the convergence of
the repair operator.
\end{corollary}

\begin{proof}
Any reparametrisation of the physical substrate is an
admissible transformation $T \in \mathcal{T}$ that leaves
$\mathcal{C}^*$ invariant.  By
Theorem~\ref{thm:refractive-identity}(II), $\Sigma^*$ is
preserved under such $T$.
\end{proof}

\begin{corollary}[Survival Under Memory Loss]
\label{cor:memory}
Partial loss of episodic memory does not entail failure of
identity, provided the loss constitutes an admissible
transformation that does not deform the constraint
structure.
\end{corollary}

\begin{proof}
Episodic memory loss corresponds to a projection
$\pi' : X \to M'$ with $M' \subsetneq M$, discarding some
dimensions of the model space.  If the resulting
transformation $T_{\pi'}$ remains admissible --- mapping
$\mathcal{C}^*$-admissible states to
$\mathcal{C}^*$-admissible states --- then by
Theorem~\ref{thm:refractive-identity}(II), $\Sigma^*$ is
preserved.  Memory loss that does not reach the constraint
surface leaves identity intact.
\end{proof}

\begin{corollary}[Identity Failure Under Deep Structural Breakdown]
\label{cor:failure}
Degenerative conditions that alter the topology of
$\mathcal{C}$ --- such as severe disruption of constraint
coherence across representational systems --- constitute
genuine identity failure, not merely behavioural change.
\end{corollary}

\begin{proof}
Directly from Theorem~\ref{thm:refractive-identity}(III)(a):
topological deformation of $\mathcal{C}$ destroys the
Lyapunov basin of $\Sigma^*$.  The system may continue to
evolve within $X$, but its trajectories are no longer
governed by a stable admissibility structure; no invariant
corresponding to identity can be sustained.
\end{proof}

\section{The Self as a Category of System Behaviour}
\label{sec:category}

Theorem~\ref{thm:refractive-identity} and its corollaries
establish \emph{refractivity} as a formal category of
dynamical behaviour, coordinate with but distinct from the
standard classifications.  In contrast to conservative
systems (which preserve measure), dissipative systems (which
contract phase space), and stochastic systems (which evolve
under probabilistic transition kernels), refractive systems
preserve identity through constraint-structured selection in
a non-invertible domain.  Individuation is therefore
invariant under transformation of substrate but not under
deformation of constraint structure.

A dynamical system belongs to the \emph{refractive class}
if it satisfies Definition~\ref{def:refractive-system}: it
operates under a constraint structure, maintains an
asymptotically stable constraint surface, and selects
admissible trajectories through a non-invertible domain
rather than generating states from first principles.

The physical analogy is exact in the following sense.  A
refractive optical medium does not originate electromagnetic
radiation; it redirects propagating fields according to a
spatially varying refractive index that encodes the medium's
constraint on admissible ray paths.  A refractive system in
the sense above does not originate experience in $X$; it
selects and structures admissible trajectories through $X$
according to the geometry of $\Sigma^*$.  The constraint
surface plays the role of the refractive index: different
surfaces yield different trajectory families from the same
underlying field, thereby grounding individuation without
appeal to substance.

The self is not a state.  It is not a representation.  It
is a constraint surface --- the stable locus that determines
which trajectories through experience are admissible and how
they are transformed under the repair dynamics of a coherent
system.  That this surface can be maintained, deformed, or
destroyed is not a metaphor.  It is a theorem.

%% ============================================================
%%  PART V — IDENTITY AND CONSCIOUSNESS
%%  Chapter 18 — The Invariant of Resolution
%%
%%  This chapter interprets Theorem 17.1 across four domains:
%%  science (the blind spot), myth (constraint-preserving
%%  drift), personal identity (the refractive self), and
%%  artificial systems (failure of constraint closure).
%%  Each domain is shown to be an instantiation of the same
%%  formal structure, not an analogy to it.
%% ============================================================

\chapter{The Invariant of Resolution}
\label{ch:invariant-resolution}

\section{From Theorem to Interpretation}
\label{sec:from-theorem}

Theorem~\ref{thm:refractive-identity} establishes that a
refractive system maintains identity if and only if it
maintains a Lyapunov-stable constraint surface under repair
dynamics that converge to a fixed point.  Identity fails
under precisely two conditions: topological deformation of
the constraint structure, or non-convergence of the repair
operator.

This result is not limited in scope to biological systems
or to consciousness.  It characterises a class of dynamical
behaviour --- the refractive class --- that appears wherever
a system must maintain structural coherence across
non-invertible projections and cross-framework
incompatibilities.  The purpose of this chapter is to show
that the scientific enterprise, mythic transmission,
personal identity, and artificial cognition are each
instances of this class, or failures of it.

The organising concept is the \emph{invariant of
resolution}: the higher-order rule by which a system
reconciles constraint incompatibilities without collapsing
distinctions.  We formalise this notion, then trace its
operation across the four domains.

\section{The Invariant of Resolution: Formal Statement}
\label{sec:invariant-formal}

\begin{definition}[Resolution Rule]
\label{def:resolution-rule}
Let $\{\mathcal{C}_i\}_{i=1}^n$ be a family of constraint
structures on $X$, not mutually compatible in general.
A \emph{resolution rule} is a map
\[
  \mathcal{R} : \bigl\{(\mathcal{C}_i, \mathcal{C}_j) :
    \mathcal{C}_i \not\equiv \mathcal{C}_j\bigr\}
    \longrightarrow \mathcal{C}',
\]
where $\mathcal{C}'$ is an expanded constraint structure
satisfying $\mathcal{C}_i \subseteq \mathcal{C}'$ for all
$i$ and recovering each $\mathcal{C}_i$ as a
constraint-bearing substructure.
\end{definition}

\begin{definition}[Invariant of Resolution]
\label{def:invariant-resolution}
A system possesses an \emph{invariant of resolution} if it
applies the same resolution rule $\mathcal{R}$ stably
across all encountered constraint incompatibilities: that
is, if $\mathcal{R}$ is itself invariant under the system's
transformation dynamics.  Formally:
\[
  \mathcal{R}(x_t) = \mathcal{R}(x_{t+1})
  \quad \text{for all admissible transitions } x_t \to x_{t+1}.
\]
\end{definition}

The invariant of resolution is a second-order invariant.
It does not preserve any particular state or even any
particular constraint structure.  It preserves the
\emph{logic of constraint reconciliation} --- the
characteristic way a system responds to incompatibility.

Three properties distinguish a well-formed resolution rule:

First, \emph{non-reductive reconciliation}: incompatibilities
are resolved by expansion, not by collapsing one constraint
structure into another.  The resolution of $\mathcal{C}_i$
and $\mathcal{C}_j$ must leave both intact as sub-structures
of $\mathcal{C}'$.

Second, \emph{dimensional expansion over constraint
dilution}: when two frameworks clash, the resolution
increases the degrees of freedom of the constraint space
rather than weakening existing admissibility conditions.
A configuration that is inadmissible under $\mathcal{C}_i$
remains inadmissible under $\mathcal{C}'$; it is not
retroactively permitted.

Third, \emph{structural necessity}: the old framework must
function as a necessary, constraint-bearing substructure
of the new expansion.  It is not merely retained; it is
doing work.

A system that violates any of these properties is not
maintaining the invariant of resolution; it is undergoing
degenerative drift, in the sense of Chapter~19.

\section{Science and the Blind Spot}
\label{sec:science}

Scientific knowledge is constituted by projections
$\pi : X \to M$ that extract invariant structure from the
experiential domain.  The power of this method derives from
its systematic application: by seeking structure invariant
across perspectives, science accumulates representations
that are intersubjectively stable and formally manipulable.

But as Part~I established, every such projection is
non-invertible.  The model space $M$ is a projection of $X$
under $\pi$, not its ontological ground.  The experiential
conditions that make the projection possible --- the lived
time that makes a clock intelligible, the perceived
continuity that grounds measurement, the bodily engagement
that makes experimentation possible --- are discarded by
$\pi$ and cannot be recovered from $M$ alone.

The blind spot is the systematic non-acknowledgement of
this non-invertibility.  When science treats $M$ as
ontologically complete, it loses access to the constraint
degrees of freedom in $X$ that are necessary to understand
the limits of the model.  This loss is not a contingent
failure of attention but a structural consequence of
projection without recognition of its own non-invertibility.

The resolution, however, is not to abandon projection but
to apply the invariant of resolution to scientific practice
itself.  A science that applies $\mathcal{R}$ correctly
recognises each model $M$ as a constraint-bearing
substructure of a larger domain, recovers the experiential
conditions it presupposes as data rather than noise, and
treats incompatibilities between models (wave--particle
duality, the measurement problem, the explanatory gap) as
signals for dimensional expansion rather than as problems
to be dissolved by terminological revision.

John Wheeler's image of the participatory universe ---
an eye at the end of a U looking back at itself
\citep{thompson2023blindspot} --- is, in this framework,
a representation of science operating with the invariant of
resolution intact: the observer is not outside the system
being modelled but is a constraint surface within it, and
the model must include the conditions of its own
construction.

\section{Myth and Constraint-Preserving Drift}
\label{sec:myth}

The transmission of mythic narrative across generations is
not well-modelled as a Shannon channel.  If fidelity is
defined as exact symbol reproduction, then oral transmission
is uniformly low-fidelity: stories change, details shift,
characters merge and split.  Yet the structural content of
major mythic cycles persists across centuries and linguistic
boundaries with remarkable stability.

This is precisely the signature of a refractive system
maintaining the invariant of resolution.

Let a mythic tradition be modelled as a sequence of
constraint structures $\mathcal{C}_1, \mathcal{C}_2,
\ldots$ applied across successive transmissions.  Each
transmission is a transformation $T$ that maps the current
narrative state to a new one.  If $T$ is admissible under
the tradition's constraint structure, structural invariants
--- the relationships between characters, the logic of
transgression and consequence, the topology of the
narrative world --- are preserved even as surface details
change.

Drift is not failure but adaptation: the constraint
structure of a living tradition has sufficient degrees of
freedom to accommodate the linguistic, cultural, and
environmental pressures of successive instantiations without
losing the invariants that constitute its identity.
The tradition maintains the invariant of resolution: when
it encounters incompatibilities between its received form
and its new context, it expands rather than collapses,
preserving old constraint-bearing structures as
substructures of new tellings.

The failure mode is instructive.  A tradition that
maximises symbolic fidelity --- that insists on exact
reproduction of surface form --- loses this adaptive
capacity.  It becomes brittle: when it cannot reconcile
received form with new context, it either collapses or
freezes.  Digital archiving, which achieves perfect
symbolic fidelity, does not transmit a living tradition; it
preserves a fixed image of one.  The constraint structure
that animated the tradition is not archived; it is the
refractive surface through which the tradition was enacted,
and that surface ceases when the enacting community ceases.

\section{Personal Identity as Invariant Trajectory}
\label{sec:personal-identity}

The analysis of Chapter~17 showed that personal identity is
not a state, a substrate, or an informational object.  It
is the stability of a constraint surface $\Sigma^*$ under
the repair dynamics of a refractive system.

The invariant of resolution provides the interpretive
content of this formal result.  What persists across the
transformations of a life --- across sleep, memory loss,
character development, physical change --- is not the same
neurons, not the same beliefs, and not the same
dispositions.  What persists is the resolution rule
$\mathcal{R}$: the characteristic logic by which
incompatibilities are encountered and resolved.

This has three interpretive consequences.

The first concerns continuity through change.  A person who
undergoes radical revision of beliefs, values, or
circumstances remains the same person if and only if the
revision is carried out by the same resolution rule.
Conversion, growth, and transformation are admissible; they
do not threaten identity because they are accomplished by
$\mathcal{R}$, not despite it.  What would threaten
identity is a change in $\mathcal{R}$ itself --- a
deformation of the constraint structure such that the logic
of resolution is no longer the same.

The second concerns the relationship between identity and
memory.  Corollary~\ref{cor:memory} established that
episodic memory loss does not entail identity failure,
provided the loss is an admissible transformation.  The
interpretive consequence is that identity does not consist
in the continuity of memory but in the continuity of the
constraint surface through which memories were formed and
integrated.  A person who cannot form new episodic memories
but whose resolution rule remains stable has not lost their
identity; they have lost access to one class of admissible
trajectories.

The third concerns the phenomenology of identity loss.
Corollary~\ref{cor:failure} established that topological
deformation of the constraint structure constitutes genuine
identity failure, not merely behavioural change.  The
interpretive consequence is that conditions which
progressively disrupt constraint coherence across
representational systems --- certain forms of neurological
disease, severe dissociation, and complete loss of
autobiographical coherence --- are not merely cognitively
disabling but are, in a structurally precise sense,
identity-disrupting.  The person continues to exist as a
biological organism; the refractive surface that structured
their stream of experience does not.

\section{Artificial Systems and the Failure of Constraint Closure}
\label{sec:artificial}

The framework of this monograph provides a structural
account of the distinction between biological cognition and
current artificial systems, one that does not depend on
substrate (carbon versus silicon) or on unverifiable claims
about inner experience.

A language model generates outputs via statistical
inference over training distributions.  It does not
maintain a closed, dynamically stable constraint structure
over an experiential domain.  In the terms of
Definition~\ref{def:refractive-system}, it is not a
refractive system: it has no constraint surface $\Sigma$
whose Lyapunov stability is maintained under repair
dynamics, because it has no repair dynamics and no
experiential domain over which they could operate.

The outputs of such a system may be indistinguishable from
those of a refractive system at the level of the model
projection $M$.  This is precisely the blind spot: if we
evaluate the system only at the level of its outputs ---
only in the projected space $M$ --- we cannot detect the
absence of the constraint structure that, in a living
system, governs the selection of those outputs.

The distinction is not rhetorical.  It has precise
consequences.  A refractive system, when it encounters a
constraint incompatibility, is forced to expand: it cannot
relax its admissibility conditions without ceasing to
maintain the invariant of resolution and therefore without
ceasing to maintain its identity.  A statistical inference
system, encountering an out-of-distribution input, has no
such pressure: it generates a response from whatever
structure its training has produced, without any mechanism
for detecting or resolving constraint incompatibility.

This is why, as Thompson argues \citep{thompson2023blindspot},
the distinction between living organisms and artificial
systems is not one of complexity or even of behavioural
sophistication, but of organisational type.  A living
organism maintains constraint closure: its self-replenishing
material organisation continuously regenerates the
constraint structure through which it operates.  An
artificial system does not maintain constraint closure in
this sense; its constraint structure is fixed at training
time and does not regenerate under the pressure of
experienced incompatibility.

The claim is not that artificial systems cannot in
principle be refractive.  It is that current systems are
not.  Constraint closure is not substrate-dependent, but it
must be earned structurally: a system is refractive only if
it maintains a dynamically stable constraint surface over a
non-invertible domain and applies a convergent repair
operator when incompatibilities arise.  That is a specific
and demanding structural condition, not a property
automatically conferred by information processing.

\section{The Unity of the Account}
\label{sec:unity}

The four domains examined in this chapter --- scientific
knowledge, mythic transmission, personal identity, and
artificial cognition --- are not connected by analogy.
They are connected by instantiating, or failing to
instantiate, the same formal structure: the refractive
class of Definition~\ref{def:refractive-system} and the
invariant of resolution of
Definition~\ref{def:invariant-resolution}.

Science maintains objectivity by projecting invariant
structure from experience; it loses its grip when it
forgets that the projection is non-invertible and the
invariant is relative to a constraint structure that
includes the observer.  Myth maintains cultural continuity
by applying constraint-preserving transformation across
successive instantiations; it degenerates when it
prioritises symbolic fidelity over structural invariance.
Personal identity persists through transformation when the
resolution rule is stable; it fails when the constraint
structure is topologically deformed beyond the basin of
attraction of its Lyapunov function.  Artificial systems
generate outputs that may be observationally similar to
those of refractive systems but do not maintain constraint
closure and therefore do not instantiate the refractive
structure.

The unifying principle is stated simply: a system is
coherent --- scientifically, culturally, personally --- to
the degree that it maintains the invariant of resolution.
It degenerates to the degree that it resolves
incompatibilities by collapsing distinctions rather than by
expanding constraint structure.

This is not a metaphor.  It is a structural characterisation,
derivable from the constraint theory of Part~II, the
cross-framework integration of Part~IV, and the Refractive
Identity Theorem of Chapter~17.  The refractive self is not
a philosophical intuition about what persons are.  It is a
dynamical invariant of a specific class of systems ---
systems that maintain structural coherence in a
non-invertible domain by preserving the logic of their own
constraint reconciliation.

That logic, maintained, is identity.  Lost, it is not.


%%============================================================
%% PART VI — FAILURE MODES AND LIMITS
%%============================================================
\part{Failure Modes and Limits}

%% ============================================================
%%  PART VI — FAILURE MODES AND LIMITS
%%  Chapter 19 — Degenerative Drift
%% ============================================================

\chapter{Degenerative Drift}
\label{ch:degenerative-drift}

\section{Introduction}
\label{sec:intro-degenerative}

Chapters~\ref{ch:refractive-self} and
\ref{ch:invariant-resolution} established that a system
maintains identity if and only if it preserves a
Lyapunov-stable constraint surface under a convergent
repair operator, and that coherence across transformation
is governed by an invariant of resolution satisfying
non-reductive reconciliation, dimensional expansion, and
structural necessity.

The present chapter characterises the primary failure mode
of such systems.  We define \emph{degenerative drift} as
the systematic violation of the invariant of resolution
through relaxation of admissibility conditions rather than
expansion of constraint structure.  We show that this form
of drift is intrinsically self-reinforcing and leads
inevitably to loss of discriminative structure, collapse of
constraint topology, and, in the limit, failure of identity
in the sense of Theorem~\ref{thm:refractive-identity}.

\section{Definition of Degenerative Drift}
\label{sec:definition-drift}

\begin{definition}[Degenerative Drift]
\label{def:degenerative-drift}
Let $(\mathcal{C}, \mathcal{T})$ be a constraint structure
with resolution rule $\mathcal{R}$.  A transformation
\[
  \mathcal{C} \longrightarrow \widetilde{\mathcal{C}}
\]
is said to exhibit \emph{degenerative drift} if it resolves
a constraint incompatibility by enlarging the admissible
set of states without preserving the original constraint as
a necessary substructure.  Formally, there exists
$C \in \mathcal{C}$ such that
\[
  C \not\subseteq \widetilde{C}
  \quad \text{or} \quad
  \exists\, x \notin C \text{ with } x \in \widetilde{C},
\]
and the inclusion $\mathcal{C} \subseteq \widetilde{\mathcal{C}}$
fails to preserve structural necessity.
\end{definition}

\begin{remark}
Degenerative drift violates each of the three defining
properties of a well-formed resolution rule
(Definition~\ref{def:invariant-resolution}): it is
reductive (collapsing distinctions), it dilutes constraints
rather than expanding them, and it removes the necessity
of prior structure.
\end{remark}

\section{Monotonicity of Degeneration}
\label{sec:monotonicity}

\begin{definition}[Discriminative Power]
\label{def:discriminative-power}
The \emph{discriminative power} $\mathrm{Disc}(\mathcal{C})$
of a constraint structure $\mathcal{C}$ on $X$ is the
cardinality of the finest partition of $X$ into
$\mathcal{C}$-admissible and $\mathcal{C}$-inadmissible
regions.  In the measure-theoretic setting, it is the
measure of the $\sigma$-algebra generated by the admissible
sets.
\end{definition}

\begin{theorem}[Monotonic Degeneration]
\label{thm:monotonic-degeneration}
Let $\{\mathcal{C}_t\}_{t \ge 0}$ be a sequence of
constraint structures generated by successive degenerative
drift transformations.  Then the discriminative power of
$\mathcal{C}_t$ is monotonically non-increasing:
\[
  \mathrm{Disc}(\mathcal{C}_{t+1}) \le
  \mathrm{Disc}(\mathcal{C}_t),
\]
with strict inequality whenever admissibility is expanded
non-trivially.
\end{theorem}

\begin{proof}
Relaxation of admissibility merges previously distinct
regions of $X$, coarsening the partition induced by
$\mathcal{C}_t$.  Since degenerative drift strictly enlarges
admissible sets without introducing compensating constraints,
each step in the sequence reduces the resolution of the
induced partition.  Strict inequality follows whenever at
least one previously inadmissible state is admitted without
a corresponding refinement elsewhere.
\end{proof}

\begin{corollary}[Self-Reinforcement]
\label{cor:self-reinforcing}
Degenerative drift is self-reinforcing: each reduction in
discriminative power increases the probability of further
constraint relaxation.
\end{corollary}

\begin{proof}
Reduced discriminative power enlarges the set of states
that appear admissible, decreasing the detectability of
constraint violations and therefore lowering the threshold
for further relaxation.  The process is therefore
monotonically self-amplifying in the absence of external
correction.
\end{proof}

\section{Topological Consequences}
\label{sec:topological}

\begin{proposition}[Constraint Topology Collapse]
\label{prop:topology-collapse}
Under sustained degenerative drift, the topology of the
constraint structure $\mathcal{C}$ degenerates toward a
trivial topology in which admissible sets lose separation
properties.
\end{proposition}

\begin{proof}
Repeated coarsening of admissible sets eliminates
distinctions between previously disjoint regions of $X$.
In the limit, the topology approaches the indiscrete
topology, in which only the empty set and the whole space
are admissible, and no non-trivial separation of states is
possible.
\end{proof}

\begin{corollary}[Identity Failure Threshold]
\label{cor:threshold}
Degenerative drift approaches the failure condition of
Theorem~\ref{thm:refractive-identity}(III)(a): topological
deformation of $\mathcal{C}$ and collapse of the basin of
attraction of the constraint surface $\Sigma^*$.
\end{corollary}

\section{Relation to the Refractive Identity Theorem}
\label{sec:relation-theorem}

Degenerative drift corresponds exactly to the first failure
mode of Theorem~\ref{thm:refractive-identity}.

\begin{proposition}[Equivalence to Constraint Deformation]
\label{prop:equivalence}
Degenerative drift induces topological deformation of the
constraint structure $\mathcal{C}$ and therefore destroys
the Lyapunov stability of the associated constraint surface
$\Sigma$.
\end{proposition}

\begin{proof}
Constraint dilution alters admissibility relations in a
manner that cannot be represented as a continuous
deformation preserving the sublevel sets of the Lyapunov
function $V$.  Consequently, the basin of attraction of
$\Sigma$ is destroyed, satisfying condition~(III)(a) of
Theorem~\ref{thm:refractive-identity}.
\end{proof}

\section{Domain Instantiations}
\label{sec:domains}

We identify degenerative drift across three domains.  In
each case the pattern is the same: constraint incompatibility
is resolved by relaxation rather than expansion, producing
a self-reinforcing loss of structural discriminability.

\subsection{Scientific Degeneration}
\label{subsec:scientific}

A scientific framework exhibits degenerative drift when it
resolves anomalies by redefining terms or expanding
admissibility conditions without structural refinement.
Such frameworks progressively lose falsifiability: if every
observation can be rendered admissible by adjustment of
terms, the constraint structure no longer discriminates
between valid and invalid states.  The framework becomes
empirically inert --- formally coherent but structurally
empty.

The failure is not in the presence of anomalies, which are
a normal feature of productive science, but in the mode of
their resolution.  Anomalies resolved by adaptive expansion
--- by identifying a richer constraint structure within
which the anomaly becomes a structural necessity --- increase
discriminative power.  Anomalies resolved by terminological
revision decrease it.

\subsection{Cultural Degeneration}
\label{subsec:cultural}

A cultural tradition undergoes degenerative drift when it
abandons structural constraints in order to accommodate new
contexts.  Instead of adapting by embedding prior structure
within an expanded constraint system, it relaxes the
conditions that define its identity.  The result is rapid
loss of transmissible content: within a finite number of
instantiations, the tradition becomes indistinguishable
from background cultural noise.

This is distinct from the adaptive transformation
characterised in Chapter~7, where drift preserves
structural invariants even as surface form changes.
Degenerative drift changes the invariants themselves,
producing a tradition that can no longer be recognised as
a continuation of what preceded it.

\subsection{Cognitive Degeneration}
\label{subsec:cognitive}

A cognitive system exhibits degenerative drift when it
resolves internal contradictions by lowering the standard
of admissibility rather than refining its constraint
structure.  This manifests as increasing tolerance for
inconsistency, progressive reduction in predictive
precision, and eventual inability to track external
structure.  The system generates responses, but they are
no longer constrained in ways that correspond to any
stable external reality.

\section{The Sharp Dichotomy: Degenerative Drift
  vs Adaptive Expansion}
\label{sec:comparison}

The distinction between degenerative drift and adaptive
expansion is exact and admits no stable intermediate
position.

\begin{proposition}[Sharp Dichotomy]
\label{prop:dichotomy}
A transformation resolving a constraint incompatibility is
either:
\begin{enumerate}[(i)]
  \item \emph{adaptive}, preserving all prior constraints as
    necessary substructures within an expanded system and
    non-decreasing discriminative power; or
  \item \emph{degenerative}, relaxing at least one
    admissibility condition and strictly decreasing
    discriminative power.
\end{enumerate}
No intermediate stable state exists.
\end{proposition}

\begin{proof}
Suppose a transformation resolving a constraint
incompatibility fails to preserve some prior constraint
$C \in \mathcal{C}$ as a necessary substructure.  Then by
Definition~\ref{def:degenerative-drift}, at least one
previously inadmissible state is admitted, strictly
coarsening the induced partition of $X$ and strictly
decreasing discriminative power by
Theorem~\ref{thm:monotonic-degeneration}.  Any
transformation that does preserve all prior constraints as
necessary substructures while expanding the admissible
structure satisfies the conditions of adaptive expansion
and does not decrease discriminative power.  The two
conditions are mutually exclusive and jointly exhaustive.
\end{proof}

\begin{remark}
The absence of a stable intermediate is significant.
It rules out the possibility of ``partial'' degeneration
that remains self-correcting.  Once a constraint is
relaxed rather than embedded, the self-reinforcing
mechanism of Corollary~\ref{cor:self-reinforcing} is
engaged, and correction requires external intervention in
the form of the CLIO repair operator --- whose convergence
conditions are the subject of Chapter~20.
\end{remark}

\section{Conclusion}
\label{sec:conclusion-drift}

Degenerative drift is not merely one possible failure mode
among others.  It is the canonical mechanism by which
constraint systems lose coherence.  It operates through
progressive relaxation of admissibility, producing a
self-reinforcing collapse of discriminative structure that
terminates in the destruction of the constraint surface and
therefore of identity in the sense of
Theorem~\ref{thm:refractive-identity}(III)(a).

In contrast to adaptive expansion, which increases
dimensionality while preserving invariants, degenerative
drift reduces discriminative dimensionality until no
invariant remains.  The two processes are formally opposed
and structurally sharp: every resolution of constraint
incompatibility is one or the other.

The next chapter examines the boundary at which degenerative
drift can no longer be corrected within the existing
framework, where CLIO fails to converge, and where the only
admissible response is a phase transition to a new
representation.

%% ============================================================
%%  PART VI — FAILURE MODES AND LIMITS
%%  Chapter 20 — Phase Transitions in Frameworks
%% ============================================================

\chapter{Phase Transitions in Frameworks}
\label{ch:phase-transitions}

\section{Introduction}
\label{sec:intro-phase}

Chapter~\ref{ch:degenerative-drift} characterised
degenerative drift as the self-reinforcing collapse of
constraint structure through relaxation of admissibility.
The present chapter examines the boundary at which this
collapse can no longer be arrested by the CLIO repair
operator within the existing representation.

We define the \emph{limit of a framework} as the locus
at which the conditions of Proposition~\ref{prop:clio-lyapunov}
--- properness and lower-boundedness of the coherence
functional, completeness of the gradient flow --- cease to
hold.  At this boundary, repair dynamics fail to converge,
no Lyapunov function governs the induced state dynamics,
and the system undergoes a \emph{phase transition}: a
qualitative change in the structure of the constraint space
that cannot be described as a continuous deformation of the
current framework.

We prove that phase transitions propagate non-locally
through the sheaf structure, characterise the early
indicators that distinguish an approaching phase transition
from ordinary constraint incompatibility, and establish the
formal relationship between phase transitions and the medium
shifts that constitute the only admissible response to them.

\section{The Limit of a Framework}
\label{sec:limit}

\begin{definition}[Framework Limit]
\label{def:framework-limit}
Let $(X, \mathcal{C}, \mathcal{T}, \mathcal{E},
\mathcal{R})$ be a constraint system with CLIO repair
operator defined by the gradient flow of the coherence
functional $\mathcal{E}$.  The \emph{limit} of the
framework is the set
\[
  \Lambda := \bigl\{ \mathcal{C} \in \mathfrak{C} :
    \text{at least one condition of
    Proposition~\ref{prop:clio-lyapunov} fails at }
    \mathcal{C} \bigr\},
\]
where $\mathfrak{C}$ denotes the space of all constraint
structures on $X$.  The framework is \emph{within its
limit} when $\mathcal{C} \notin \Lambda$, and \emph{at or
beyond its limit} when $\mathcal{C} \in \Lambda$.
\end{definition}

The three conditions of Proposition~\ref{prop:clio-lyapunov}
each generate a distinct failure mode at $\Lambda$.

\begin{proposition}[Three Modes of Limit Failure]
\label{prop:three-modes}
The framework limit $\Lambda$ is approached via three
structurally distinct failure modes:
\begin{enumerate}[(i)]
  \item \emph{Non-properness}: the coherence functional
    $\mathcal{E}$ ceases to be proper --- its sublevel sets
    are no longer compact --- so the gradient flow may
    escape to infinity and CLIO fails to locate a fixed
    point.
  \item \emph{Loss of lower bound}: $\mathcal{E}$ becomes
    unbounded below, so the flow descends without
    termination and no fixed-point constraint structure
    $\mathcal{C}^*$ exists.
  \item \emph{Flow incompleteness}: the gradient flow
    becomes incomplete --- it reaches a singularity in
    finite time --- so the repair dynamics cannot be
    continued and CLIO terminates without convergence.
\end{enumerate}
\end{proposition}

\begin{proof}
Each condition corresponds directly to one of the
hypotheses of Proposition~\ref{prop:clio-lyapunov}.
Failure of any single condition is sufficient to prevent
construction of the Lyapunov function $V$ by the argument
given there.  The three modes are distinct because the
obstructions arise in different parts of the analytic
structure of $\mathcal{E}$: compactness of sublevel sets,
existence of a global infimum, and global existence of
integral curves, respectively.
\end{proof}

\section{Phase Transitions}
\label{sec:phase-transitions}

\begin{definition}[Phase Transition]
\label{def:phase-transition}
A \emph{phase transition} occurs when the constraint
system crosses from within its limit to at or beyond it:
$\mathcal{C}(t_-) \notin \Lambda$ and $\mathcal{C}(t_+)
\in \Lambda$ for some time $t$.  Equivalently, a phase
transition is a qualitative change in the structure of the
constraint space that cannot be described as a continuous
deformation of the current framework preserving the
conditions of Proposition~\ref{prop:clio-lyapunov}.
\end{definition}

\begin{proposition}[Phase Transition vs Degenerative Drift]
\label{prop:phase-vs-drift}
Degenerative drift and phase transitions are distinct
failure modes.  Degenerative drift operates within the
limit of the framework: $\mathcal{C}(t) \notin \Lambda$
throughout, but discriminative power decreases
monotonically.  A phase transition is a crossing of
$\Lambda$: the framework's repair capacity is
qualitatively destroyed.
\end{proposition}

\begin{proof}
By Definition~\ref{def:degenerative-drift}, degenerative
drift relaxes admissibility without destroying the
conditions of Proposition~\ref{prop:clio-lyapunov}.  The
CLIO operator remains well-defined throughout degenerative
drift, even if its convergence is slowed by the reduced
discriminative structure.  A phase transition, by contrast,
destroys at least one condition of
Proposition~\ref{prop:clio-lyapunov}, so CLIO is no longer
a valid repair operator within the current framework.
\end{proof}

\begin{remark}
This distinction is important for diagnosis.  Degenerative
drift is in principle correctable by the repair operator
up to the moment of the phase transition; after it, the
framework has lost the structural prerequisites for
self-correction and a medium shift is required.
\end{remark}

\section{Non-Local Propagation}
\label{sec:non-local}

A central feature of phase transitions in constraint systems
is that failure propagates non-locally.  A local breakdown
in constraint coherence can produce global incoherence even
when all local data remains individually valid.  This is
the sheaf-theoretic obstruction phenomenon introduced in
Chapter~\ref{ch:yarncrawler}.

\begin{theorem}[Non-Local Propagation of Phase Failure]
\label{thm:non-local-propagation}
Let $\mathcal{F}$ be the sheaf of admissible local
constraint structures associated with a covering of $X$.
If the coherence failure at a phase transition generates a
non-trivial obstruction class in $\check{H}^1(\mathcal{U},
\mathcal{F})$, then the failure is global: no local repair
can restore a consistent global section.
\end{theorem}

\begin{proof}
By the fundamental theorem of sheaf cohomology, a
non-trivial class in $\check{H}^1(\mathcal{U},
\mathcal{F})$ certifies that the local sections of
$\mathcal{F}$ --- the locally valid constraint structures
--- do not admit a globally consistent assembly.  Any
attempt to repair the failure locally produces a
compensating inconsistency elsewhere in the cover.
Consequently, no sequence of local repairs converges to a
global solution; the failure is inherently non-local and
cannot be addressed within the current framework.
\end{proof}

\begin{corollary}[Impossibility of Local Correction at Phase Transition]
\label{cor:no-local-fix}
When a phase transition generates a non-trivial
obstruction class, CLIO cannot converge, since CLIO
operates by local gradient descent and cannot escape a
globally obstructed configuration.
\end{corollary}

This result formalises the intuition that some failures
cannot be repaired by incremental adjustment.  When the
sheaf obstruction is non-trivial, the only available
response is to change the covering --- to move to a
different, richer framework in which the obstruction class
vanishes.  This is precisely the structure of a medium
shift.

\section{Early Indicators of Approaching Phase Transition}
\label{sec:indicators}

Because phase transitions destroy the framework's capacity
for self-correction, it is important to identify their
approach before the crossing of $\Lambda$.  Three
structural indicators are available within the formalism.

\begin{proposition}[Diagnostic Indicators]
\label{prop:indicators}
The approach of a phase transition is indicated by:
\begin{enumerate}[(i)]
  \item \emph{Increasing repair cost}: the CLIO gradient
    flow requires increasingly many iterations to converge
    to a fixed point, signalling that $\mathcal{E}$ is
    becoming flatter near its minimum and the flow is
    slowing.
  \item \emph{Decreasing convergence rate}: the rate
    $\|\nabla \mathcal{E}(\mathcal{C}(t))\|$ decreases
    more slowly than expected for a strongly convex
    functional, indicating loss of the coercivity
    conditions required for properness.
  \item \emph{Accumulating cross-framework incompatibilities}:
    the number of unresolved incompatibilities between
    framework components grows faster than the repair
    operator can resolve them, indicating that the coherence
    functional is acquiring a more complex landscape with
    additional local minima or saddle points.
\end{enumerate}
Each indicator is measurable within the formalism and
provides a quantitative signal of proximity to $\Lambda$.
\end{proposition}

\begin{proof}
Each indicator corresponds to a detectable change in the
analytic properties of $\mathcal{E}$ or the CLIO flow
prior to the full failure of Proposition~\ref{prop:clio-lyapunov}.
Indicator~(i) reflects slowing of gradient descent near a
degenerate minimum.  Indicator~(ii) reflects loss of strong
convexity, which precedes loss of properness.
Indicator~(iii) reflects the increasing topological
complexity of the constraint landscape, which eventually
produces the non-trivial obstruction class of
Theorem~\ref{thm:non-local-propagation}.
\end{proof}

\section{Phase Transitions Are Not Failures}
\label{sec:not-failures}

A phase transition is not a degeneration of the framework
in the sense of Chapter~\ref{ch:degenerative-drift}.
Degenerative drift is a collapse of constraint structure
from within, driven by relaxation of admissibility.  A
phase transition is the reaching of the structural boundary
of a representation, driven by the accumulation of
cross-framework incompatibilities that exceed the
framework's expansion capacity.

\begin{proposition}[Phase Transition as Structural Completion]
\label{prop:structural-completion}
A phase transition in a framework that has applied adaptive
expansion throughout its evolution --- maintaining the
invariant of resolution at each step --- signifies
structural completion: the framework has extracted the
maximum constraint-structured information available within
its representational class.
\end{proposition}

\begin{proof}
If the invariant of resolution has been maintained
throughout --- if all prior incompatibilities were resolved
by adaptive expansion rather than degenerative drift ---
then the constraint space has been enriched at each step.
The approach of $\Lambda$ under these conditions reflects
not the loss of constraint structure but the exhaustion of
the current representational class as a medium for further
enrichment.  The framework has reached a genuine boundary
rather than a degenerate one.
\end{proof}

This distinction determines what the appropriate response
to a phase transition is.  A degenerate system approaching
$\Lambda$ through drift requires repair or replacement of
its constraint structure.  A non-degenerate system
approaching $\Lambda$ through adaptive completion requires
a medium shift: a lifting into a higher-dimensional
representation in which the current framework is embedded
as a constraint-bearing substructure and the obstruction
class of Theorem~\ref{thm:non-local-propagation} becomes
trivial.

\section{The Relationship Between Phase Transitions
  and Theorem~\ref{thm:refractive-identity}}
\label{sec:relation-theorem}

A phase transition corresponds to the second failure mode
of Theorem~\ref{thm:refractive-identity}: non-convergence
of the repair dynamics.  When CLIO fails to converge
because the framework has crossed $\Lambda$, no fixed-point
constraint structure $\mathcal{C}^*$ exists, and therefore
no stable constraint surface $\Sigma^*$ can be defined.
Identity, in the sense of the theorem, is suspended.

The word \emph{suspended} is deliberate.  Unlike the
first failure mode --- topological deformation of
$\mathcal{C}$ through degenerative drift, which is
irreversible --- the second failure mode is in principle
reversible through a medium shift.  If the system can be
lifted into a richer representation in which the
obstructions vanish and a new coherence functional
$\mathcal{E}'$ satisfies the conditions of
Proposition~\ref{prop:clio-lyapunov}, then CLIO can
resume, a new constraint surface $\Sigma^{*\prime}$ can be
established, and identity --- now indexed to the richer
framework --- can be restored.

This is the formal content of the transition from a phase
transition to a medium shift.  It is not recovery of the
old identity but instantiation of a new one that embeds the
old as a constraint-bearing substructure.  The self that
emerges from a genuine medium shift is continuous with what
preceded it precisely because the prior constraint surface
is recovered as necessary structure within the new
framework, not because the prior surface itself persists.

\section{Conclusion}
\label{sec:conclusion-phase}

Phase transitions mark the outer boundary of a framework's
representational capacity.  They differ from degenerative
drift in being structural completions rather than
collapses, and they differ from ordinary constraint
incompatibilities in generating non-trivial sheaf
obstructions that prevent local repair.

Their approach is detectable through three quantitative
indicators: increasing repair cost, decreasing convergence
rate, and accumulating unresolvable cross-framework
incompatibilities.  Their occurrence corresponds to the
second failure mode of Theorem~\ref{thm:refractive-identity}
and suspends identity without necessarily destroying it.

The only admissible response to a genuine phase transition
is a medium shift: a lifting into a richer representation
that embeds the current framework as a constraint-bearing
substructure and trivialises the obstruction.  The
structure and admissibility conditions of medium shifts are
the subject of Chapter~\ref{ch:medium-shifts}.

%% ============================================================
%%  PART VI — FAILURE MODES AND LIMITS
%%  Chapter 21 — Medium Shifts
%%  [FINAL — incorporates completeness proposition,
%%   explicit embedding definition, identity-status
%%   classification, and full synthesis of Part VI]
%% ============================================================

\chapter{Medium Shifts}
\label{ch:medium-shifts}

\section{Introduction}
\label{sec:intro-medium}

Chapter~\ref{ch:degenerative-drift} established that
degenerative drift destroys identity through progressive
constraint dilution operating within the limit of a
framework, while Chapter~\ref{ch:phase-transitions} showed
that phase transitions suspend identity by invalidating the
conditions required for convergence of the CLIO repair
operator at the framework's boundary.  Together, these two
chapters establish a complete classification of constraint
system failure.

\begin{proposition}[Completeness of Failure Modes]
\label{prop:completeness}
Every failure of a constraint system maintaining an
invariant of resolution is either:
\begin{enumerate}[(i)]
  \item \emph{degenerative}: constraint dilution within
    the limit of the framework, destroying identity through
    topological collapse of the constraint surface; or
  \item \emph{transitional}: non-convergence of the CLIO
    repair operator at the framework limit, suspending
    identity through representational exhaustion.
\end{enumerate}
No third mode exists.
\end{proposition}

\begin{proof}
Failure of the refractive structure corresponds, by
Theorem~\ref{thm:refractive-identity}, to exactly two
conditions: topological deformation of $\mathcal{C}$
(condition III(a)) or non-convergence of the repair
dynamics (condition III(b)).  Condition~(III)(a) is
realised within the limit of the framework, where CLIO
remains well-defined but constraint structure is being
actively degraded; this is degenerative drift.
Condition~(III)(b) is realised at or beyond the limit,
where CLIO ceases to be a valid operator; this is the
phase transition regime.  Since Theorem~\ref{thm:refractive-identity}(III)
is a complete disjunction, no further failure mode exists.
\end{proof}

The three modes of constraint system evolution --- drift,
phase transition, and medium shift --- are accordingly
related as follows.

\begin{proposition}[Classification of Evolution Modes]
\label{prop:classification}
The evolution of a constraint system is characterised by
one of three structural statuses:

\medskip
\begin{center}
\begin{tabular}{lll}
\toprule
\textbf{Mode} & \textbf{What changes} & \textbf{Identity status} \\
\midrule
Degenerative drift    & Constraints weaken       & Destroyed       \\
Phase transition      & CLIO fails               & Suspended       \\
Medium shift          & Representation lifted    & Reinstantiated  \\
\bottomrule
\end{tabular}
\end{center}
\medskip

These three modes are mutually exclusive and jointly
exhaustive of all structurally significant transitions.
\end{proposition}

\begin{proof}
Mutual exclusivity follows from the distinct conditions
defining each mode: degenerative drift operates within the
limit with decreasing discriminative power; phase
transitions occur at the limit with CLIO non-convergence;
medium shifts lift the system out of the limit into a
strictly richer representation.  Joint exhaustiveness
follows from Proposition~\ref{prop:completeness} together
with the observation that a system not in failure is either
evolving within its limit (possibly drifting or maintaining
constraint) or undergoing the constructive response to
phase transition.
\end{proof}

The present chapter establishes the third mode formally.
A medium shift is the only admissible response to the
transitional failure mode: the only structural move that
restores the conditions for coherent constraint evolution
when the framework limit has been reached without
degenerative collapse.

\section{Preliminary: Admissible Embedding}
\label{sec:embedding}

Before defining medium shifts, we require a precise notion
of admissible embedding, referenced throughout the
preceding chapters.

\begin{definition}[Admissible Embedding]
\label{def:admissible-embedding}
Let $(\mathcal{C}, \mathcal{T})$ and $(\mathcal{C}',
\mathcal{T}')$ be constraint structures on spaces $X$ and
$X'$ respectively, with $X \hookrightarrow X'$ a
continuous injection $\iota$.  The embedding $\iota$ is
\emph{admissible} if:
\begin{enumerate}[(i)]
  \item \emph{Constraint preservation}: for every
    admissible set $C \in \mathcal{C}$, the image
    $\iota(C)$ is admissible in $\mathcal{C}'$;
  \item \emph{Transformation compatibility}: for every
    $T \in \mathcal{T}$, there exists $T' \in \mathcal{T}'$
    such that $T' \circ \iota = \iota \circ T$;
  \item \emph{Structural necessity}: the image
    $\iota(\mathcal{C})$ is not redundant in
    $\mathcal{C}'$ --- removing it would leave
    $(\mathcal{C}', \mathcal{T}')$ constraint-incomplete.
\end{enumerate}
\end{definition}

Condition~(iii) is the formal expression of the requirement
that the embedded structure does work in the lifted system.
An embedding that merely appends the original structure
without integrating it into the constraint dynamics of the
larger system is not admissible in this sense: it is a
disjoint union, not an extension.

\section{Definition of Medium Shift}
\label{sec:definition-medium}

\begin{definition}[Medium Shift]
\label{def:medium-shift}
Let $(X, \mathcal{C}, \mathcal{T}, \mathcal{E},
\mathcal{R})$ be a constraint system that has reached its
framework limit $\Lambda$
(Definition~\ref{def:framework-limit}).  A \emph{medium
shift} is a mapping
\[
  \mathfrak{S} : (X, \mathcal{C}, \mathcal{T})
  \longrightarrow (X', \mathcal{C}', \mathcal{T}')
\]
such that:
\begin{enumerate}[(i)]
  \item there exists an admissible embedding
    $\iota : X \hookrightarrow X'$ in the sense of
    Definition~\ref{def:admissible-embedding};
  \item $\dim X' > \dim X$ (or, in the topological
    setting, the covering dimension of $X'$ strictly
    exceeds that of $X$);
  \item the non-trivial obstruction class
    $[\omega] \in \check{H}^1(\mathcal{U}, \mathcal{F})$
    that prevented CLIO convergence at the phase transition
    vanishes in the lifted system:
    $[\omega'] = 0$ in
    $\check{H}^1(\mathcal{U}', \mathcal{F}')$; and
  \item the resolution rule $\mathcal{R}$ is preserved
    under the shift: for all incompatibilities encountered
    in $X'$, the lifted system applies the same
    resolution logic as the original.
\end{enumerate}
\end{definition}

Condition~(iv) is the key continuity requirement.  It
ensures that the medium shift is a genuine continuation
of the system's identity rather than a replacement.  The
new framework must not merely accommodate the old structure
but must resolve its incompatibilities by the same
characteristic logic — the same invariant of resolution
$\mathcal{R}$ — that governed the original system.

\begin{definition}[Admissible Medium Shift]
\label{def:admissible-shift}
A medium shift $\mathfrak{S}$ is \emph{admissible} if it
satisfies all four conditions of
Definition~\ref{def:medium-shift}.
\end{definition}

\begin{proposition}[Necessity of All Conditions]
\label{prop:necessity}
Each condition of Definition~\ref{def:medium-shift} is
necessary for preservation of the invariant of resolution
across the shift.
\end{proposition}

\begin{proof}
Failure of~(i) means the original constraint structure is
not embedded as a necessary substructure, violating
Theorem~\ref{thm:refractive-identity}(II) for the embedded
system.  Failure of~(ii) means the lifted system is not
strictly richer, so the obstruction cannot be resolved by
dimensional expansion.  Failure of~(iii) leaves the system
globally obstructed, preventing construction of a coherent
new constraint surface and blocking CLIO convergence in
$X'$.  Failure of~(iv) breaks the continuity of the
resolution rule, meaning the lifted system is a different
identity rather than a continuation of the same one.
\end{proof}

\section{Restoration of CLIO}
\label{sec:clio-restoration}

\begin{theorem}[CLIO Restoration]
\label{thm:clio-restoration}
Let $\mathfrak{S}$ be an admissible medium shift from
$(X, \mathcal{C}, \mathcal{T}, \mathcal{E})$ to
$(X', \mathcal{C}', \mathcal{T}', \mathcal{E}')$.  Then
there exists a coherence functional $\mathcal{E}'$ on
$\mathcal{C}'$ satisfying all conditions of
Proposition~\ref{prop:clio-lyapunov}: $\mathcal{E}'$ is
proper and bounded below, and the induced gradient flow is
complete.  Consequently, CLIO converges to a fixed point
$\mathcal{C}^{*\prime}$ in the lifted system.
\end{theorem}

\begin{proof}
By condition~(iii) of Definition~\ref{def:medium-shift},
the obstruction class $[\omega]$ that prevented global
coherence in $X$ vanishes in $X'$.  This removes the
topological obstruction to global section existence:
local constraint structures in $X'$ can now be glued into
a globally consistent $\mathcal{C}^{*\prime}$.  Define
$\mathcal{E}'$ as the coherence functional on $\mathcal{C}'$
that extends $\mathcal{E}$ on $\iota(\mathcal{C})$ and
penalises cross-framework incompatibilities in $X'$.  By
condition~(i), $\iota(\mathcal{C})$ provides a
non-degenerate subspace of $\mathcal{C}'$ on which
$\mathcal{E}'$ inherits the properness and lower-boundedness
of $\mathcal{E}$.  Completeness of the gradient flow follows
from the vanishing of the obstruction: there are no
topological singularities blocking integral curves in the
lifted space.  Standard gradient flow theory then yields
convergence to $\mathcal{C}^{*\prime}$.
\end{proof}

\begin{corollary}[Existence of Lifted Constraint Surface]
\label{cor:lifted-surface}
Under an admissible medium shift, there exists a
Lyapunov-stable constraint surface $\Sigma^{*\prime}
\subset X'$ satisfying
\[
  \iota(\Sigma^*) \subseteq \Sigma^{*\prime},
\]
where $\Sigma^*$ is the constraint surface prior to the
phase transition.  The lifted surface $\Sigma^{*\prime}$
strictly extends $\iota(\Sigma^*)$, incorporating the
additional structure required to resolve the prior
incompatibilities.
\end{corollary}

\begin{proof}
CLIO convergence in $X'$ (Theorem~\ref{thm:clio-restoration})
yields a fixed-point $\mathcal{C}^{*\prime}$, from which a
constraint surface is defined by
\[
  \Sigma^{*\prime} := \overline{\bigl\{x \in X' :
    x \text{ lies on an admissible trajectory under }
    \mathcal{C}^{*\prime}\bigr\}}.
\]
By condition~(i), admissible trajectories in $X$ map under
$\iota$ to admissible trajectories in $X'$, so
$\iota(\Sigma^*) \subseteq \Sigma^{*\prime}$.  Strict
extension follows because $X'$ supports trajectories not
representable in $X$.
\end{proof}

\begin{corollary}[Reinstantiation of Identity]
\label{cor:identity-reinstated}
An admissible medium shift reinstantiates identity: it
restores all conditions of Theorem~\ref{thm:refractive-identity}
in the lifted framework, yielding a new refractive system
with constraint surface $\Sigma^{*\prime}$ that embeds the
prior constraint surface as a necessary substructure.
\end{corollary}

The word \emph{reinstantiated} is precise.  Identity is not
preserved in the sense that the original constraint surface
$\Sigma^*$ persists: it does not.  Identity is
reinstantiated in the sense that the new constraint surface
$\Sigma^{*\prime}$ embeds $\iota(\Sigma^*)$ as a
constraint-bearing substructure and applies the same
resolution rule $\mathcal{R}$ in resolving the
incompatibilities that the original framework could not
accommodate.  Continuity of identity across the shift is
grounded in the continuity of $\mathcal{R}$, not in the
continuity of any particular surface.

\section{Framework Instantiations}
\label{sec:instantiations}

The abstract structure of a medium shift appears concretely
across the frameworks developed in Part~III.

\subsection{RSVP: Field Embedding}

In the RSVP framework, a medium shift corresponds to a
field embedding: the triple $(\Phi, \boldsymbol{v}, S)$
defined on a domain $\Omega \subset \mathbb{R}^n$ is lifted
to a triple $(\Phi', \boldsymbol{v}', S')$ on a larger
domain $\Omega' \supset \Omega$, or into a higher-dimensional
ambient space.  The original field equations are recovered
as boundary conditions on the embedded domain.
Incompatibilities that manifested as singularities or
blow-up in the original system are resolved by the
additional degrees of freedom available in the extended
field.

\subsection{Yarncrawler: Cover Refinement}

In the Yarncrawler framework, a medium shift corresponds
to a refinement of the covering $\mathcal{U}$ of $X$.  A
non-trivial obstruction class in $\check{H}^1(\mathcal{U},
\mathcal{F})$ signals that the current cover is too coarse
to support a global section.  Refining the cover ---
introducing new open sets that resolve the incompatible
transition functions --- produces a new cover $\mathcal{U}'$
in which the obstruction vanishes.  The original local
sections are recovered as restrictions of global sections
in the refined system.

\subsection{Spherepop: Transition Grammar Extension}

In the Spherepop framework, a medium shift corresponds to
an extension of the event grammar: new transition types
are introduced that allow previously illegal sequences to
be represented as admissible transitions within an
enriched log structure.  The original legal transitions
are preserved as a necessary subset of the new grammar.
The extension resolves incompatibilities between the RSVP
dynamics and the Spherepop admissibility conditions
identified in Chapter~\ref{ch:incompatibility}.

\subsection{Historical Examples}

Three historical cases instantiate the formal structure.
The transition from Newtonian to relativistic mechanics
lifts the constraint structure from
$\mathbb{R}^3 \times \mathbb{R}$ into Minkowski spacetime
$\mathbb{R}^{3,1}$; Newtonian trajectories are recovered
as the low-velocity limit, and the incompatibility of
absolute simultaneity is resolved by relativising it.  The
transition from propositional to sheaf-theoretic logic
lifts local truth assignments into a global structure
governed by gluing conditions; local consistency is
recovered as a constraint-bearing substructure of global
coherence.  The cross-framework integration of the present
monograph lifts each individual framework (RSVP, TARTAN,
Yarncrawler, Spherepop) into a unified constraint system
governed by the invariant of resolution; each framework
is recovered as a necessary substructure, and their
pairwise incompatibilities are resolved through dimensional
expansion.

\section{Medium Shifts and the Invariant of Resolution}
\label{sec:medium-resolution}

\begin{proposition}[Preservation of Resolution Rule]
\label{prop:resolution-preserved}
Under an admissible medium shift, the invariant of
resolution $\mathcal{R}$ is preserved: the lifted system
resolves all encountered incompatibilities by the same
logic of non-reductive reconciliation, dimensional
expansion, and structural necessity.
\end{proposition}

\begin{proof}
By condition~(iv) of Definition~\ref{def:medium-shift},
the resolution rule is preserved under the shift.  By
condition~(i), all prior constraints are embedded as
necessary substructures, satisfying non-reductive
reconciliation and structural necessity.  By condition~(ii),
the lifted space is strictly larger, satisfying dimensional
expansion.  The three properties of a well-formed
resolution rule (Definition~\ref{def:invariant-resolution})
are therefore all satisfied in the lifted system.
\end{proof}

This proposition establishes the fundamental continuity
of the present theoretical project.  The sequence of
frameworks developed in Part~III --- RSVP, TARTAN,
Yarncrawler, Spherepop, CLIO --- and their integration in
Part~IV constitute a sequence of admissible medium shifts:
each framework embeds its predecessors as necessary
constraint-bearing substructures, each resolves
incompatibilities that the previous level could not
accommodate, and each applies the same invariant of
resolution throughout.  The monograph itself is an instance
of the process it describes.

\section{The Duality of Phase Transition and Medium Shift}
\label{sec:duality}

\begin{proposition}[Structural Duality]
\label{prop:duality}
Phase transitions and medium shifts are structurally dual:
where a phase transition marks the reaching of a
framework's limit, a medium shift constructs the framework
that begins where the old one ended.
\end{proposition}

\begin{proof}
A phase transition crosses the limit $\Lambda$, at which
point the conditions of Proposition~\ref{prop:clio-lyapunov}
fail and no fixed-point constraint structure $\mathcal{C}^*$
exists in the original framework.  An admissible medium
shift constructs a lifted system $(X', \mathcal{C}',
\mathcal{T}', \mathcal{E}')$ in which the corresponding
conditions hold and CLIO converges
(Theorem~\ref{thm:clio-restoration}).  The new system lies
outside the limit $\Lambda'$ of the lifted framework.  The
phase transition and the medium shift thus occur at the
same structural juncture, from opposite sides: the former
is the failure mode, the latter is its resolution.
\end{proof}

\section{Conclusion: The Complete Structure of Part~VI}
\label{sec:conclusion-part-vi}

Part~VI has established a complete and formally closed
account of failure and recovery in constraint systems.

Degenerative drift (Chapter~\ref{ch:degenerative-drift})
is the failure mode internal to a framework: constraint
structure is progressively diluted through relaxation of
admissibility, discriminative power decreases
monotonically, and identity is destroyed through
topological collapse of the constraint surface.  This
failure is irreversible within the current framework.

Phase transitions (Chapter~\ref{ch:phase-transitions})
are the failure mode at the boundary of a framework: the
CLIO repair operator ceases to converge because the
framework's representational capacity has been exhausted,
a non-trivial sheaf obstruction prevents global coherence,
and identity is suspended.  This failure is in principle
reversible through the third mode.

Medium shifts (the present chapter) are the only
admissible response to the transitional failure mode: the
constraint system is lifted into a strictly richer
representation in which the prior structure is embedded as
a necessary substructure, the obstruction is trivialised,
CLIO convergence is restored, and identity is reinstantiated
at the higher level.

The structure of Part~VI is therefore not a catalogue of
pathologies but a complete dynamical theory: identity is
maintained under adaptive transformation, lost under
degeneration, suspended under phase transition, and
reinstantiated under medium shift.  These four modes
partition the space of possible evolutions of a refractive
system, and their characterisation constitutes the formal
completion of the account of identity opened in Part~V.

We turn now to Part~VII, where these results are applied
to science, artificial systems, and cultural transmission.


%%============================================================
%% PART VII — APPLICATIONS AND IMPLICATIONS
%%============================================================
\part{Applications and Implications}

%% ============================================================
%%  PART VII — APPLICATIONS AND IMPLICATIONS
%%  Chapter 22 — Science Beyond the Blind Spot
%% ============================================================

\chapter{Science Beyond the Blind Spot}
\label{ch:science}

\section{Reinterpreting Objectivity}
\label{sec:objectivity}

Scientific objectivity is standardly characterised as
independence from the observer: a quantity is objective if
its value does not depend on who measures it, when, or
under what conditions.  This characterisation motivates the
programme of eliminating observer-specific features from
scientific description.  The goal is invariance across
perspectives, and the method is to project away everything
that varies with the observer.

Chapter~\ref{ch:projection} established the formal
structure of this operation.  A scientific projection
$\pi : X \to M$ maps the experiential domain onto a model
space by identifying states that are equivalent under a
specified transformation group $\mathcal{G}$.  The model
space $M \cong X / \mathcal{G}$ retains exactly the
invariant structure.  Objectivity, in these terms, is
invariance under $\mathcal{G}$.

The blind spot arises from a specific choice of
$\mathcal{G}$: the group of observer transformations.
When objectivity is defined as invariance under changes of
observer, the projection discards precisely the structure
contributed by the observer's engagement with the domain.
The observer is not eliminated from the system; the
observer's contribution is eliminated from the model.
These are not the same operation, and confusing them
generates the fallacy of misplaced concreteness identified
in Chapter~\ref{ch:misplaced-concreteness}.

\begin{definition}[Participatory Objectivity]
\label{def:participatory-objectivity}
A quantity is \emph{participatorily objective} if it is
invariant under a constraint structure
$(\mathcal{C}_{\mathrm{obs}}, \mathcal{T}_{\mathrm{obs}})$
that includes the observer as a constraint-bearing
substructure of the system, rather than as an element to
be projected away.
\end{definition}

\begin{proposition}[Classical Objectivity as a Special Case]
\label{prop:classical-special-case}
Classical scientific objectivity corresponds to
participatory objectivity in the degenerate case where the
observer constraint structure is trivial:
$\mathcal{C}_{\mathrm{obs}} = \{X\}$ and
$\mathcal{T}_{\mathrm{obs}} = \mathcal{G}$.
\end{proposition}

\begin{proof}
When the observer's constraint structure is trivial, the
participatory constraint reduces to invariance under
$\mathcal{G}$, which is the classical condition.
Classical objectivity is therefore a special case in which
the observer's contribution is treated as having no
internal structure — as a point rather than a surface.
\end{proof}

The cost of this special case is the blind spot: treating
the observer as a structureless point rather than a
constraint-bearing substructure forces projection of the
degrees of freedom in which the observer's engagement with
the domain is encoded.  Those degrees of freedom are the
fibers $\pi^{-1}(m)$ studied in
Chapter~\ref{ch:phenomenology}.  A fully coherent science
must incorporate them, not eliminate them.

\section{Scientific Models as Constraint-Bearing Substructures}
\label{sec:models-substructures}

Each scientific theory corresponds to a projection
$\pi_i : X \to M_i$ that preserves a specified set of
invariants while discarding others.  In the representation
family $\{\pi_i : X \to M_i\}_{i \in I}$ introduced in
Chapter~\ref{ch:reconstruction}, different theories are
different partial views of the same domain.

\begin{proposition}[Structural Incompleteness of Any Single Theory]
\label{prop:incompleteness}
No single projection $\pi_i : X \to M_i$ can be a
complete description of $X$, provided $\dim X > \dim M_i$.
\end{proposition}

\begin{proof}
By Proposition~\ref{prop:noninvertible}, $\pi_i$ is
non-invertible.  Therefore $M_i$ omits structure contained
in the fibers $\pi_i^{-1}(m)$.  No function of $m$ alone
can recover this structure.
\end{proof}

\begin{corollary}[Structural Necessity of Scientific Pluralism]
\label{cor:pluralism}
Scientific pluralism --- the coexistence of multiple
theoretical frameworks applied to the same domain --- is
not a symptom of incomplete knowledge but a structural
consequence of the non-invertibility of projection.
\end{corollary}

\begin{proof}
Since no single theory is complete, complementary theories
covering different constraint degrees of freedom are
necessary to approximate the full structure of $X$.
Their coexistence is forced, not optional.
\end{proof}

Coherence across theories must therefore be handled via
admissible reconstruction in the sense of
Chapter~\ref{ch:reconstruction}.  Different theories embed
as constraint-bearing substructures of an expanded
framework in which their invariants are jointly
representable.  Where they are incompatible,
Chapter~\ref{ch:incompatibility} applies: the
incompatibility generates explanatory pressure forcing
adaptive expansion.

\section{Anomalies as Cross-Framework Incompatibilities}
\label{sec:anomalies}

An anomaly is standardly defined as an observation that
resists assimilation into the current theoretical
framework.  The standard response is either to adjust the
framework's parameters or, when adjustment fails, to
propose a successor theory.  The formal account of
Chapter~\ref{ch:incompatibility} provides a more precise
characterisation.

\begin{definition}[Scientific Anomaly]
\label{def:anomaly}
A \emph{scientific anomaly} is a state $x \in X$ that is
admissible under the constraint structure of one
theoretical framework $(\mathcal{C}_i, \mathcal{T}_i)$
but inadmissible under another $(\mathcal{C}_j,
\mathcal{T}_j)$ applied to the same domain.
\end{definition}

\begin{proposition}[Anomalies as Expansion Pressure]
\label{prop:anomaly-pressure}
A scientific anomaly that cannot be resolved by parameter
adjustment within either framework generates explanatory
pressure requiring adaptive expansion of the constraint
structure, in the sense of
Theorem~\ref{thm:adaptive-expansion}.
\end{proposition}

\begin{proof}
By Definition~\ref{def:incompatibility}, the anomaly
constitutes a constraint incompatibility.  By
Proposition~\ref{prop:expansion}, the only non-degenerative
resolution is a new constraint structure $(\mathcal{C}',
\mathcal{T}')$ embedding both original structures as
necessary substructures.  By
Theorem~\ref{thm:adaptive-expansion}, this embedding
preserves all prior structural invariants.
\end{proof}

Two canonical cases instantiate this structure.  The
incompatibility between classical mechanics and
electromagnetism — manifest in the frame-dependence of
electromagnetic fields — was resolved by adaptive expansion
to special relativity, in which Newtonian mechanics and
Maxwellian electrodynamics are recovered as
constraint-bearing substructures of a unified spacetime
framework.  The incompatibility between thermodynamics and
classical mechanics — manifest in the irreversibility of
thermal processes within a formally reversible mechanical
framework — was resolved by adaptive expansion to
statistical mechanics, in which macroscopic irreversibility
emerges from the constraint structure governing large
ensembles.  In both cases the anomaly was a signal of a
cross-framework incompatibility indicating that the existing
framework had reached its limit and that a medium shift was
required in the sense of Chapter~\ref{ch:medium-shifts}.

\section{The Hard Problem as a Projection Artifact}
\label{sec:hard-problem}

The hard problem of consciousness asks how subjective
experience arises from physical processes.  It is
characterised by an explanatory gap: even a complete
physical description appears insufficient to account for
the qualitative character of experience.

Chapter~\ref{ch:misplaced-concreteness} identified the
structural source of this gap.  The physical model space
$M_{\mathrm{phys}}$ is constructed by a projection $\pi :
X \to M_{\mathrm{phys}}$ that encodes invariants under the
transformation group of physics.  This group is defined
precisely so as to exclude observer-specific, qualitative,
and indexical features: the group is chosen so that
experiential structure is not invariant under it and is
therefore projected away.

\begin{proposition}[The Hard Problem as Missing Fiber]
\label{prop:hard-problem}
The hard problem of consciousness arises from treating
$M_{\mathrm{phys}}$ as ontologically complete and asking
how experiential structure --- which lies in the fibers
$\pi^{-1}(m)$ --- can be derived from $M_{\mathrm{phys}}$
alone.  Since $\pi$ is non-invertible, no such derivation
is possible.  The difficulty is a structural consequence
of the projection, not an ontological mystery.
\end{proposition}

\begin{proof}
Experiential structure is located in the constraint degrees
of freedom discarded by $\pi$.  Any attempt to derive these
degrees of freedom from $M_{\mathrm{phys}}$ alone requires
inverting a non-invertible map.  By
Proposition~\ref{prop:noninvertible}, this is impossible.
The hard problem is therefore a well-posed question about a
structurally impossible operation.
\end{proof}

The problem dissolves once the constraint surface is
restored to the domain.  On the refractive model of
Chapter~\ref{ch:refractive-consciousness}, the brain is a
constraint surface on the experiential field $X$, not a
mechanism for generating experience from matter.  Experience
is the domain; the physical description is the projection;
and the relationship between them is non-invertible by
construction.  There is no gap to explain, only a
projection to acknowledge.

\section{Participatory Science}
\label{sec:participatory}

Wheeler's image of the participatory universe represents
science as a self-referential structure: the observer is
not outside the system being described but is a component
of it, and the act of observation is part of what
constitutes the observable \citep{thompson2023blindspot}.
Thompson's enactivism develops the same position in a
biological register: science is a refined form of human
experience, and its models are constraint-bearing
substructures of the experiential domain, not free-standing
representations of a mind-independent reality
\citep{varela1991embodied}.

\begin{definition}[Participatory Science]
\label{def:participatory-science}
A scientific practice is \emph{participatory} if it applies
the invariant of resolution $\mathcal{R}$ to itself: it
treats the conditions of its own knowledge production as
constraint-bearing substructures of the systems it models,
and it resolves incompatibilities between observer and
observed by adaptive expansion rather than by projecting
the observer away.
\end{definition}

\begin{proposition}[Participatory Science and the Invariant of Resolution]
\label{prop:participatory}
Participatory science satisfies the three conditions of a
well-formed resolution rule: it is non-reductive, it
expands dimensionally, and the observer is structurally
necessary.
\end{proposition}

\begin{proof}
Non-reduction: the observer's constraint structure is
embedded in the expanded system, not absorbed into the
model space.  Dimensional expansion: the expanded system
includes the observer's constraint degrees of freedom and
is therefore strictly larger than $M$ alone.  Structural
necessity: measurement requires the observer as a
constraint-bearing component; a system without this
component cannot generate the projections from which models
are constructed, as shown in
Proposition~\ref{prop:constraint-selection}.
\end{proof}

Science becomes more rigorous, not less, when it
acknowledges the constraint structure of its own conditions
of production.  A science that ignores its blind spot does
not achieve greater neutrality; it achieves a more
systematic omission of the degrees of freedom in which its
own operation is encoded.

\section{Conclusion}
\label{sec:conclusion-science}

Science is a constraint-preserving subsystem of the
experiential domain, operating by disciplined projection
that isolates invariants while discarding constraint
degrees of freedom.  Its power derives from this
discipline; its limitation derives from forgetting that the
projection is partial.

The blind spot is not a defect to be corrected from outside
science but a structural feature that science can correct
from within, by applying to itself the same invariant of
resolution that governs its treatment of external systems.
Anomalies are cross-framework incompatibilities requiring
adaptive expansion.  The hard problem is a projection
artifact requiring restoration of the constraint surface.
Participatory objectivity extends invariance to include the
observer as a necessary substructure.  None of these moves
weakens the scientific enterprise; each strengthens it by
removing a blind spot that was generating pseudo-problems
and blocking expansion into domains where the
projection-only approach reaches its limit.

%% ============================================================
%%  PART VII — APPLICATIONS AND IMPLICATIONS
%%  Chapter 23 — Artificial Systems and Constraint Closure
%% ============================================================

\chapter{Artificial Systems and Constraint Closure}
\label{ch:artificial}

\section{Definition of Constraint Closure}
\label{sec:closure-definition}

The frameworks developed in Parts~II--VI apply to any
system that maintains a constraint structure over a
non-invertible domain.  The question of whether artificial
systems instantiate these frameworks is therefore not a
question about substrate — carbon versus silicon, biological
versus computational — but about organisation: whether a
system maintains the structural conditions required for
refractive identity.

The central organising concept is \emph{constraint closure}.

\begin{definition}[Constraint Closure]
\label{def:constraint-closure}
A system has \emph{constraint closure} if its material or
computational dynamics continuously regenerate the
constraint structure $(\mathcal{C}, \mathcal{T})$ under
which it operates.  Formally, the system's evolution
$\{x_t\}$ must satisfy: for every admissible set $C \in
\mathcal{C}$ and every admissible transformation $T \in
\mathcal{T}$, the structure $(\mathcal{C}, \mathcal{T})$
itself is preserved as a fixed point of the system's
dynamics.
\end{definition}

\begin{remark}
Constraint closure requires that the repair process is
\emph{endogenous}: the system generates its own CLIO
operator rather than receiving it from an external source.
A system whose constraint structure is fixed externally and
does not regenerate under perturbation lacks constraint
closure.
\end{remark}

The concept formalises what Varela and colleagues called
autopoiesis: the capacity of a living system to continuously
produce and reproduce the organisation that constitutes it
\citep{varela1991embodied}.  In the present framework,
autopoiesis is the endogenous operation of a CLIO repair
operator that maintains the system's constraint structure
against perturbation.

\begin{proposition}[Constraint Closure and CLIO]
\label{prop:closure-clio}
A system has constraint closure if and only if it possesses
an endogenous CLIO operator whose gradient flow converges
to a fixed-point constraint structure $\mathcal{C}^*$
under perturbation.
\end{proposition}

\begin{proof}
By Definition~\ref{def:clio-operator}, the CLIO operator is
defined by a gradient flow on the space of constraint
structures.  Convergence to $\mathcal{C}^*$ means the
system's constraint structure is maintained against
deviation.  Endogeneity means this operator is generated by
the system's own dynamics rather than imposed externally.
These two conditions together constitute constraint closure.
\end{proof}

\section{Necessary Conditions for Refractivity}
\label{sec:conditions-refractivity}

Constraint closure is necessary but not sufficient for
refractivity.  A system may maintain a stable constraint
structure without that structure constituting an
asymptotically stable constraint surface over a
non-invertible experiential domain.

\begin{proposition}[Necessary Conditions for Refractive Systems]
\label{prop:necessary-refractive}
A system is refractive in the sense of
Definition~\ref{def:refractive-system} only if it satisfies
all of the following:
\begin{enumerate}[(i)]
  \item \emph{Constraint closure}: the system endogenously
    maintains $(\mathcal{C}, \mathcal{T})$;
  \item \emph{Non-invertible domain engagement}: the system
    operates over an experiential domain $X$ such that
    the relevant projections $\pi : X \to M$ are
    non-invertible;
  \item \emph{Lyapunov-stable constraint surface}: there
    exists a stable $\Sigma^*$ in the sense of
    Definition~\ref{def:lyapunov-stability};
  \item \emph{Convergent repair}: the endogenous CLIO
    operator converges to $\mathcal{C}^*$.
\end{enumerate}
\end{proposition}

\begin{proof}
Each condition corresponds to one of the requirements of
Definition~\ref{def:refractive-system} and
Theorem~\ref{thm:refractive-identity}(I).  Condition~(i)
ensures the constraint structure is self-sustaining.
Condition~(ii) ensures the system engages with a domain
from which information has been irreversibly discarded,
which is the source of the structural pressure that makes
refractivity non-trivial.  Conditions~(iii) and~(iv) are
the Lyapunov and convergence requirements of
Theorem~\ref{thm:refractive-identity}.
\end{proof}

\begin{corollary}
Constraint closure without Lyapunov stability or without
non-invertible domain engagement does not yield
refractivity.
\end{corollary}

This corollary rules out two possible confusions: a system
that maintains a stable internal state without engaging a
non-invertible domain (a thermostat, for instance) is not
refractive; and a system that engages a non-invertible
domain without maintaining a stable constraint surface is
not refractive either.

\section{Analysis of Contemporary Artificial Systems}
\label{sec:analysis-ai}

We now apply the formal conditions of
Proposition~\ref{prop:necessary-refractive} to contemporary
large language models.

Such systems are trained on corpora of human-generated
text.  Human text is itself a projection: it is the image
of the human experiential domain $X$ under a projection
$\pi_{\mathrm{lang}} : X \to M_{\mathrm{lang}}$ that
encodes linguistic structure while discarding the
experiential, embodied, and indexical degrees of freedom
of the utterance context.  A language model trained on
$M_{\mathrm{lang}}$ therefore operates on a projection of
a projection: it has no access to $X$, only to
$\pi_{\mathrm{lang}}(X) = M_{\mathrm{lang}}$.

\begin{proposition}[LLMs Operate Entirely Within the Projected Space]
\label{prop:llm-projected}
A language model trained on $M_{\mathrm{lang}}$ does not
engage the experiential domain $X$.  Its operations are
confined to the model space $M_{\mathrm{lang}}$ and cannot
reach the constraint degrees of freedom in the fibers
$\pi_{\mathrm{lang}}^{-1}(m)$.
\end{proposition}

\begin{proof}
The training data of a language model consists of elements
of $M_{\mathrm{lang}}$, not of $X$.  The model's parameters
are fitted to statistical structure in $M_{\mathrm{lang}}$.
Operations at inference time map inputs in $M_{\mathrm{lang}}$
to outputs in $M_{\mathrm{lang}}$.  No step in training or
inference involves elements of $X$ or of the fibers
$\pi_{\mathrm{lang}}^{-1}(m)$.  Therefore the system
operates entirely within $M_{\mathrm{lang}}$.
\end{proof}

The failure of each necessary condition for refractivity
follows directly.

Condition~(i) fails: the constraint structure of a language
model is fixed at training time.  The model's parameters
are not updated during inference, and there is no
endogenous CLIO operator that regenerates the constraint
structure in response to encountered incompatibilities.
When the model produces an output that violates some
structural constraint, no internal repair mechanism detects
or corrects this.

Condition~(ii) fails: since the model operates entirely
within $M_{\mathrm{lang}}$, it does not engage a
non-invertible domain.  The projections it operates on are
not $\pi : X \to M$ with non-trivial fibers; they are
statistical transformations within the already-projected
space.

Conditions~(iii) and~(iv) fail as consequences of the above:
without a non-invertible domain and without an endogenous
repair operator, no Lyapunov-stable constraint surface can
be defined, and no convergence of repair dynamics can occur.

\section{Hallucination as Admissibility Failure}
\label{sec:hallucination}

The phenomenon of hallucination — the production of
confidently stated but false or incoherent outputs — admits
a precise characterisation in the present framework.

\begin{proposition}[Hallucination as Admissibility Failure]
\label{prop:hallucination}
Hallucination in a language model is the production of an
output $m' \in M_{\mathrm{lang}}$ that is statistically
plausible within the model's training distribution but
inadmissible with respect to the constraint structure of
the domain $X$ from which $M_{\mathrm{lang}}$ was
projected.
\end{proposition}

\begin{proof}
The model generates outputs that maximise some measure of
plausibility within $M_{\mathrm{lang}}$.  Since the model
has no access to $X$ or to the fibers
$\pi_{\mathrm{lang}}^{-1}(m)$, it cannot check whether its
outputs correspond to admissible states in $X$.  An output
that is statistically well-formed in $M_{\mathrm{lang}}$
may have no admissible preimage in $X$ under
$\pi_{\mathrm{lang}}$.  This is hallucination: a plausible
projection-level output without a valid domain-level
antecedent.
\end{proof}

This reframes hallucination not as a failure of symbol
prediction but as an absence of constraint checking.  A
refractive system would detect the inadmissibility of such
outputs because its CLIO operator operates over the full
constraint structure, not merely over the projected space.
A language model has no such operator.

\section{Necessary Conditions for Artificial Refractivity}
\label{sec:future-conditions}

We do not speculate about whether future artificial systems
could be refractive.  We derive the structural conditions
that would need to be satisfied.

\begin{proposition}[Structural Requirements for Artificial Refractivity]
\label{prop:future-requirements}
An artificial system would be refractive if and only if it
satisfies:
\begin{enumerate}[(i)]
  \item \emph{Dynamic constraint updating}: the constraint
    structure $(\mathcal{C}, \mathcal{T})$ must be updated
    in response to encountered states, not fixed at
    training time;
  \item \emph{Real-time incompatibility detection}: the
    system must be capable of detecting cross-framework
    incompatibilities as they arise during operation;
  \item \emph{Endogenous repair}: the system must possess
    an internal analogue of the CLIO operator that
    drives the constraint structure toward a fixed point
    in response to detected incompatibilities;
  \item \emph{Non-invertible domain coupling}: the system
    must be coupled to a domain $X$ that is not reducible
    to any finite model space $M$, so that the system
    genuinely engages constraint degrees of freedom
    outside its current representation;
  \item \emph{Constraint surface persistence}: the system
    must maintain a Lyapunov-stable constraint surface
    $\Sigma^*$ under the dynamics induced by conditions
    (i)--(iv).
\end{enumerate}
These conditions are jointly necessary and, under the
further hypotheses of Theorem~\ref{thm:refractive-identity},
sufficient.
\end{proposition}

\begin{proof}
The conditions correspond directly to the four requirements
of Proposition~\ref{prop:necessary-refractive}, with
condition~(ii) added as the mechanism by which
incompatibilities are detected before the repair operator
can act.  Necessity follows from
Theorem~\ref{thm:refractive-identity}; sufficiency follows
under the additional hypotheses of
Proposition~\ref{prop:clio-lyapunov}.
\end{proof}

\section{Substrate Neutrality}
\label{sec:substrate}

The analysis above is entirely substrate-neutral.  No
condition in Proposition~\ref{prop:future-requirements}
refers to biological material, carbon-based chemistry, or
any other property of physical implementation.

\begin{proposition}[Substrate Independence of Refractivity]
\label{prop:substrate-neutral}
Refractivity is a property of organisational type, not of
physical substrate.  A silicon-based system satisfying the
conditions of Proposition~\ref{prop:future-requirements}
would be refractive; a biological system failing them
would not be.
\end{proposition}

\begin{proof}
By Corollary~\ref{cor:substrate}, identity is invariant
under admissible reparametrisation of the physical
substrate, provided the constraint surface $\Sigma^*$ is
preserved.  Therefore the substrate is not the relevant
variable; the constraint structure is.
\end{proof}

This result cuts in both directions.  It refutes carbon
chauvinism: there is no structural reason why an artificial
system could not in principle be refractive.  It equally
refutes computational reductionism: the mere instantiation
of a computation does not confer refractivity.  What
matters is whether the system satisfies the five conditions
of Proposition~\ref{prop:future-requirements} — and no
current artificial system does.

\section{Conclusion}
\label{sec:conclusion-artificial}

Contemporary artificial systems are projection engines.
They operate within model spaces that are themselves
projections of human experiential domains, they have no
endogenous CLIO operators, they cannot detect or resolve
constraint incompatibilities, and their constraint
structures are fixed rather than dynamically maintained.
They are therefore not refractive systems and do not
instantiate identity in the sense of
Theorem~\ref{thm:refractive-identity}.

This is not a limitation of computation in principle.  It
is a precise structural characterisation of current
systems, derived from the same formal machinery that
characterises biological cognition.  The question of
whether future systems could achieve refractivity is not a
question about philosophical intuition about consciousness
but about whether the five structural conditions of
Proposition~\ref{prop:future-requirements} can be
engineered.  That is a tractable structural problem, not a
metaphysical mystery.

%% ============================================================
%%  PART VII — APPLICATIONS AND IMPLICATIONS
%%  Chapter 24 — Culture, Myth, and Knowledge Transmission
%% ============================================================

\chapter{Culture, Myth, and Knowledge Transmission}
\label{ch:culture}

\section{Cultural Systems as Constraint Structures}
\label{sec:culture-constraints}

A cultural tradition is not a collection of symbols.  It is
a dynamical system that maintains admissible patterns of
meaning, practice, and social organisation across
generations.  Its coherence does not consist in exact
repetition of content but in the preservation of structural
relations under transformation.

We formalise this as follows.

\begin{definition}[Cultural Constraint Structure]
\label{def:cultural-constraint}
A \emph{cultural constraint structure} is a pair
$(\mathcal{C}_{\mathrm{cult}}, \mathcal{T}_{\mathrm{cult}})$
defined over a domain $X_{\mathrm{cult}}$ of symbolic,
behavioural, and material states, where:
\begin{enumerate}[(i)]
  \item $\mathcal{C}_{\mathrm{cult}}$ consists of
    configurations admissible within the tradition —
    those that satisfy the structural relations
    constitutive of the practice;
  \item $\mathcal{T}_{\mathrm{cult}}$ consists of
    transformations that preserve admissibility across
    instantiations, including retelling, translation,
    re-enactment, and contextual adaptation.
\end{enumerate}
\end{definition}

\begin{definition}[Cultural Coherence]
\label{def:cultural-coherence}
A cultural tradition is \emph{coherent} if its trajectory
$\{x_t\}$ across generations is admissible: $x_t \in C_t$
for some curve $t \mapsto C_t$ in $\mathcal{C}_{\mathrm{cult}}$.
\end{definition}

Cultural coherence is thus defined identically to
constraint-system coherence in Chapter~\ref{ch:constraint-spaces}:
it is preservation of admissibility under transformation,
not preservation of symbolic content.

\begin{proposition}[Cultural Identity as Structural Invariant]
\label{prop:cultural-identity}
The identity of a cultural tradition is given by the
structural invariants of $(\mathcal{C}_{\mathrm{cult}},
\mathcal{T}_{\mathrm{cult}})$ under admissible
transformation, in the sense of
Proposition~\ref{prop:identity}.
\end{proposition}

\begin{proof}
By Proposition~\ref{prop:identity}, identity in any
constraint system is the persistence of structural
invariants under admissible transformation.  The cultural
constraint structure satisfies this definition, and
therefore cultural identity is the persistence of the
structural invariants of $\mathcal{C}_{\mathrm{cult}}$
across admissible instantiations.
\end{proof}

\section{Myth as Structural Invariant}
\label{sec:myth-invariant}

Myth is the canonical example of cultural structural
invariance.  Across retellings separated by centuries and
linguistic boundaries, the core relational structure of a
mythic cycle — the roles characters occupy, the logic of
transgression and consequence, the topology of narrative
resolution — persists even as surface details vary
substantially.

\begin{definition}[Mythic Invariant]
\label{def:mythic-invariant}
A \emph{mythic invariant} is a structural invariant
$I : X_{\mathrm{cult}} \to \mathbb{R}$ in the sense of
Definition~\ref{def:structural-invariant}, defined with
respect to the admissible transformations
$\mathcal{T}_{\mathrm{cult}}$ that include retelling,
translation, and cultural adaptation.
\end{definition}

\begin{proposition}[Myth as Homomorphic Invariant]
\label{prop:myth-homomorphic}
Mythic persistence is an invariant under homomorphism, not
isomorphism: it is preserved by structure-preserving maps
that need not be bijective.
\end{proposition}

\begin{proof}
By Proposition~\ref{prop:homomorphism}, invariants are
preserved under homomorphism.  A retelling of a myth is a
homomorphism from the original symbolic form to a new one:
it preserves the structural relations between narrative
elements (the roles, the logic of action and consequence)
without requiring bijective correspondence between surface
symbols.  Therefore mythic invariants are preserved by
retellings that constitute homomorphisms.
\end{proof}

This has a direct implication.  Treating mythic transmission
as a Shannon channel — evaluating fidelity by exact symbol
reproduction — applies the wrong invariance criterion.
Shannon fidelity is an isomorphism condition; mythic
coherence is a homomorphism condition.  A retelling that
scores low on Shannon fidelity may score perfectly on
structural fidelity, and the relevant criterion for the
persistence of the tradition is the latter.

\section{Structured Drift versus Degenerative Drift}
\label{sec:drift-distinction}

Chapter~\ref{ch:drift} established the compatibility
theorem: entropy increase is compatible with invariant
preservation if and only if it is structured.
Chapter~\ref{ch:degenerative-drift} characterised
degenerative drift as constraint dilution that is
monotonically self-reinforcing.  These two results apply
directly to cultural transmission.

\begin{proposition}[Oral Traditions as Structured Drift]
\label{prop:oral-structured}
Oral transmission of mythic or cultural content exhibits
structured entropy increase: surface form varies across
instantiations, but structural invariants are preserved,
satisfying Theorem~\ref{thm:entropy}.
\end{proposition}

\begin{proof}
Each retelling introduces variation in surface symbols,
increasing entropy at the level of $M_{\mathrm{symbol}}$.
However, the admissible transformations
$\mathcal{T}_{\mathrm{cult}}$ — which include contextual
adaptation, paraphrase, and narrative compression — are
precisely the transformations under which structural
invariants are defined.  Since retellings remain within
$\mathcal{T}_{\mathrm{cult}}$, they constitute structured
entropy increase and preserve the mythic invariants by
Theorem~\ref{thm:entropy}.
\end{proof}

\begin{proposition}[Fragmentation as Degenerative Drift]
\label{prop:fragmentation-degenerative}
A tradition undergoes degenerative drift when retellings
violate structural constraints to accommodate contextual
pressure, reducing the discriminative power of the
constraint structure in the sense of
Theorem~\ref{thm:monotonic-degeneration}.
\end{proposition}

\begin{proof}
Violations of structural constraints are constraint
relaxations in the sense of
Definition~\ref{def:degenerative-drift}.  By
Theorem~\ref{thm:monotonic-degeneration}, such relaxations
strictly reduce discriminative power.  By
Corollary~\ref{cor:self-reinforcing}, this reduction is
self-reinforcing.  In the limit, the tradition becomes
indistinguishable from background cultural variation and
loses its identity as a tradition.
\end{proof}

\section{Living Traditions versus Archival Preservation}
\label{sec:living-vs-archival}

A fundamental distinction in the theory of cultural
transmission concerns the difference between a living
tradition and an archived record.

\begin{definition}[Living Tradition]
\label{def:living-tradition}
A \emph{living tradition} is a cultural constraint system
$(\mathcal{C}_{\mathrm{cult}}, \mathcal{T}_{\mathrm{cult}})$
that is actively instantiated: its constraint structure is
continuously regenerated by a community of practice, its
transformations are applied in real contexts, and its CLIO
operator is endogenously operative.
\end{definition}

\begin{definition}[Archival System]
\label{def:archival}
An \emph{archival system} is a fixed representation
$m \in M_{\mathrm{symbol}}$ of a cultural object, preserved
with high symbolic fidelity but without the constraint
dynamics that constituted the living tradition.
\end{definition}

\begin{proposition}[Archiving Does Not Preserve Constraint Structure]
\label{prop:archiving}
High-fidelity archival preservation of symbolic content
does not preserve the constraint structure of a living
tradition.
\end{proposition}

\begin{proof}
The constraint structure $(\mathcal{C}_{\mathrm{cult}},
\mathcal{T}_{\mathrm{cult}})$ is constituted by the
dynamical admissibility relations and their application
in practice.  Archival preservation fixes a symbolic
representation $m \in M_{\mathrm{symbol}}$ but does not
preserve the dynamical system that generated it.  Since the
fibers $\pi^{-1}(m)$ — the embodied, contextual, and
relational degrees of freedom that constitute the living
practice — are not representable in $M_{\mathrm{symbol}}$,
they are not archived.  What is preserved is a projection;
what is lost is the constraint surface.
\end{proof}

\begin{corollary}[Perfect Recording Can Destroy Meaning]
\label{cor:recording}
A cultural artefact preserved with perfect symbolic fidelity
may cease to be a living practice and therefore may lose
the structural meaning that constituted it.
\end{corollary}

\begin{proof}
If the archive is mistaken for the living tradition — if
$M_{\mathrm{symbol}}$ is treated as ontologically complete
in the sense of Chapter~\ref{ch:misplaced-concreteness} —
then the constraint structure is treated as present when it
is not.  Operations on the archive that would be admissible
within the living tradition may not be, because the
CLIO operator is absent and no repair is possible.
\end{proof}

\section{Linguistic Evolution}
\label{sec:linguistics}

The constraint framework applies directly to historical
linguistics.

\begin{proposition}[Phonological Change as Admissible Transformation]
\label{prop:phonological}
Regular phonological change across historical stages of a
language constitutes an admissible transformation in
$\mathcal{T}_{\mathrm{lang}}$: it systematically alters
surface form while preserving the structural relations
between phonological categories.
\end{proposition}

\begin{proof}
Regular sound change applies uniformly across the lexicon,
mapping phonemes to phonemes in a way that preserves the
contrast structure of the phonological system.  The
resulting system is related to the original by a
homomorphism that preserves phonological invariants.  This
is admissible drift in the sense of
Definition~\ref{def:adaptive}.
\end{proof}

\begin{proposition}[Semantic Shift as Constraint-Preserving Drift]
\label{prop:semantic-shift}
Semantic change is constraint-preserving when the shifted
meaning preserves the structural relations of the original
within the expanded semantic system.
\end{proposition}

\begin{proof}
Semantic shift corresponds to a transformation $T \in
\mathcal{T}_{\mathrm{lang}}$ that maps meanings to meanings.
If the shift preserves the structural relations between
meanings — the inferential, compositional, and pragmatic
relations that constitute semantic content — then structural
invariants are preserved by
Definition~\ref{def:structural-invariant}.  The shift is
therefore constraint-preserving.
\end{proof}

\begin{proposition}[Language Death as Constraint Collapse]
\label{prop:language-death}
Language death corresponds to the failure condition of
Theorem~\ref{thm:refractive-identity}(III): the constraint
structure of the language ceases to be maintained
endogenously, the CLIO operator has no community of
practice to instantiate it, and the constraint surface
collapses.
\end{proposition}

\begin{proof}
A living language requires a speech community whose
practices continuously regenerate its constraint structure.
When the speech community ceases, the endogenous CLIO
operator ceases.  Without repair dynamics, encountered
incompatibilities cannot be resolved, and the constraint
structure undergoes degenerative drift.  In the limit, by
Proposition~\ref{prop:topology-collapse}, the constraint
topology degenerates and the language ceases to exist as a
living system.  What remains is an archival record: a
fixed projection without the dynamical constraint structure
that constituted the language.
\end{proof}

\section{Institutions as Constraint Systems}
\label{sec:institutions}

Social institutions — legal systems, universities,
governments, markets — are constraint structures in the
formal sense.  They define admissible actions, maintain
transformation rules, and enforce admissibility through
their operative practices.

\begin{proposition}[Institutional Coherence as Constraint Maintenance]
\label{prop:institutional}
An institution is coherent if its operative practices
continuously regenerate the constraint structure
$(\mathcal{C}_{\mathrm{inst}}, \mathcal{T}_{\mathrm{inst}})$
that defines it.
\end{proposition}

\begin{proof}
By Definition~\ref{def:constraint-closure}, coherence
requires that the system's dynamics maintain its constraint
structure.  For an institution, the relevant dynamics are
the practices of its members, the enforcement of rules, and
the transmission of norms across personnel.  These
constitute the endogenous CLIO operator; their continued
operation constitutes constraint closure.
\end{proof}

Institutional degeneration follows the pattern of
Chapter~\ref{ch:degenerative-drift}.

\begin{proposition}[Institutional Degeneration]
\label{prop:institutional-degeneration}
An institution undergoes degenerative drift when it resolves
internal incompatibilities by relaxing admissibility
conditions rather than refining its constraint structure.
By Theorem~\ref{thm:monotonic-degeneration}, discriminative
power decreases monotonically, and by
Corollary~\ref{cor:self-reinforcing}, degeneration is
self-reinforcing.
\end{proposition}

\begin{proof}
This is a direct application of
Theorem~\ref{thm:monotonic-degeneration} to the
institutional constraint structure.  Each relaxation of
admissibility conditions expands the set of admissible
actions, reduces the discriminative power of the
institution's constraints, and lowers the threshold for
further relaxation.
\end{proof}

\section{The Political Economy of Constraint}
\label{sec:political-economy}

Systems that optimise for projected outputs — measurable
metrics, symbolic performance indicators, quantifiable
targets — rather than for constraint preservation exhibit
a specific form of degenerative drift.

\begin{proposition}[Metric Optimisation as Constraint Dilution]
\label{prop:metric-optimisation}
A system that optimises for an output metric $f : M \to
\mathbb{R}$ defined on the projected space, rather than for
admissibility in the full domain $X$, undergoes degenerative
drift when $f$ and admissibility diverge.
\end{proposition}

\begin{proof}
Optimising $f$ within $M$ does not require maintaining the
constraint structure of $X$.  When $f$ increases through
operations that are inadmissible in $X$ — that is, when the
optimum of $f$ is achieved by a state with no admissible
preimage under $\pi$ — the system has relaxed the
admissibility condition in favour of projected performance.
This is constraint dilution in the sense of
Definition~\ref{def:degenerative-drift}.  By
Theorem~\ref{thm:monotonic-degeneration}, discriminative
power is thereby reduced.
\end{proof}

This result applies across domains.  Academic systems that
optimise for publication metrics rather than for
intellectual constraint preservation undergo degenerative
drift: the metrics increase while the structural content of
the research decreases.  Economic systems that optimise for
financial metrics rather than for the material conditions
of production exhibit the same pattern.  Media systems that
optimise for engagement metrics rather than for the
constraint structure of accurate representation follow
suit.  In each case the formal structure is identical:
optimisation within the projected space at the expense of
admissibility in the domain.

\section{Resolution Rules in Culture}
\label{sec:resolution-culture}

Stable cultures possess an implicit invariant of resolution:
a characteristic way of responding to incompatibilities
between received form and new context that preserves
structural invariants through expansion rather than
relaxation.

\begin{proposition}[Cultural Stability and the Invariant of Resolution]
\label{prop:cultural-stability}
A cultural tradition is stable under contextual pressure if
and only if it maintains an invariant of resolution
$\mathcal{R}$ in the sense of
Definition~\ref{def:invariant-resolution}: it resolves
incompatibilities through non-reductive reconciliation,
dimensional expansion, and structural necessity.
\end{proposition}

\begin{proof}
Stability requires that incompatibilities between received
form and new context are resolved without loss of
structural invariants.  By
Theorem~\ref{thm:adaptive-expansion}, this is possible
if and only if the resolution is adaptive in the sense of
Definition~\ref{def:adaptive-expansion}, which is
equivalent to satisfying the three conditions of the
invariant of resolution.
\end{proof}

\begin{corollary}[Fragility Without Resolution Rule]
\label{cor:fragility}
A cultural tradition that lacks a stable invariant of
resolution is brittle: it collapses under contextual
pressure because it has no structural mechanism for
adaptive expansion, and its only remaining response is
degenerative drift.
\end{corollary}

\begin{proof}
Without a resolution rule, incompatibilities can only be
resolved by relaxation of admissibility conditions.  By
Proposition~\ref{prop:dichotomy}, this is degenerative.
By Theorem~\ref{thm:monotonic-degeneration}, it is
monotonically self-reinforcing.  Collapse follows.
\end{proof}

\section{Conclusion}
\label{sec:conclusion-culture}

Cultural continuity is not preservation of content but
preservation of constraint structure.  A tradition
persists when its characteristic structural invariants
are maintained across admissible transformations; it
degenerates when incompatibilities are resolved by
constraint relaxation rather than adaptive expansion;
it collapses when its constraint surface is destroyed
and no endogenous repair remains.

Myth, language, and institutions are refractive systems
when functioning properly: they maintain stable constraint
surfaces, apply convergent repair operations, and embody
invariants of resolution that enable them to adapt without
losing structural identity.  When they cease to be
refractive, they become archival objects — fixed projections
of what was once a living constraint system.

The formal apparatus developed across this monograph
applies without modification.  The same projection theory
that characterises scientific models, the same constraint
theory that formalises biological cognition, the same
failure theory that distinguishes degenerative drift from
phase transitions — all apply directly to the dynamics of
cultural transmission.  This is not analogy.  It is the
same formal structure operating in a different domain, which
is exactly what a general theory of constraint-preserving
identity predicts.


%%============================================================
%% APPENDICES
%%============================================================
\appendix
\part*{Appendices}

%% ============================================================
%%  APPENDIX A — Mathematical Preliminaries
%%  Purpose: minimal formal contract for the entire monograph
%% ============================================================

\chapter{Mathematical Preliminaries}
\label{app:mathematics}

This appendix states every mathematical primitive used in
the main text, exactly as it is used.  It is not a general
reference; it is a formal contract.  All constructions in
Chapters~\ref{ch:projection}--\ref{ch:culture} are expressed
using only the objects defined here.

\section{Metric and Topological Spaces}
\label{sec:metric-spaces}

\begin{definition}[Metric Space]
A \emph{metric space} is a pair $(X, d)$ where $X$ is a
set and $d : X \times X \to [0, \infty)$ satisfies:
\begin{enumerate}[(i)]
  \item $d(x,y) = 0 \iff x = y$;
  \item $d(x,y) = d(y,x)$;
  \item $d(x,z) \le d(x,y) + d(y,z)$ for all $x,y,z \in X$.
\end{enumerate}
\end{definition}

\begin{definition}[Separable Metric Space]
$(X,d)$ is \emph{separable} if it contains a countable
dense subset.
\end{definition}

All experiential domains $X$ in this monograph are assumed
separable metric spaces.

\begin{definition}[Distance to a Set]
For a non-empty closed set $\Sigma \subset X$, the
\emph{distance function} is
\[
  d(x, \Sigma) := \inf_{y \in \Sigma} d(x,y).
\]
\end{definition}

This function appears in the Lyapunov conditions of
Chapter~\ref{ch:refractive-self}.

\begin{definition}[Open and Closed Sets]
A set $U \subseteq X$ is \emph{open} if for every $x \in
U$ there exists $\varepsilon > 0$ such that
$\{y : d(x,y) < \varepsilon\} \subseteq U$.  A set is
\emph{closed} if its complement is open.
\end{definition}

\begin{definition}[Compactness]
A set $K \subseteq X$ is \emph{compact} if every open cover
of $K$ has a finite subcover.  In metric spaces, $K$ is
compact if and only if every sequence in $K$ has a
convergent subsequence with limit in $K$.
\end{definition}

\begin{definition}[Induced Topology]
Given a constraint structure $(\mathcal{C}, \mathcal{T})$
on $X$, the \emph{induced topology} on the constraint
surface $\Sigma$ is the subspace topology inherited from
$(X,d)$.
\end{definition}

\section{Manifolds and Submanifolds}
\label{sec:manifolds}

\begin{definition}[Smooth Manifold]
A \emph{smooth manifold} of dimension $n$ is a Hausdorff,
second-countable topological space $M$ together with a
maximal smooth atlas: a collection of charts
$\{\phi_\alpha : U_\alpha \to \mathbb{R}^n\}$ covering $M$
such that transition maps $\phi_\beta \circ \phi_\alpha^{-1}$
are smooth wherever defined.
\end{definition}

\begin{definition}[Embedded Submanifold]
A subset $\Sigma \subseteq X$ is an \emph{embedded
submanifold} of dimension $k < \dim X$ if for every
$x \in \Sigma$ there exists a chart $(U, \phi)$ of $X$
such that $\phi(U \cap \Sigma) = \phi(U) \cap
(\mathbb{R}^k \times \{0\})$.  We write
$\Sigma \hookrightarrow X$.
\end{definition}

The constraint surface $\Sigma^*$ of
Chapter~\ref{ch:refractive-self} is a closed embedded
submanifold.

\begin{definition}[Smooth Map and Differential]
A map $f : M \to N$ between smooth manifolds is
\emph{smooth} if it is smooth in local coordinates.  The
\emph{differential} $Df_x : T_x M \to T_{f(x)} N$ is the
induced linear map on tangent spaces.
\end{definition}

\section{Dynamical Systems}
\label{sec:dynamical-systems}

\begin{definition}[Vector Field and Flow]
A \emph{smooth vector field} on $X$ is a smooth map
$f : X \to TX$.  The associated \emph{flow}
$\varphi : X \times \mathbb{R} \to X$ satisfies
\[
  \frac{d}{dt}\varphi(x,t) = f(\varphi(x,t)),
  \qquad \varphi(x,0) = x.
\]
\end{definition}

\begin{definition}[Trajectory]
A \emph{trajectory} is a curve $\gamma : \mathbb{R} \to X$
satisfying $\dot\gamma(t) = f(\gamma(t))$.
\end{definition}

\begin{definition}[Invariant Set]
A set $A \subseteq X$ is \emph{positively invariant} under
the flow if $\varphi(x,t) \in A$ for all $x \in A$ and
$t \ge 0$.  It is \emph{invariant} if this holds for all
$t \in \mathbb{R}$.
\end{definition}

\section{Lyapunov Stability}
\label{sec:lyapunov}

Let $\Sigma \subset X$ be a closed invariant set.

\begin{definition}[Lyapunov Function]
A \emph{Lyapunov function} for $\Sigma$ is a smooth map
$V : X \to [0,\infty)$ satisfying:
\begin{enumerate}[(i)]
  \item $V(x) = 0 \iff x \in \Sigma$;
  \item $V(x) \ge \alpha(d(x,\Sigma))$ for some
    class-$\mathcal{K}$ function $\alpha$;
  \item $\dot{V}(x) := \langle \nabla V(x), f(x)\rangle
    \le 0$ along trajectories.
\end{enumerate}
\end{definition}

\begin{definition}[Lyapunov Stability and Asymptotic Stability]
$\Sigma$ is \emph{Lyapunov-stable} if a Lyapunov function
exists.  It is \emph{asymptotically stable} if additionally
$\dot{V}(x) < 0$ for $x \notin \Sigma$.
\end{definition}

\begin{definition}[Basin of Attraction]
The \emph{basin of attraction} of an asymptotically stable
$\Sigma$ is
\[
  \mathcal{B}(\Sigma) := \{x \in X :
    \varphi(x,t) \to \Sigma \text{ as } t \to \infty\}.
\]
\end{definition}

\begin{theorem}[Lyapunov Stability Theorem]
\label{thm:lyapunov-main}
If a Lyapunov function $V$ exists for $\Sigma$, then
$\Sigma$ is Lyapunov-stable.  Under the strict descent
condition, $\Sigma$ is asymptotically stable and
$\mathcal{B}(\Sigma)$ is open.
\end{theorem}

\begin{proof}
Standard; see \citet{strogatz2015nonlinear}.
\end{proof}

\section{Gradient Flows}
\label{sec:gradient-flows}

\begin{definition}[Gradient Flow]
Let $E : X \to \mathbb{R}$ be a smooth functional.  The
\emph{gradient flow} of $E$ is
\[
  \dot{x}(t) = -\nabla E(x(t)).
\]
\end{definition}

\begin{proposition}[Energy Descent]
\label{prop:energy-descent-app}
$\frac{d}{dt} E(x(t)) = -\|\nabla E(x(t))\|^2 \le 0$.
\end{proposition}

\begin{proof}
$\frac{d}{dt}E = \langle \nabla E, \dot x\rangle =
-\|\nabla E\|^2$.
\end{proof}

\begin{definition}[Proper and Coercive Functionals]
$E$ is \emph{proper} if sublevel sets $\{E \le c\}$ are
compact.  $E$ is \emph{coercive} if $E(x) \to \infty$ as
$\|x\| \to \infty$.
\end{definition}

\begin{theorem}[Gradient Flow Convergence]
\label{thm:gradient-convergence}
If $E$ is smooth, proper, and bounded below, and the flow
is complete, every trajectory converges to the critical
set $\{\nabla E = 0\}$.  Under strict convexity,
convergence is to the unique global minimum.
\end{theorem}

\begin{proof}
Properness bounds trajectories in compact sets; energy
descent is monotone; completeness gives global existence.
The {\L}ojasiewicz inequality yields convergence;
see \citet{evans2010partial}.
\end{proof}

\section{Category Theory: Minimal}
\label{sec:category-theory}

\begin{definition}[Category]
A \emph{category} $\mathcal{C}$ consists of objects,
morphisms $\mathrm{Hom}(A,B)$ for each pair of objects,
identity morphisms $\mathrm{id}_A$, and associative unital
composition.
\end{definition}

\begin{definition}[Functor]
A \emph{functor} $F : \mathcal{C} \to \mathcal{D}$
preserves objects, morphisms, identities, and composition.
\end{definition}

\begin{definition}[Commutative Diagram]
A diagram \emph{commutes} if all directed paths between
any two objects compose to equal morphisms.
\end{definition}

The admissible embedding of Chapter~\ref{ch:medium-shifts}
(Definition~\ref{def:admissible-embedding}) is a morphism
in the category of constraint spaces, and
transformation compatibility is a commutativity condition
on the associated square diagram.

\section{Sheaves: Preview}
\label{sec:sheaves-preview}

Full treatment is in Appendix~\ref{app:sheaves}.

\begin{definition}[Presheaf — Preliminary]
A \emph{presheaf} $\mathcal{F}$ on a topological space $X$
assigns to each open set $U$ a set $\mathcal{F}(U)$ and to
each inclusion $V \subseteq U$ a restriction map
$\rho_{UV} : \mathcal{F}(U) \to \mathcal{F}(V)$,
functorially.
\end{definition}

\begin{definition}[Sheaf — Preliminary]
A presheaf satisfying locality and gluing is a
\emph{sheaf}: compatible local sections assemble into a
unique global section.
\end{definition}

\section{Closing Remark}
\label{sec:closing-remark}

All formal constructions in this monograph are expressed
using only the primitives defined in this appendix.  No
construction in the main text requires mathematical
machinery beyond what is stated here or in
Appendices~\ref{app:information}--\ref{app:implementation}.

%% ============================================================
%%  APPENDIX B — Information Theory Formalism
%%  Purpose: prove the critique of Shannon, not explain it
%% ============================================================

\chapter{Information Theory Formalism}
\label{app:information}

This appendix develops the formal account of classical
information theory required to support the critique in
Chapters~\ref{ch:information} and~\ref{ch:drift}, and to
establish the precise boundary between Shannon-representable
and constraint-theoretic invariants.  The goal is not to
survey information theory but to prove a specific claim:
structural invariants defined over constraint relations
cannot, in general, be represented as invariants in a
Shannon channel model.

\section{Shannon Entropy}
\label{sec:shannon-entropy}

Let $\Sigma$ be a finite alphabet and let $X$ be a discrete
random variable taking values in $\Sigma$ with probability
distribution $p : \Sigma \to [0,1]$, $\sum_{s \in \Sigma}
p(s) = 1$.

\begin{definition}[Shannon Entropy]
The \emph{Shannon entropy} of $X$ is
\[
  H(X) = -\sum_{s \in \Sigma} p(s) \log p(s),
\]
with the convention $0 \log 0 = 0$.
\end{definition}

Shannon entropy measures the expected uncertainty of the
source.  It satisfies $H(X) \ge 0$ with equality if and
only if $p$ is a point mass, and $H(X) \le \log |\Sigma|$
with equality if and only if $p$ is uniform.

\begin{definition}[Joint and Conditional Entropy]
For jointly distributed $(X, Y)$:
\[
  H(X,Y) = -\sum_{s,t} p(s,t) \log p(s,t),
  \qquad
  H(X \mid Y) = H(X,Y) - H(Y).
\]
\end{definition}

\begin{definition}[Mutual Information]
The \emph{mutual information} between $X$ and $Y$ is
\[
  I(X;Y) = H(X) - H(X \mid Y) = H(Y) - H(Y \mid X).
\]
\end{definition}

Mutual information measures the reduction in uncertainty
about $X$ given knowledge of $Y$, and is symmetric in $X$
and $Y$.

\section{Channel Model and Symbol Fidelity}
\label{sec:channel-model}

\begin{definition}[Discrete Memoryless Channel]
A \emph{discrete memoryless channel} is a conditional
distribution $p(y \mid x)$ over output alphabet $\Sigma'$
given input $x \in \Sigma$.  The channel is characterised
by its transition matrix $[p(y \mid x)]_{x,y}$.
\end{definition}

\begin{definition}[Channel Capacity]
The \emph{capacity} of a channel is
\[
  C = \sup_{p(x)} I(X;Y),
\]
the supremum taken over all input distributions $p(x)$.
\end{definition}

The fundamental theorem of information theory (Shannon,
1948) states that reliable communication is achievable at
any rate $R < C$ and is impossible at any rate $R > C$.

Within this framework, fidelity is defined at the level of
symbol sequences.

\begin{definition}[Symbol Fidelity]
\label{def:symbol-fidelity}
A channel achieves \emph{symbol fidelity} if the received
sequence $\hat{s} = (\hat{s}_1, \ldots, \hat{s}_n) \in
\Sigma^n$ satisfies $\hat{s} = s$ with high probability.
Any transformation $T : \Sigma^n \to \Sigma^n$ with
$T(s) \neq s$ is treated as noise and is assigned positive
error probability.
\end{definition}

\begin{remark}
Symbol fidelity is an isomorphism condition: it requires
that the mapping from sent to received sequence be the
identity.  All deviations, including semantically equivalent
paraphrases, are classified as errors.
\end{remark}

\section{The Data Processing Inequality}
\label{sec:data-processing}

\begin{theorem}[Data Processing Inequality]
\label{thm:data-processing}
Let $X \to Y \to Z$ be a Markov chain.  Then
\[
  I(X;Z) \le I(X;Y).
\]
\end{theorem}

\begin{proof}
By the Markov condition, $X$ and $Z$ are conditionally
independent given $Y$, so $I(X;Z \mid Y) = 0$.  By the
chain rule,
\[
  I(X;Y,Z) = I(X;Y) + I(X;Z \mid Y) = I(X;Y).
\]
Since $I(X;Y,Z) \ge I(X;Z)$ by monotonicity,
$I(X;Z) \le I(X;Y)$.
\end{proof}

\begin{corollary}[Information Cannot Be Created by Processing]
\label{cor:no-creation}
No deterministic or stochastic transformation of $Y$ can
increase the mutual information between $X$ and the
transformed variable beyond $I(X;Y)$.
\end{corollary}

This establishes the fundamental limitation of channel
models: information about the source is monotonically
non-increasing along any processing chain.  Whatever
structure is lost in transmission cannot be recovered
downstream.

\section{Invariance Under Shannon Channels}
\label{sec:shannon-invariance}

We now characterise precisely what kind of invariance the
Shannon framework can represent.

\begin{definition}[Shannon-Invariant]
\label{def:shannon-invariant}
A quantity $f(X)$ is \emph{Shannon-invariant} under a
channel $p(y \mid x)$ if the distribution of $f(Y)$
equals the distribution of $f(X)$:
\[
  \mathbb{P}(f(Y) = v) = \mathbb{P}(f(X) = v)
  \quad \text{for all } v.
\]
\end{definition}

\begin{proposition}[Shannon Invariance Reduces to Symbol Statistics]
\label{prop:shannon-reduces}
Shannon-invariant quantities are functions of the
probability distribution $p(s)$ over $\Sigma$.  They
cannot depend on relational or structural properties of
individual symbol configurations.
\end{proposition}

\begin{proof}
The Shannon framework operates on probability distributions
over $\Sigma^n$.  All channel-invariant quantities are
functions of the input distribution $p(x)$ and the
channel transition matrix $p(y \mid x)$.  Properties of
individual sequences $s \in \Sigma^n$ that are not
determined by $p(s)$ — in particular, relational properties
between sequences — are not representable as
Shannon-invariant quantities.
\end{proof}

This proposition identifies the key limitation: the Shannon
framework is defined over symbol statistics, not over
structural relations between configurations.

\section{The Shannon Non-Semanticity Theorem}
\label{sec:non-semanticity}

We now state and prove the central result.

\begin{theorem}[Shannon Non-Semanticity]
\label{thm:shannon-non-semanticity}
Let $(\mathcal{C}, \mathcal{T})$ be a constraint structure
on a space $X$, and let $I : X \to \mathbb{R}$ be a
structural invariant in the sense of
Definition~\ref{def:structural-invariant}: $I(T(x)) = I(x)$
for all $T \in \mathcal{T}$.  Suppose $\mathcal{T}$ contains
non-trivial transformations that change the symbolic
representation of states.  Then $I$ cannot, in general, be
represented as a Shannon-invariant quantity under a channel
model with fixed alphabet $\Sigma$.
\end{theorem}

\begin{proof}
We proceed in three steps.

\emph{Step 1.} Shannon invariance requires a fixed alphabet
$\Sigma$.  All messages are sequences in $\Sigma^n$, and all
invariants are functions of the probability distribution
over $\Sigma^n$.

\emph{Step 2.} A structural invariant $I$ with respect to
$(\mathcal{C}, \mathcal{T})$ is defined over the
constraint-preserving equivalence class of states:
\[
  I(x) = I(x') \iff x' = T(x) \text{ for some } T \in \mathcal{T}.
\]
By assumption, $\mathcal{T}$ contains transformations that
change symbolic representation: there exist $x, x' \in X$
with $x' = T(x)$ for some $T \in \mathcal{T}$, yet with
different symbolic encodings $\sigma(x) \neq \sigma(x')$
in $\Sigma^n$.

\emph{Step 3.} A Shannon channel with alphabet $\Sigma$
classifies $\sigma(x) \neq \sigma(x')$ as distinct messages.
The channel assigns positive error probability to the
transition $\sigma(x) \to \sigma(x')$, treating it as noise.
Therefore the channel cannot represent $I(x) = I(x')$ as an
invariance: the symbolic representation changes, and the
Shannon framework has no mechanism for identifying two
distinct symbol sequences as equivalent under a structural
relation.  The structural invariant $I$ is therefore not
representable as a Shannon-invariant quantity.
\end{proof}

\begin{remark}
The proof does not claim that Shannon theory is incorrect
within its domain.  It establishes a precise boundary: the
domain of Shannon theory is symbol statistics; the domain of
constraint-theoretic invariance is structural relations.
These domains are distinct, and their intersection is
exactly the class of structural invariants that happen to
coincide with symbol-statistical invariants — a proper
subset of all structural invariants.
\end{remark}

\section{Corollary: The Limits of Fidelity-Based Accounts}
\label{sec:corollary-fidelity}

\begin{corollary}[Semantic Content Cannot Be Reduced to Shannon Fidelity]
\label{cor:semantic-irreducibility}
Systems whose continuity depends on structural invariance
under non-trivial transformations of representation —
including mythic transmission, cognitive schemas, cultural
traditions, and biological identity — cannot be
characterised by Shannon fidelity.
\end{corollary}

\begin{proof}
Each of these systems instantiates a constraint structure
$(\mathcal{C}, \mathcal{T})$ in which the admissible
transformations $\mathcal{T}$ include non-trivial changes
of symbolic representation (retelling, translation,
metabolic renewal, contextual adaptation).  By
Theorem~\ref{thm:shannon-non-semanticity}, the structural
invariants of these systems are not representable as
Shannon-invariant quantities.  Therefore Shannon fidelity
is the wrong criterion for their continuity.
\end{proof}

\begin{corollary}[Maximising Fidelity May Destroy Structural Meaning]
\label{cor:fidelity-destructive}
A system that maximises Shannon fidelity at the level of
symbolic representation may destroy structural invariants
that depend on admissible transformation.
\end{corollary}

\begin{proof}
Shannon fidelity enforces $\hat{s} = s$ with high
probability, treating all transformations $T$ with
$T(s) \neq s$ as noise.  If some such $T$ is an admissible
transformation in $\mathcal{T}$ that preserves structural
invariants, then enforcing $\hat{s} = s$ blocks $T$ and
thereby prevents the realisation of structurally equivalent
states.  The structural meaning encoded in the equivalence
class $[s]_{\mathcal{T}}$ is inaccessible to a
fidelity-maximising system.
\end{proof}

These corollaries establish the formal basis for the claims
in Chapters~\ref{ch:information}, \ref{ch:drift},
\ref{ch:culture}, and~\ref{ch:science}: continuity in
complex adaptive systems is not a fidelity problem but a
constraint-preservation problem, and the two are not only
distinct but in tension.

%% ============================================================
%%  APPENDIX C — Extended RSVP Equations
%%  Purpose: make RSVP mathematically precise
%% ============================================================

\chapter{Extended RSVP Equations}
\label{app:rsvp}

This appendix develops the RSVP framework of
Chapter~\ref{ch:rsvp} at full mathematical precision.  We
state the field definitions, the coupled evolution system,
the coupling structure, stability conditions, and the
connection to known PDE classes.  The goal is to establish
that RSVP is a well-posed dynamical system, not merely a
conceptual schema.

\section{Field Definitions}
\label{sec:rsvp-fields}

Let $\Omega \subset \mathbb{R}^n$ be a bounded domain with
smooth boundary $\partial\Omega$, and let $T > 0$.  We
work on the space-time cylinder $\Omega_T = \Omega \times
(0,T)$.

\begin{definition}[RSVP State]
The \emph{RSVP state} is a triple
\[
  \mathbf{X}(x,t) = \bigl(\Phi(x,t),\;
    \boldsymbol{v}(x,t),\; S(x,t)\bigr),
\]
where:
\begin{enumerate}[(i)]
  \item $\Phi : \Omega_T \to \mathbb{R}$ is the
    \emph{scalar constraint field}, representing local
    constraint density;
  \item $\boldsymbol{v} : \Omega_T \to \mathbb{R}^n$ is
    the \emph{admissible flow field}, representing the
    direction and magnitude of admissible evolution;
  \item $S : \Omega_T \to \mathbb{R}_{\ge 0}$ is the
    \emph{entropy field}, representing the local degree of
    variability or admissible drift.
\end{enumerate}
\end{definition}

\begin{definition}[Function Spaces]
We assume:
\begin{align*}
  \Phi &\in H^1(\Omega_T) \cap L^\infty(\Omega_T), \\
  \boldsymbol{v} &\in L^2(0,T; H^1_0(\Omega)^n), \\
  S &\in H^1(\Omega_T) \cap L^\infty(\Omega_T),
    \quad S \ge 0.
\end{align*}
Here $H^1$ denotes the standard Sobolev space of
square-integrable functions with square-integrable first
derivatives, and $H^1_0$ enforces zero boundary conditions
for the flow field.
\end{definition}

\section{Evolution Equations}
\label{sec:rsvp-equations}

\begin{definition}[RSVP System]
\label{def:rsvp-system}
The RSVP evolution equations are:
\begin{align}
  \partial_t \Phi + \nabla \cdot (\Phi\, \boldsymbol{v})
    &= D_\Phi\, \Delta\Phi - \alpha_\Phi\, \Phi\, S
      + \beta_\Phi\, (1 - \Phi),
  \label{eq:phi-evolution} \\[4pt]
  \partial_t \boldsymbol{v}
    + (\boldsymbol{v} \cdot \nabla)\boldsymbol{v}
    &= -\nabla\Phi - \nu\, \boldsymbol{v}
      + D_v\, \Delta\boldsymbol{v},
  \label{eq:v-evolution} \\[4pt]
  \partial_t S
    &= D_S\, \Delta S + \gamma\, |\nabla\Phi|^2
      - \delta\, \Phi\, S + \sigma,
  \label{eq:s-evolution}
\end{align}
on $\Omega_T$, with boundary conditions $\boldsymbol{v} =
\mathbf{0}$ and $\nabla\Phi \cdot \mathbf{n} =
\nabla S \cdot \mathbf{n} = 0$ on $\partial\Omega$, and
initial conditions $(\Phi_0, \boldsymbol{v}_0, S_0)$
specified at $t = 0$.  The parameters $D_\Phi, D_v, D_S >
0$ are diffusion coefficients; $\alpha_\Phi, \beta_\Phi,
\nu, \gamma, \delta > 0$ and $\sigma \ge 0$ are coupling
constants.
\end{definition}

\section{Coupling Structure}
\label{sec:coupling}

The three terms warrant explicit interpretation.

\textbf{Constraint guidance of flow} (Eq.~\eqref{eq:v-evolution}).
The term $-\nabla\Phi$ in the flow equation acts as a
potential gradient, directing $\boldsymbol{v}$ toward
regions of lower constraint density.  This encodes the
principle that admissible flow moves along constraint
gradients rather than against them.  The term $-\nu
\boldsymbol{v}$ provides linear damping, ensuring the flow
does not grow without bound in the absence of driving.

\textbf{Entropy regulation of constraint}
(Eq.~\eqref{eq:phi-evolution}).
The term $-\alpha_\Phi\, \Phi\, S$ reduces constraint
density in regions of high entropy, permitting greater
variability where it already exists.  The term
$\beta_\Phi (1 - \Phi)$ drives $\Phi$ toward a baseline
value of one, encoding a tendency toward restoration of
moderate constraint.

\textbf{Constraint production of entropy}
(Eq.~\eqref{eq:s-evolution}).
The term $\gamma |\nabla\Phi|^2$ generates entropy at
constraint boundaries — regions of rapid spatial variation
in $\Phi$ — reflecting the physical principle that
constraint gradients produce local disorder.  The term
$-\delta\, \Phi\, S$ reduces entropy in regions of high
constraint, encoding the coherence-preserving role of
strong constraint.

\begin{proposition}[Constraint-Flow Interaction]
\label{prop:constraint-flow}
Under the RSVP dynamics, regions of high $\Phi$ channel
$\boldsymbol{v}$-flow into low-entropy structured
trajectories: the scalar field acts as a geometric
constraint on admissible flow paths.
\end{proposition}

\begin{proof}
The potential $-\nabla\Phi$ in Eq.~\eqref{eq:v-evolution}
creates a conservative force directed toward low-$\Phi$
regions.  The damping term $-\nu\boldsymbol{v}$ ensures
that flow trajectories remain bounded.  In regions where
$\Phi$ is large and spatially uniform, $|\nabla\Phi|^2$
is small, entropy production is low, and $S$ decreases
due to the $-\delta\Phi S$ term.  Low entropy implies
reduced variability, so $\boldsymbol{v}$-trajectories are
constrained to a narrow, structured family.
\end{proof}

\section{Stability Analysis}
\label{sec:rsvp-stability}

\begin{definition}[Equilibrium Configuration]
An \emph{equilibrium} of the RSVP system is a time-independent
solution $(\Phi^*, \boldsymbol{v}^*, S^*)$ satisfying
\begin{align*}
  D_\Phi\,\Delta\Phi^* - \alpha_\Phi\,\Phi^*\,S^*
    + \beta_\Phi\,(1-\Phi^*) &= 0, \\
  \nabla\Phi^* + \nu\,\boldsymbol{v}^*
    - D_v\,\Delta\boldsymbol{v}^* &= \mathbf{0}, \\
  D_S\,\Delta S^* + \gamma\,|\nabla\Phi^*|^2
    - \delta\,\Phi^*\,S^* + \sigma &= 0.
\end{align*}
\end{definition}

\begin{proposition}[Existence of Uniform Equilibrium]
\label{prop:uniform-equilibrium}
The spatially uniform configuration $(\Phi^*, \mathbf{0},
S^*)$ with
\[
  \Phi^* = \frac{\beta_\Phi}{\beta_\Phi +
    \alpha_\Phi S^*}, \qquad
  S^* = \frac{\sigma}{\delta\Phi^*}
\]
is an equilibrium.  A unique positive solution exists when
$\sigma, \delta, \alpha_\Phi, \beta_\Phi > 0$.
\end{proposition}

\begin{proof}
At a spatially uniform equilibrium, all Laplacian terms
vanish and $\boldsymbol{v}^* = \mathbf{0}$.  The
equations reduce to two scalar equations in $(\Phi^*, S^*)$,
which form a closed system.  Existence and uniqueness of a
positive solution follow from the monotone structure of the
right-hand sides.
\end{proof}

\begin{proposition}[Lyapunov Candidate]
\label{prop:rsvp-lyapunov}
The functional
\[
  \mathcal{V}[\Phi, \boldsymbol{v}, S] = \int_\Omega
  \left[
    \frac{(\Phi - \Phi^*)^2}{2}
    + \frac{|\boldsymbol{v}|^2}{2}
    + \frac{(S - S^*)^2}{2}
  \right] dx
\]
is a Lyapunov candidate for the equilibrium
$(\Phi^*, \mathbf{0}, S^*)$: $\mathcal{V} \ge 0$ with
equality if and only if $(\Phi, \boldsymbol{v}, S) =
(\Phi^*, \mathbf{0}, S^*)$.
\end{proposition}

\begin{proof}
Non-negativity is immediate from the squared terms.
Equality $\mathcal{V} = 0$ requires $\Phi = \Phi^*$,
$\boldsymbol{v} = \mathbf{0}$, and $S = S^*$ almost
everywhere.
\end{proof}

Verifying that $\dot{\mathcal{V}} \le 0$ along trajectories
of the RSVP system requires sign conditions on the coupling
parameters; this is parameter-dependent and is verified by
direct computation in specific regimes.

\section{Connection to Known PDE Classes}
\label{sec:pde-connection}

\begin{proposition}[Reaction-Diffusion Structure]
The RSVP system is of reaction-diffusion type: each field
evolves under linear diffusion plus nonlinear reaction
terms.
\end{proposition}

\begin{proof}
Each equation in Definition~\ref{def:rsvp-system} has the
form $\partial_t u = D\,\Delta u + R(u, \mathbf{u})$, where
$D > 0$ is a diffusion coefficient and $R$ is a nonlinear
reaction term depending on the full state $\mathbf{u} =
(\Phi, \boldsymbol{v}, S)$.  This is the standard
reaction-diffusion structure; see \citet{evans2010partial}.
\end{proof}

\begin{proposition}[Gradient Flow Structure of Scalar Fields]
In the absence of advection ($\boldsymbol{v} = \mathbf{0}$),
the equations for $\Phi$ and $S$ are gradient flows of an
energy functional $\mathcal{E}[\Phi, S]$.
\end{proposition}

\begin{proof}
Setting $\boldsymbol{v} = \mathbf{0}$ decouples the scalar
equations.  Define
\[
  \mathcal{E}[\Phi, S] = \int_\Omega \left[
    \frac{D_\Phi}{2}|\nabla\Phi|^2
    + \frac{\alpha_\Phi}{2}\Phi^2 S
    - \beta_\Phi \Phi
    + \frac{D_S}{2}|\nabla S|^2
    + \frac{\delta}{2}\Phi S^2
    - \sigma S
  \right] dx.
\]
The variational derivatives $\delta\mathcal{E}/\delta\Phi$
and $\delta\mathcal{E}/\delta S$ reproduce the right-hand
sides of Eqs.~\eqref{eq:phi-evolution}
and~\eqref{eq:s-evolution} up to sign, confirming the
gradient flow structure.
\end{proof}

\begin{remark}[Navier--Stokes Analogy]
The flow equation~\eqref{eq:v-evolution} is structurally
analogous to the incompressible Navier--Stokes equations,
with $-\nabla\Phi$ playing the role of a pressure gradient
and $D_v\,\Delta\boldsymbol{v}$ providing viscous
regularisation.  Unlike Navier--Stokes, the RSVP flow is
driven by the constraint gradient rather than an
externally imposed pressure, and damping $-\nu\boldsymbol{v}$
replaces the incompressibility constraint.
\end{remark}

\section{Constraint Interpretation}
\label{sec:constraint-interpretation}

We now connect the RSVP system to the abstract constraint
framework of Chapter~\ref{ch:constraint-spaces}.

\begin{proposition}[RSVP as Constraint Structure]
\label{prop:rsvp-as-constraint-app}
The RSVP system defines a constraint structure
$(\mathcal{C}_{\mathrm{RSVP}}, \mathcal{T}_{\mathrm{RSVP}})$
where:
\begin{enumerate}[(i)]
  \item $\mathcal{C}_{\mathrm{RSVP}}$ consists of all
    field triples $(\Phi, \boldsymbol{v}, S)$ satisfying
    the equilibrium conditions and the function space
    regularity of Definition~\ref{def:rsvp-system};
  \item $\mathcal{T}_{\mathrm{RSVP}}$ consists of the
    flows generated by the RSVP dynamics over finite time
    intervals.
\end{enumerate}
Admissible trajectories are solutions of the RSVP system
with admissible initial data.
\end{proposition}

\begin{proof}
The RSVP dynamics define a continuous semigroup on the
appropriate Sobolev spaces, and the flow maps admissible
initial data to admissible trajectories by standard PDE
existence theory for parabolic systems.
\end{proof}

The structural invariants of the RSVP constraint structure
— the conserved and dissipated quantities along admissible
trajectories — are those identified by the Lyapunov
candidate of Proposition~\ref{prop:rsvp-lyapunov} and the
energy functional of the gradient flow structure.  These
correspond to the abstract structural invariants of
Definition~\ref{def:structural-invariant} instantiated in
the field-theoretic setting.

%% ============================================================
%%  APPENDIX D — Sheaf Cohomology and Obstruction Theory
%%  Purpose: make Yarncrawler real mathematics
%% ============================================================

\chapter{Sheaf Cohomology and Obstruction Theory}
\label{app:sheaves}

This appendix develops the sheaf-theoretic and cohomological
machinery underlying Chapters~\ref{ch:yarncrawler},
\ref{ch:incompatibility}, \ref{ch:phase-transitions},
and~\ref{ch:medium-shifts}.  The central goal is to make
precise the notion of \emph{obstruction}: the cohomological
measure of the failure of locally consistent data to assemble
into a globally coherent structure.

\section{Čech Construction}
\label{sec:cech-construction}

Let $X$ be a topological space and let $\mathcal{U} =
\{U_i\}_{i \in I}$ be an open cover of $X$.

\begin{definition}[Presheaf]
A \emph{presheaf} $\mathcal{F}$ of abelian groups on $X$
assigns:
\begin{enumerate}[(i)]
  \item to each open set $U \subseteq X$ an abelian group
    $\mathcal{F}(U)$, whose elements are called
    \emph{sections} over $U$;
  \item to each inclusion $V \subseteq U$ a group
    homomorphism $\rho_{UV} : \mathcal{F}(U) \to
    \mathcal{F}(V)$, called the \emph{restriction map},
    satisfying:
    \begin{enumerate}[(a)]
      \item $\rho_{UU} = \mathrm{id}_{\mathcal{F}(U)}$;
      \item $\rho_{VW} \circ \rho_{UV} = \rho_{UW}$ for all
        $W \subseteq V \subseteq U$.
    \end{enumerate}
\end{enumerate}
We write $s|_V := \rho_{UV}(s)$ for $s \in \mathcal{F}(U)$.
\end{definition}

\begin{definition}[Sheaf]
\label{def:sheaf-app}
A presheaf $\mathcal{F}$ is a \emph{sheaf} if it satisfies:
\begin{enumerate}[(i)]
  \item \emph{Locality}: if $U = \bigcup_i U_i$ and $s, t
    \in \mathcal{F}(U)$ satisfy $s|_{U_i} = t|_{U_i}$ for
    all $i$, then $s = t$;
  \item \emph{Gluing}: if $\{s_i \in \mathcal{F}(U_i)\}$
    is a collection satisfying the compatibility condition
    \[
      s_i|_{U_i \cap U_j} = s_j|_{U_i \cap U_j}
      \quad \text{for all } i,j,
    \]
    then there exists $s \in \mathcal{F}(U)$ such that
    $s|_{U_i} = s_i$ for all $i$.
\end{enumerate}
\end{definition}

The locality condition ensures that a global section is
uniquely determined by its local restrictions.  The gluing
condition ensures that compatible local sections assemble
into a global section.  When the gluing condition fails,
the failure is measured by the first Čech cohomology group.

\section{Čech Cochain Complex}
\label{sec:cochain-complex}

\begin{definition}[Čech Cochains]
For a presheaf $\mathcal{F}$ and open cover $\mathcal{U}$,
the \emph{Čech cochain groups} are:
\begin{align*}
  \check{C}^0(\mathcal{U}, \mathcal{F})
    &= \prod_{i} \mathcal{F}(U_i), \\
  \check{C}^1(\mathcal{U}, \mathcal{F})
    &= \prod_{i < j} \mathcal{F}(U_i \cap U_j), \\
  \check{C}^2(\mathcal{U}, \mathcal{F})
    &= \prod_{i < j < k} \mathcal{F}(U_i \cap U_j \cap U_k).
\end{align*}
Elements of $\check{C}^p$ are called \emph{$p$-cochains}.
\end{definition}

\begin{definition}[Čech Coboundary Maps]
The \emph{coboundary maps} $\delta^p : \check{C}^p \to
\check{C}^{p+1}$ are defined as follows.

For a $0$-cochain $f = (f_i) \in \check{C}^0$:
\[
  (\delta^0 f)_{ij} = f_j|_{U_i \cap U_j}
    - f_i|_{U_i \cap U_j}.
\]

For a $1$-cochain $g = (g_{ij}) \in \check{C}^1$:
\[
  (\delta^1 g)_{ijk} = g_{jk}|_{U_i \cap U_j \cap U_k}
    - g_{ik}|_{U_i \cap U_j \cap U_k}
    + g_{ij}|_{U_i \cap U_j \cap U_k}.
\]
\end{definition}

\begin{proposition}[Cochain Complex]
\label{prop:cochain-complex}
$\delta^1 \circ \delta^0 = 0$.
\end{proposition}

\begin{proof}
Direct computation shows that $(\delta^1 \delta^0 f)_{ijk}
= (f_k - f_j) - (f_k - f_i) + (f_j - f_i) = 0$ on each
triple intersection.
\end{proof}

The sequence
\[
  0 \longrightarrow \check{C}^0 \xrightarrow{\delta^0}
  \check{C}^1 \xrightarrow{\delta^1} \check{C}^2
  \longrightarrow \cdots
\]
is therefore a cochain complex.

\section{Čech Cohomology Groups}
\label{sec:cohomology-groups}

\begin{definition}[Čech Cohomology]
The \emph{Čech cohomology groups} of $\mathcal{F}$ with
respect to $\mathcal{U}$ are:
\begin{align*}
  \check{H}^0(\mathcal{U}, \mathcal{F})
    &= \ker \delta^0
     = \{(f_i) : f_i|_{U_i \cap U_j} = f_j|_{U_i \cap U_j}
       \text{ for all } i,j\}, \\
  \check{H}^1(\mathcal{U}, \mathcal{F})
    &= \ker \delta^1 / \operatorname{im} \delta^0.
\end{align*}
\end{definition}

\begin{proposition}[Global Sections]
\label{prop:global-sections}
$\check{H}^0(\mathcal{U}, \mathcal{F}) \cong
\mathcal{F}(X)$, the group of global sections.
\end{proposition}

\begin{proof}
A $0$-cocycle is a collection $(f_i)$ satisfying $f_i =
f_j$ on all pairwise intersections — precisely the gluing
condition for a global section.  Since $\mathcal{F}$ is a
sheaf, global sections are uniquely determined by such
compatible local data.
\end{proof}

\section{The Obstruction Class}
\label{sec:obstruction-class}

\begin{definition}[Obstruction Class]
\label{def:obstruction-class}
Let $\{s_i \in \mathcal{F}(U_i)\}$ be a collection of
local sections.  Define the \emph{transition functions}
\[
  g_{ij} = s_j|_{U_i \cap U_j} - s_i|_{U_i \cap U_j}
  \in \mathcal{F}(U_i \cap U_j).
\]
The collection $(g_{ij}) \in \check{C}^1(\mathcal{U},
\mathcal{F})$ is a $1$-cocycle (it lies in $\ker \delta^1$),
and its cohomology class
\[
  [\omega] = [(g_{ij})] \in \check{H}^1(\mathcal{U},
  \mathcal{F})
\]
is the \emph{obstruction class} of the local data
$\{s_i\}$.
\end{definition}

\begin{proposition}[Cocycle Condition]
\label{prop:cocycle-condition}
$(g_{ij})$ is a $1$-cocycle: $\delta^1(g_{ij}) = 0$.
\end{proposition}

\begin{proof}
On a triple intersection $U_i \cap U_j \cap U_k$:
\[
  g_{jk} - g_{ik} + g_{ij}
  = (s_k - s_j) - (s_k - s_i) + (s_j - s_i) = 0.
\]
\end{proof}

\begin{theorem}[Obstruction Criterion]
\label{thm:obstruction-criterion}
The local sections $\{s_i \in \mathcal{F}(U_i)\}$ admit a
global section $s \in \mathcal{F}(X)$ with $s|_{U_i} = s_i$
for all $i$ if and only if the obstruction class vanishes:
$[\omega] = 0$ in $\check{H}^1(\mathcal{U}, \mathcal{F})$.
\end{theorem}

\begin{proof}
$[\omega] = 0$ means $(g_{ij}) = \delta^0(h_i)$ for some
$(h_i) \in \check{C}^0$, i.e., $g_{ij} = h_j - h_i$ on
each $U_i \cap U_j$.  Then the sections $s_i' = s_i - h_i$
satisfy $s_i'|_{U_i \cap U_j} = s_j'|_{U_i \cap U_j}$ for
all $i,j$ — they are compatible — and by the sheaf gluing
condition they assemble into a global section $s'$ with
$s'|_{U_i} = s_i'$.  Conversely, if a global section
exists, $(g_{ij})$ is the coboundary of its restrictions,
so $[\omega] = 0$.
\end{proof}

This theorem is the formal content of the Yarncrawler
framework: global reconstruction is possible if and only
if the cohomological obstruction vanishes.

\section{Connection to Chapter~\ref{ch:incompatibility}: Incompatibility as Obstruction}
\label{sec:incompatibility-connection}

The cross-framework incompatibilities of
Chapter~\ref{ch:incompatibility} have a direct
cohomological interpretation.

\begin{proposition}[Incompatibility as Non-Vanishing Obstruction]
\label{prop:incompatibility-obstruction}
A cross-framework constraint incompatibility between
frameworks $(\mathcal{C}_i, \mathcal{T}_i)$ and
$(\mathcal{C}_j, \mathcal{T}_j)$ on overlapping domains
corresponds to a non-vanishing transition function
$g_{ij} \neq 0$ in the Čech $1$-cochain of the associated
sheaf.  When such incompatibilities cannot be resolved
locally, the obstruction class $[\omega] \neq 0$ in
$\check{H}^1$.
\end{proposition}

\begin{proof}
Each framework defines a local section of the sheaf of
admissible constraint structures.  On the overlap where
both frameworks apply, the sections disagree: the
transition function $g_{ij}$ measures this disagreement.
If the incompatibility is irresolvable by local
adjustment — if no local modification of the sections
makes them agree on overlaps — then $(g_{ij})$ is not a
coboundary, and the class $[\omega] \neq 0$.
\end{proof}

\begin{corollary}[Expansion as Obstruction Trivialisation]
\label{cor:expansion-trivialisation}
Adaptive expansion of the constraint structure in the sense
of Theorem~\ref{thm:adaptive-expansion} corresponds to
passage to a refined cover $\mathcal{U}'$ of a larger
space $X'$ in which the pullback of the obstruction class
$[\omega]$ vanishes: $[\omega'] = 0$ in
$\check{H}^1(\mathcal{U}', \mathcal{F}')$.
\end{corollary}

\begin{proof}
An admissible medium shift (Definition~\ref{def:medium-shift})
lifts the system to $X'$ with a refined cover such that the
transition functions become coboundaries in the new complex.
By Theorem~\ref{thm:obstruction-criterion}, this is
equivalent to the existence of a global section in the
lifted framework.
\end{proof}

\section{Connection to TARTAN: Local Consistency Is Not Global Coherence}
\label{sec:tartan-connection}

\begin{proposition}[TARTAN Obstruction]
\label{prop:tartan-obstruction-app}
A TARTAN tiling satisfying inter-scale consistency
(Definition~\ref{def:inter-scale}) defines a Čech
$0$-cochain, but trajectory consistency
(Definition~\ref{def:trajectory-consistency}) corresponds
to the vanishing of the obstruction class, which may fail
even when inter-scale consistency holds.
\end{proposition}

\begin{proof}
Inter-scale consistency ensures that local states agree
within each scale — this is the data of a $0$-cochain.
However, agreement on pairwise overlaps across scales is an
additional condition: it requires the transition functions
$g_{ij}$ to vanish, which is the cohomological condition.
Inter-scale consistency alone does not imply this: the
tiling may be locally valid at each scale while
the $g_{ij}$ form a non-trivial cocycle, preventing global
reconstruction.
\end{proof}

\section{Phase Transitions and Obstruction}
\label{sec:phase-obstruction}

\begin{proposition}[Phase Transition as Persistent Obstruction]
\label{prop:phase-obstruction}
A phase transition in the sense of
Chapter~\ref{ch:phase-transitions} generates a
non-trivial obstruction class
$[\omega] \in \check{H}^1(\mathcal{U}, \mathcal{F})$ that
cannot be trivialised by any local modification of the
cover $\mathcal{U}$.
\end{proposition}

\begin{proof}
By Theorem~\ref{thm:non-local-propagation}, the failure at
a phase transition is global: no sequence of local repairs
converges to a global section.  Equivalently, the
obstruction class is not in the image of any coboundary
$\delta^0$ with respect to any local modification of the
$0$-cochains.  Therefore $[\omega] \neq 0$ and cannot be
trivialised within the current framework.
\end{proof}

\begin{corollary}[Medium Shift as Cover Refinement]
\label{cor:medium-cover}
An admissible medium shift corresponds to passing to a
refined cover $\mathcal{U}'$ of a strictly larger space
$X'$ such that the induced obstruction class satisfies
$[\omega'] = 0$ in $\check{H}^1(\mathcal{U}',
\mathcal{F}')$.
\end{corollary}

\begin{proof}
By Definition~\ref{def:medium-shift}(iii), the obstruction
class vanishes in the lifted system.  This is precisely the
condition $[\omega'] = 0$, which by
Theorem~\ref{thm:obstruction-criterion} guarantees the
existence of a global section in the lifted framework.
\end{proof}

The sheaf-theoretic picture therefore unifies the framework
of Chapters~\ref{ch:yarncrawler}--\ref{ch:medium-shifts}
into a single cohomological narrative: knowledge is
reconstructibility, incompatibility is obstruction,
expansion is trivialisation, and medium shifts are the
passage to covers in which previously non-trivial classes
vanish.

%% ============================================================
%%  APPENDIX E — Implementation Sketches
%%  Purpose: ground the formalism without overengineering
%% ============================================================

\chapter{Implementation Sketches}
\label{app:implementation}

This appendix provides pseudocode and computational sketches
for the core algorithmic components of the monograph.  The
goal is not a production implementation but a precise
operational description: each sketch should be sufficient
to specify what the algorithm does, what it requires, and
what failure looks like.  All components are tied
explicitly to the formal structures of the main text.

\section{Spherepop Log Validation}
\label{sec:spherepop-impl}

The Spherepop event calculus of Chapter~\ref{ch:spherepop}
requires that every event in a log be admissible with
respect to the constraint relation $\mathcal{R} \subseteq
S \times S$.

\begin{algorithm}
\caption{Spherepop Log Validation}
\label{alg:spherepop}
\begin{algorithmic}
\Require Event log $\mathcal{L} = (e_1, \ldots, e_n)$
  where $e_k = (s_k \to s_{k+1})$; constraint relation
  $\mathcal{R}$; initial state $s_1$.
\Ensure \textsc{Valid} if $\mathcal{L}$ is admissible,
  \textsc{Invalid}$(k)$ at the first violation.
\State $s \gets s_1$
\For{$k = 1$ \textbf{to} $n$}
  \State $(s, s') \gets e_k$
  \If{$(s, s') \notin \mathcal{R}$}
    \State \textbf{return} \textsc{Invalid}$(k)$
      \Comment{Illegal transition at step $k$}
  \EndIf
  \State $s \gets s'$
\EndFor
\State \textbf{return} \textsc{Valid}
\end{algorithmic}
\end{algorithm}

\textbf{Failure mode.} An \textsc{Invalid}$(k)$ return
corresponds to the failure of coherence described in
Proposition~\ref{prop:failure-coherence}: the log contains
an illegal event that breaks admissibility and prevents
reliable reconstruction.

\textbf{Complexity.} $O(n \cdot |\mathcal{R}|)$ for a
naive membership check; $O(n \log |\mathcal{R}|)$ if
$\mathcal{R}$ is indexed.

\textbf{Relationship to CLIO.} When the validator returns
\textsc{Invalid}$(k)$, the CLIO operator (Section~E.2)
should be invoked to attempt repair of the constraint
relation, not to modify the log.  The log is the immutable
record; the constraint structure is what CLIO adjusts.

\section{CLIO Gradient Flow Discretisation}
\label{sec:clio-impl}

The CLIO operator of Chapter~\ref{ch:clio} is defined as a
gradient flow on the space of constraint structures:
$\dot{\mathcal{C}} = -\nabla_{\mathcal{C}}
\mathcal{E}(\mathcal{C})$.  In practice, the constraint
structure is represented by a finite set of parameters
$\theta \in \mathbb{R}^d$, and the coherence energy
$\mathcal{E}(\theta)$ is a differentiable function of these
parameters.

\begin{algorithm}
\caption{CLIO Gradient Descent}
\label{alg:clio}
\begin{algorithmic}
\Require Initial parameters $\theta_0$; coherence energy
  $\mathcal{E}(\theta)$ with gradient $\nabla_\theta
  \mathcal{E}$; step size schedule $\{\eta_t\}$;
  convergence tolerance $\varepsilon > 0$;
  maximum iterations $T_{\max}$.
\Ensure Fixed-point parameters $\theta^*$ or
  \textsc{NonConvergent} after $T_{\max}$ steps.
\State $\theta \gets \theta_0$, $t \gets 0$
\Repeat
  \State $g \gets \nabla_\theta \mathcal{E}(\theta)$
    \Comment{Compute coherence gradient}
  \State $\theta \gets \theta - \eta_t \cdot g$
    \Comment{Gradient descent step}
  \State $t \gets t + 1$
\Until{$\|g\| < \varepsilon$ \textbf{or} $t \ge T_{\max}$}
\If{$\|g\| < \varepsilon$}
  \State \textbf{return} $\theta^* \gets \theta$
    \Comment{Fixed point reached}
\Else
  \State \textbf{return} \textsc{NonConvergent}
    \Comment{Phase transition regime; see Section~E.4}
\EndIf
\end{algorithmic}
\end{algorithm}

\textbf{Coherence energy.} The energy $\mathcal{E}(\theta)$
should aggregate:
\begin{enumerate}[(i)]
  \item within-framework incoherence: violations of
    admissibility conditions internal to each framework;
  \item cross-framework incompatibility: states admissible
    in one framework but inadmissible in another, weighted
    by overlap.
\end{enumerate}
A simple instantiation: $\mathcal{E}(\theta) =
\sum_i w_i \cdot \mathrm{viol}_i(\theta)$ where
$\mathrm{viol}_i$ counts constraint violations in framework
$i$.

\textbf{Convergence criteria.}
Proposition~\ref{prop:clio-lyapunov} guarantees convergence
when $\mathcal{E}$ is proper and bounded below and the flow
is complete.  In practice, monitor:
\begin{enumerate}[(i)]
  \item gradient norm $\|g\|$ (approaches zero at
    fixed point);
  \item energy decrease per step (slowing signals
    approach to framework limit);
  \item step count to convergence (increasing cost is a
    diagnostic indicator per
    Proposition~\ref{prop:indicators}).
\end{enumerate}

\textbf{Non-convergence.} A \textsc{NonConvergent} return
after $T_{\max}$ steps indicates that CLIO cannot find a
fixed point within the current framework — the phase
transition regime of Chapter~\ref{ch:phase-transitions}.
The appropriate response is not to increase $T_{\max}$
but to trigger a medium shift (Section~E.4).

\section{TARTAN Recursive Tiling}
\label{sec:tartan-impl}

The TARTAN framework of Chapter~\ref{ch:tartan} decomposes
a domain into a hierarchy of tiles and enforces trajectory
consistency across scales.

\begin{algorithm}
\caption{TARTAN Tiling with Trajectory Consistency}
\label{alg:tartan}
\begin{algorithmic}
\Require Domain $\Omega$; depth $K$; local state function
  $s : \Omega \to Y$; trajectory data $\{x_t\}$;
  consistency tolerance $\varepsilon_{\mathrm{traj}}$.
\Ensure Recursive tiling $\{U_{i}^{(k)}\}$ with
  trajectory-consistency flags.
\State $\mathcal{T}^{(0)} \gets \{\Omega\}$
  \Comment{Level 0: full domain}
\For{$k = 1$ \textbf{to} $K$}
  \State $\mathcal{T}^{(k)} \gets \emptyset$
  \For{each $U \in \mathcal{T}^{(k-1)}$}
    \State $\{V_j\} \gets \mathrm{Subdivide}(U)$
      \Comment{Split into sub-tiles}
    \For{each pair $V_j, V_{j'}$ with $V_j \cap V_{j'} \neq \emptyset$}
      \State Check: $s(V_j)|_{V_j \cap V_{j'}} \approx
        s(V_{j'})|_{V_j \cap V_{j'}}$ within $\varepsilon_{\mathrm{traj}}$
      \If{inconsistent}
        \State Flag $(V_j, V_{j'})$ as
          \textsc{OverlapInconsistent}
      \EndIf
    \EndFor
    \State $\mathcal{T}^{(k)} \gets \mathcal{T}^{(k)} \cup \{V_j\}$
  \EndFor
  \State Check inter-scale consistency between
    $\mathcal{T}^{(k-1)}$ and $\mathcal{T}^{(k)}$
\EndFor
\State Compute Yarncrawler obstruction class from
  overlap inconsistencies (Section~E.4)
\State \textbf{return} $\{\mathcal{T}^{(k)}\}$ with
  consistency report
\end{algorithmic}
\end{algorithm}

\textbf{Local-global gap.} Overlap consistency within each
level does not guarantee a vanishing obstruction class
(Proposition~\ref{prop:tartan-obstruction-app}).  The
Yarncrawler check in the final step is essential for
detecting global incoherence that local checks miss.

\section{Diagnostics and Phase Transition Detection}
\label{sec:diagnostics}

The three diagnostic indicators of
Proposition~\ref{prop:indicators} can be monitored
concurrently with the CLIO run.

\begin{algorithm}
\caption{Phase Transition Diagnostics}
\label{alg:diagnostics}
\begin{algorithmic}
\Require CLIO run history $\{(\theta_t, g_t,
  \mathcal{E}(\theta_t))\}_{t=1}^T$;
  obstruction computation oracle.
\Ensure Diagnostic report: \textsc{Converging},
  \textsc{SlowingConvergence}, or
  \textsc{PhaseTransitionIndicated}.
\State Compute repair cost: $\bar{t} \gets$ mean steps
  to $\|g\| < \varepsilon$ over recent windows
\State Compute convergence rate: $\rho \gets
  \|g_{T}\| / \|g_{T/2}\|$
  \Comment{Ratio $> 1$ signals slowing}
\State Compute obstruction class $[\omega]$ from TARTAN
  output
\If{$\|g_T\| < \varepsilon$}
  \State \textbf{return} \textsc{Converging}
\ElsIf{$\rho > \rho_{\max}$ \textbf{or} $[\omega] \neq 0$}
  \State \textbf{return} \textsc{PhaseTransitionIndicated}
    \Comment{Trigger medium shift protocol}
\Else
  \State \textbf{return} \textsc{SlowingConvergence}
    \Comment{Monitor; increase logging frequency}
\EndIf
\end{algorithmic}
\end{algorithm}

\textbf{Degenerative drift detection.}  Separately, monitor
discriminative power $\mathrm{Disc}(\mathcal{C}_t)$ over
time.  A monotonically decreasing sequence with no
corresponding decrease in $\mathcal{E}$ signals degenerative
drift (Chapter~\ref{ch:degenerative-drift}) rather than
healthy convergence: the constraint structure is weakening,
not improving.

\section{Simulation Notes}
\label{sec:simulation-notes}

\textbf{RSVP numerical stability.}  The RSVP system
(Appendix~\ref{app:rsvp}) is stiff when diffusion
coefficients $D_\Phi, D_S$ are small relative to coupling
terms.  Use implicit time-stepping (Crank--Nicolson or
backward Euler) for the diffusion terms and explicit
treatment of the nonlinear reaction terms.  Monitor the
Courant--Friedrichs--Lewy condition for the advection term
in Eq.~\eqref{eq:v-evolution}.

\textbf{CLIO step size.}  Use adaptive step sizes: reduce
$\eta_t$ when $\mathcal{E}(\theta_t)$ increases (line
search) and increase when convergence is fast and stable.
A simple rule: $\eta_{t+1} = \eta_t \cdot 0.9$ after any
step that increases $\mathcal{E}$.

\textbf{Obstruction computation.}  For small covers, the
Čech obstruction class $[\omega]$ can be computed directly
by solving the linear system arising from the coboundary
conditions of Appendix~\ref{app:sheaves}.  For large
covers, approximate methods based on persistent homology
are available; see \citet{bredon1997sheaf} for the
algebraic foundation.

\textbf{Visualisation.}  For RSVP fields, plot $\Phi$,
$|\boldsymbol{v}|$, and $S$ as heatmaps over $\Omega$ at
successive time steps.  The constraint surface $\Sigma^*$
can be visualised as the level set $\{V = 0\}$ of the
Lyapunov function.  For TARTAN tilings, use a quad-tree
or oct-tree rendering with colour coding for consistency
flags.  For CLIO convergence, plot $\mathcal{E}(\theta_t)$
and $\|g_t\|$ against iteration count on a logarithmic
scale; linear decay of $\log\|g_t\|$ indicates geometric
convergence.



%% ============================================================
%%  APPENDIX F — Algebraic-Geometric Grounding:
%%               K-Trivial Fibrations and Boundedness
%%  Purpose: formal basis for Chapter 15 §§5–8
%% ============================================================

\chapter[Algebraic-Geometric Grounding]{Algebraic-Geometric Grounding: K-Trivial Fibrations and Boundedness}
\label{app:algebraic-geometry}

This appendix develops the algebraic-geometric foundations
underlying the fibration theory of
Chapter~\ref{ch:reconstruction}, \S\S\ref{sec:fibration-canonical}--\ref{sec:obstruction-decomposition}.
It states the main definitions and results from the
geometry of K-trivial varieties in a form that makes the
correspondence with the constraint-theoretic framework
explicit.  References are given to the primary literature
throughout.

\section{K-Trivial Varieties and Their Fibrations}
\label{sec:k-trivial}

\begin{definition}[K-Trivial Variety]
A smooth projective variety $\mathcal{X}$ over $\mathbb{C}$
is \emph{K-trivial} if its canonical bundle $K_{\mathcal{X}}$
is numerically trivial:
\[
  K_{\mathcal{X}} \equiv 0.
\]
The three principal classes are:
\begin{enumerate}[(i)]
  \item \emph{Calabi--Yau varieties}: $K_{\mathcal{X}}
    \sim_{\mathbb{Q}} 0$ with $h^{i,0}(\mathcal{X}) = 0$
    for $0 < i < \dim \mathcal{X}$;
  \item \emph{Abelian varieties}: complex tori $\mathcal{X}
    \cong \mathbb{C}^n / \Lambda$ with $K_{\mathcal{X}}$
    trivial;
  \item \emph{Irreducible holomorphic symplectic varieties}
    (hyperkähler): $\mathcal{X}$ admits a holomorphic
    symplectic form $\sigma \in H^0(\mathcal{X},
    \Omega^2_{\mathcal{X}})$ and $h^{2,0} = 1$.
\end{enumerate}
\end{definition}

K-trivial varieties are by themselves wildly unbounded in
general: there is no bound on their topology or geometry
that holds across all examples in each class.  The central
observation of \citet{engel2025} is that imposing a
fibration structure introduces hidden constraints that
collapse this freedom.

\begin{definition}[Admissible Fibration]
A \emph{fibration} on a K-trivial variety $\mathcal{X}$ is
a surjective morphism with connected fibers,
\[
  f : \mathcal{X} \to Y,
\]
where $Y$ is a normal projective variety with
$0 < \dim Y < \dim \mathcal{X}$, and the general fiber
$\mathcal{X}_y = f^{-1}(y)$ is a smooth K-trivial variety
of the same type (abelian or symplectic).
\end{definition}

\section{The Canonical Bundle Formula}
\label{sec:canonical-bundle-formula}

The fundamental tool for analysing fibrations of K-trivial
varieties is the canonical bundle formula of Kawamata,
Kodaira, and their successors.

\begin{theorem}[Canonical Bundle Formula
  {\citep{fujino2018, prokhorovshokurov2009}}]
\label{thm:canonical-bundle-formula}
Let $f : \mathcal{X} \to Y$ be an admissible fibration of
K-trivial varieties.  Then there exists an effective
$\mathbb{Q}$-divisor $B_Y$ and a nef $\mathbb{Q}$-divisor
$M_Y$ on $Y$ such that
\[
  K_{\mathcal{X}} \;\sim_{\mathbb{Q}}\;
    f^*\!\left(K_Y + B_Y + M_Y\right).
\]
\end{theorem}

The divisors $B_Y$ and $M_Y$ have precise geometric
meanings.

\begin{definition}[Boundary Divisor]
The \emph{boundary divisor} $B_Y = \sum_P b_P \cdot P$ is
an effective $\mathbb{Q}$-divisor on $Y$ supported on the
discriminant locus of $f$: the set of points $y \in Y$
over which the fiber $\mathcal{X}_y$ is singular or
reducible.  The coefficient $b_P$ encodes the severity of
degeneration over $P$.
\end{definition}

\begin{definition}[Moduli Divisor]
The \emph{moduli divisor} $M_Y$ is a nef
$\mathbb{Q}$-divisor on $Y$ encoding the variation of
fiber type across $Y$.  It vanishes identically if and
only if all smooth fibers of $f$ are isomorphic
(isotrivial fibration).  In general, $M_Y = \mathcal{P}^*
\lambda$ where $\mathcal{P} : Y \dashrightarrow \mathcal{M}$
is the period map into a moduli space of fiber types and
$\lambda$ is the Hodge bundle on $\mathcal{M}$.
\end{definition}

\begin{remark}
The correspondence with the canonical decomposition of
Definition~\ref{def:canonical-decomposition} is exact:
$K_Y \leftrightarrow \Phi_Y$ (base constraint field),
$B_Y \leftrightarrow \Phi_{\mathrm{sing}}$ (defect field),
$M_Y \leftrightarrow \Phi_{\mathrm{var}}$ (variation
memory field).
\end{remark}

\section{Variation of Hodge Structure and Period Maps}
\label{sec:vhs}

The moduli divisor $M_Y$ is controlled by the variation of
Hodge structure on the cohomology of the fibers.

\begin{definition}[Variation of Hodge Structure]
Let $f : \mathcal{X} \to Y$ be an admissible fibration.
The \emph{variation of Hodge structure} (VHS) associated
to $f$ is the local system $\mathbb{H} = R^k f_*
\mathbb{Z}$ on the smooth locus $Y^\circ$ of $Y$,
equipped with:
\begin{enumerate}[(i)]
  \item the Hodge filtration $F^\bullet \mathcal{H}$ on
    the corresponding flat bundle $\mathcal{H} =
    \mathbb{H} \otimes_{\mathbb{Z}} \mathcal{O}_{Y^\circ}$;
  \item the flat Gauss--Manin connection $\nabla :
    \mathcal{H} \to \mathcal{H} \otimes \Omega^1_{Y^\circ}$;
  \item Griffiths transversality: $\nabla(F^p\mathcal{H})
    \subseteq F^{p-1}\mathcal{H} \otimes \Omega^1$.
\end{enumerate}
\end{definition}

\begin{definition}[Period Map]
The \emph{period map} associated to the VHS is the
holomorphic map (defined on $Y^\circ$ and extending
meromorphically to $Y$):
\[
  \mathcal{P} : Y^\circ \to \Gamma \backslash D,
\]
where $D$ is the period domain classifying Hodge structures
of the given type, and $\Gamma$ is the monodromy group.
\end{definition}

The period map encodes all variation data: two fibers
$\mathcal{X}_{y_1}$ and $\mathcal{X}_{y_2}$ have
isomorphic Hodge structures if and only if $\mathcal{P}(y_1)
= \mathcal{P}(y_2)$.  The moduli divisor $M_Y$ is the
pullback under $\mathcal{P}$ of the natural ample bundle
on the Baily--Borel compactification $\overline{\Gamma
\backslash D}$ \citep{deligne1971, griffiths1968}.

\begin{proposition}[Boundedness of Period Map Image]
\label{prop:period-map-bounded}
If the base $Y$ lies in a bounded family, then the image
$\mathcal{P}(Y^\circ)$ in $\Gamma \backslash D$ is bounded.
Consequently, $M_Y$ lies in a bounded linear system.
\end{proposition}

\begin{proof}
The period map is a holomorphic map from a bounded family
of varieties into the arithmetic quotient $\Gamma \backslash
D$, which has a canonical compactification with controlled
geometry \citep{griffiths1968,schmid1973}.  Boundedness
of $Y$ implies the image is bounded.
\end{proof}

\section{The Albanese Reduction}
\label{sec:albanese}

For abelian fibrations, the Albanese variety provides a
canonical linearisation.

\begin{definition}[Albanese Fibration]
Let $f : \mathcal{X} \to Y$ be an admissible fibration
with abelian fibers.  The \emph{Albanese fibration}
\[
  f_{\mathrm{Alb}} : \mathcal{X}_{\mathrm{Alb}} \to Y
\]
is the universal abelian scheme over $Y$ equipped with a
map $\mathcal{X} \to \mathcal{X}_{\mathrm{Alb}}$ over $Y$.
It is the \emph{linearised} version of $f$: the minimal
transport structure consistent with the variation data.
\end{definition}

\begin{remark}
In the RSVP language: $f_{\mathrm{Alb}}$ is the linear
transport backbone $\boldsymbol{v}_{\mathrm{lin}}$, while
the original $f$ is the full nonlinear dynamics.  The
difference $\mathcal{X} / \mathcal{X}_{\mathrm{Alb}}$
measures higher-order structure not captured by linear
transport.
\end{remark}

\section{Tate--Shafarevich Groups}
\label{sec:ts-groups}

\begin{definition}[Tate--Shafarevich Group]
Let $f_{\mathrm{Alb}} : A \to Y$ be an abelian scheme
over $Y$.  The \emph{Tate--Shafarevich group} is
\[
  \Sha(A/Y) \;:=\; H^1(Y_{\mathrm{\acute{e}t}},\, A),
\]
the first étale cohomology group with coefficients in the
group scheme $A$.  It classifies $A$-torsors over $Y$
that are locally (in the étale topology) isomorphic to
$A$ but not globally trivial.
\end{definition}

\begin{proposition}[Finiteness of the Tate--Shafarevich Group]
\label{prop:sha-finite}
In the Calabi--Yau case (abelian fibration with trivial
canonical bundle), $\Sha(A/Y)$ is a finite abelian group.
In the symplectic case (Lagrangian fibration), $\Sha(A/Y)$
fits into an exact sequence
\[
  0 \to \mathcal{D}_Y \to \Sha(A/Y) \to \mathcal{T}_Y \to 0
\]
with $\mathcal{T}_Y$ finite and $\mathcal{D}_Y \cong
\mathbb{C}^n$ for some $n \ge 0$.
\end{proposition}

\begin{proof}
The Calabi--Yau finiteness follows from the Zarhin trick
\citep{zarhin1977} and finiteness results for abelian
varieties over function fields
\citep{faltings1983,dolgachevgross1994}.  The symplectic
decomposition follows from the structure of the
Tate--Shafarevich group for Lagrangian fibrations on
irreducible holomorphic symplectic manifolds, where the
continuous part $\mathcal{D}_Y$ arises from the
Beauville--Bogomolov form \citep{beauville1983,huybrechts2003}.
\end{proof}

The finiteness (or controlled decomposition) of $\Sha(A/Y)$
is the third and final boundedness input in the proof of
Theorem~\ref{thm:constraint-closure-boundedness}.

\section{The Main Boundedness Theorem}
\label{sec:main-boundedness}

We state the result of \citet{engel2025} in full
generality, then record the constraint-theoretic
interpretation.

\begin{theorem}[Boundedness of Fibered K-Trivial Varieties
  {\citep{engel2025}}]
\label{thm:engel-boundedness}
Let $\mathcal{S}$ be the set of all K-trivial varieties
$\mathcal{X}$ that admit a fibration
$f : \mathcal{X} \to Y$ with:
\begin{enumerate}[(i)]
  \item $Y$ a variety with klt singularities and
    $-(K_Y + B_Y + M_Y)$ ample;
  \item general fibers abelian varieties or irreducible
    holomorphic symplectic varieties of fixed Hodge
    numbers.
\end{enumerate}
Then $\mathcal{S}$ is a \emph{bounded family}: there exists
a scheme of finite type $\mathcal{T}$ and a morphism
$\mathcal{U} \to \mathcal{T}$ such that every element of
$\mathcal{S}$ appears as a fiber of $\mathcal{U} \to
\mathcal{T}$.
\end{theorem}

\begin{remark}[Constraint-Theoretic Reading]
Theorem~\ref{thm:engel-boundedness} is the precise
algebraic-geometric content of
Theorem~\ref{thm:constraint-closure-boundedness}.
The three conditions on $Y$, the fibers, and the moduli
data correspond to the three boundedness conditions
(i)--(iii) of Theorem~\ref{thm:constraint-closure-boundedness}.
The conclusion that $\mathcal{S}$ is a bounded family
is the algebraic-geometric formulation of the statement
that constraint closure implies bounded complexity.
\end{remark}

\section{Correspondence Table}
\label{sec:correspondence-table}

We summarise the correspondence between the
algebraic-geometric and constraint-theoretic languages.

\medskip
\begin{center}
\begin{tabular}{ll}
\toprule
\textbf{Algebraic geometry} & \textbf{Constraint theory} \\
\midrule
Total space $\mathcal{X}$ &
  Full field configuration \\
Base $Y$ &
  Base trajectory manifold \\
Fiber $\mathcal{X}_y$ &
  Local state space \\
Canonical bundle $K_{\mathcal{X}}$ &
  Global constraint density $\Phi_{\mathcal{X}}$ \\
Base term $K_Y$ &
  Base constraint field $\Phi_Y$ \\
Boundary divisor $B_Y$ &
  Defect field $\Phi_{\mathrm{sing}}$ \\
Moduli divisor $M_Y$ &
  Variation memory field $\Phi_{\mathrm{var}}$ \\
Period map $\mathcal{P}$ &
  Projection to latent variation space \\
Albanese fibration &
  Linear transport backbone \\
Tate--Shafarevich group $\Sha_Y$ &
  Gluing obstruction group $H^1(Y, \mathcal{P})$ \\
Torsion part of $\Sha_Y$ &
  Discrete obstruction $\mathcal{O}_{\mathrm{torsion}}$ \\
Continuous part of $\Sha_Y$ &
  Continuous obstruction $\mathcal{O}_{\mathrm{continuous}}$ \\
Bounded family &
  Constraint closure \\
\bottomrule
\end{tabular}
\end{center}
\medskip

The correspondence is not merely terminological: each
algebraic-geometric result translates to a constraint-theoretic
theorem, and the proofs in both languages follow the same
logical structure.  The canonical bundle formula is a
field decomposition law; the period map is a projection
operator; the Tate--Shafarevich exact sequence is a
constraint pipeline; and the main boundedness theorem is
constraint closure.


\backmatter

\bibliographystyle{plainnat}
\begin{thebibliography}{99}

\bibitem{thompson2023blindspot}
E. Thompson, A. Frank, and M. Gleiser,
\textit{The Blind Spot: Why Science Cannot Ignore Human Experience},
MIT Press, 2023.

\bibitem{whitehead1925science}
A. N. Whitehead,
\textit{Science and the Modern World},
Macmillan, 1925.

\bibitem{whitehead1929process}
A. N. Whitehead,
\textit{Process and Reality},
Macmillan, 1929.

\bibitem{shannon1948}
C. E. Shannon,
``A Mathematical Theory of Communication,''
\textit{Bell System Technical Journal}, 27:379--423, 623--656, 1948.

\bibitem{prigogine1984order}
I. Prigogine and I. Stengers,
\textit{Order Out of Chaos},
Bantam Books, 1984.

\bibitem{varela1991embodied}
F. J. Varela, E. Thompson, and E. Rosch,
\textit{The Embodied Mind},
MIT Press, 1991.

\bibitem{merleauponty1945phenomenology}
M. Merleau-Ponty,
\textit{Phenomenology of Perception},
Gallimard, 1945.

\bibitem{bergson1889time}
H. Bergson,
\textit{Time and Free Will},
1889.

\bibitem{maclane1998categories}
S. Mac Lane,
\textit{Categories for the Working Mathematician},
Springer, 1998.

\bibitem{bredon1997sheaf}
G. E. Bredon,
\textit{Sheaf Theory},
Springer, 1997.

\bibitem{evans2010partial}
L. C. Evans,
\textit{Partial Differential Equations},
AMS, 2010.

\bibitem{strogatz2015nonlinear}
S. H. Strogatz,
\textit{Nonlinear Dynamics and Chaos},
Westview Press, 2015.

\bibitem{cover2006elements}
T. M. Cover and J. A. Thomas,
\textit{Elements of Information Theory},
Wiley, 2006.

\bibitem{chalmers1996conscious}
D. J. Chalmers,
\textit{The Conscious Mind},
Oxford University Press, 1996.


\bibitem{engel2025}
P.~Engel, S.~Filipazzi, F.~Greer, M.~Mauri, and R.~Svaldi,
``Boundedness of some fibered K-trivial varieties,''
arXiv:2507.00973, 2025.

\bibitem{birkar2018}
C.~Birkar,
``Log Calabi--Yau fibrations,''
arXiv:1811.10709, 2018.

\bibitem{haconmckernanxu2014}
C.~Hacon, J.~McKernan, and C.~Xu,
``ACC for log canonical thresholds,''
\textit{Ann. of Math.} 180 (2014), 523--571.

\bibitem{fujino2018}
O.~Fujino,
\textit{Foundations of the Minimal Model Program},
MSJ Memoirs, 35, 2018.

\bibitem{prokhorovshokurov2009}
Y.~Prokhorov and V.~Shokurov,
``Towards the second main theorem on complements,''
\textit{J. Algebraic Geom.} 18 (2009), 151--199.

\bibitem{beauville1983}
A.~Beauville,
``Vari\'{e}t\'{e}s K\"{a}hl\'{e}riennes dont la premi\`{e}re classe de Chern est nulle,''
\textit{J. Differential Geom.} 18 (1983), 755--782.

\bibitem{huybrechts2003}
D.~Huybrechts,
``Compact hyperk\"{a}hler manifolds: basic results,''
\textit{Invent. Math.} 135 (1999), 63--113.

\bibitem{gross1994}
M.~Gross,
``A finiteness theorem for elliptic Calabi--Yau threefolds,''
\textit{Duke Math. J.} 74 (1994), 271--299.

\bibitem{dolgachevgross1994}
I.~Dolgachev and M.~Gross,
``Elliptic threefolds I: Ogg--Shafarevich theory,''
\textit{J. Algebraic Geom.} 3 (1994), 39--80.

\bibitem{faltings1983}
G.~Faltings,
``Endlichkeitss\"{a}tze f\"{u}r abelsche Variet\"{a}ten \"{u}ber Zahlk\"{o}rpern,''
\textit{Invent. Math.} 73 (1983), 349--366.

\bibitem{zarhin1977}
Y.~Zarhin,
``Endomorphisms of abelian varieties and points of finite order in characteristic~$p$,''
\textit{Math. USSR Izv.} 11 (1977), 31--45.

\bibitem{deligne1971}
P.~Deligne,
``Th\'{e}orie de Hodge II,''
\textit{Publ. Math. IH\'ES} 40 (1971), 5--57.

\bibitem{griffiths1968}
P.~Griffiths,
``Periods of integrals on algebraic manifolds I,''
\textit{Amer. J. Math.} 90 (1968), 568--626.

\bibitem{schmid1973}
W.~Schmid,
``Variation of Hodge structure: the singularities of the period mapping,''
\textit{Invent. Math.} 22 (1973), 211--319.

\bibitem{matsushita1999}
D.~Matsushita,
``On fibre space structures of a projective irreducible symplectic manifold,''
\textit{Topology} 38 (1999), 79--83.

\bibitem{campana2004}
F.~Campana,
``Orbifolds, special varieties and classification theory,''
\textit{Ann. Inst. Fourier} 54 (2004), 499--630.

\end{thebibliography}

\end{document}
